%% tutorial.tex -- a self-contained, step-by-step tutorial on the paper
%% "Leaving the Ring: Hash-Based Proofs of Shuffle and Decryption for BGV Mix-Nets" (main.tex).
%% Build:  latexmk -pdf tutorial.tex
%%
%% Every number in the examples was computed, not typed by hand:
%%   * the "mini" examples over Z_17 by a short Python script written for this tutorial;
%%   * the "toy run" (q = 7681, N = 4, two ballots) by ../toy_example.py, which runs the
%%     reference implementation (fp_core.py, mix_proof.py) unchanged over a small field.
%%     Its full output is ../docs/TOY_EXAMPLE.md.
%%   * the code parameters of the section "The commitment" were recomputed and agree with Table 3 of main.tex.
%%   * the quadratic-test example of the section "A simpler instantiation" by ../docs/tutorial_scripts/ligero_mini.py.

\documentclass[11pt,a4paper]{article}

\usepackage[margin=2.4cm]{geometry}
\usepackage{amsmath,amssymb,amsthm}
\usepackage{mathtools}
\usepackage{booktabs}
\usepackage{enumitem}
\usepackage{xcolor}
\usepackage[most]{tcolorbox}
\usepackage{array}
\usepackage{longtable}
\usepackage[hidelinks]{hyperref}

\setlength{\parskip}{0.45em}
\setlength{\parindent}{0pt}

%% ---------------------------------------------------------------- notation (same as main.tex)
\newcommand{\Z}{\mathbb{Z}}
\newcommand{\Zq}{\mathbb{Z}_q}
\newcommand{\Rq}{R_q}
\newcommand{\F}{\mathbb{F}}
\newcommand{\cube}[1]{\{0,1\}^{#1}}
\newcommand{\eq}{\mathsf{eq}}
\newcommand{\Enc}{\mathsf{Enc}}
\newcommand{\EncZ}{\mathsf{Enc}_0}
\newcommand{\Dec}{\mathsf{Dec}}
\newcommand{\coeff}{\mathsf{coeff}}
\newcommand{\Mpk}{M_{\mathsf{pk}}}
\newcommand{\Rmix}{\mathcal{R}_{\mathsf{mix}}}
\newcommand{\Rdec}{\mathcal{R}_{\mathsf{dec}}}
\newcommand{\RS}{\mathsf{RS}}
\newcommand{\ip}[2]{\langle #1, #2\rangle}
\newcommand{\mle}[1]{\widetilde{#1}}
\newcommand{\msk}[1]{#1^{\star}}
\newcommand{\epsprox}{\varepsilon_{\mathsf{prox}}}
\newcommand{\epsrbr}{\varepsilon_{\mathsf{rbr}}}
\newcommand{\nq}{n_{\mathsf{q}}}
\newcommand{\uh}{\hat{u}}
\newcommand{\hh}{\hat{h}}
\newcommand{\zh}{\hat{z}}
\newcommand{\ch}{\hat{c}}
\newcommand{\zpub}{z^{\mathsf{pub}}}
\newcommand{\id}{\mathrm{id}}

%% ---------------------------------------------------------------- boxes
\tcbset{enhanced,breakable,boxrule=0.5pt,arc=2pt,left=6pt,right=6pt,top=4pt,bottom=4pt,
        fonttitle=\bfseries\small}
\newtcolorbox{example}[1][]{colback=blue!3,colframe=blue!45!black,title={Example. #1}}
\newtcolorbox{toyrun}[1][]{colback=green!3,colframe=green!35!black,title={Toy run. #1}}
\newtcolorbox{why}[1][]{colback=orange!4,colframe=orange!60!black,title={Why? #1}}
\newtcolorbox{caution}[1][]{colback=red!3,colframe=red!50!black,title={Careful. #1}}
\newtcolorbox{recap}[1][]{colback=gray!5,colframe=gray!60!black,title={Recap. #1}}

\theoremstyle{definition}
\newtheorem{definition}{Definition}[section]
\theoremstyle{plain}
\newtheorem{lemma}[definition]{Lemma}
\newtheorem{theorem}[definition]{Theorem}
\newtheorem{proposition}[definition]{Proposition}
\newtheorem{fact}[definition]{Fact}

\title{A Complete Tutorial on\\[0.3em]
\emph{Leaving the Ring: Hash-Based Proofs of Shuffle and Decryption for BGV Mix-Nets}\\[0.4em]
\large Every definition, every step, every number}
\author{Companion to \texttt{paper/main.tex}}
\date{November 2026}

\begin{document}
\maketitle

\begin{abstract}
This document explains the paper in \texttt{main.tex} from first principles. It assumes only
school algebra: what an integer, a polynomial and a sum are. Everything else --- modular
arithmetic, rings and fields, lattice (BGV) encryption, mix-nets, zero-knowledge proofs,
multilinear polynomials, the sumcheck protocol, Reed--Solomon codes, Merkle trees, Fiat--Shamir
--- is defined before it is used, and illustrated with a small example whose numbers can be
checked by hand. The protocol of the paper is then traced message by message on a real run of
the reference implementation with two ballots. The security theorems are restated, every term of
every bound is explained, and each proof is walked through step by step.
\end{abstract}

\tableofcontents
\newpage

%% =====================================================================================
\section{How to read this tutorial}
%% =====================================================================================

\paragraph{What the paper does, in one paragraph.}
In an electronic election with a \emph{mix-net}, encrypted ballots are passed through several
\emph{mix servers}. Each server shuffles the list and changes the look of every ciphertext
(``re-randomises'' it) so that nobody can link an output ballot to an input ballot, and hence to a
voter. At the end, \emph{decryption servers} jointly decrypt the shuffled list. Every server must
\emph{prove} that it did its job honestly, without revealing its secret (its permutation, or its key
share). The paper gives such proofs for ciphertexts of the lattice-based encryption scheme BGV, as
used by the mix-net of Aranha, Baum, Gj{\o}steen and Silde (CCS 2023). The earlier proof of that
mix-net was shown to be broken by Bootle, Lyubashevsky and Merino-Gallardo (PKC 2025). The new
proofs never do algebra in the ring where the attack lives; they move everything into a
\emph{field} by a random \emph{projection}, and then use hash-based proof techniques.

\paragraph{Two running examples.}
\begin{itemize}[leftmargin=1.3em]
\item The \textbf{mini examples} (blue boxes) use the prime $q=17$. All numbers are small enough to
  check with pen and paper. They illustrate one concept at a time.
\item The \textbf{toy run} (green boxes) is a real execution of the paper's shuffle proof, produced
  by the project's reference implementation (\texttt{toy\_example.py}) with $q=7681$, ring degree
  $N=4$ and two ballots. Its complete trace is in \texttt{docs/TOY\_EXAMPLE.md}. These numbers are
  harder to check by hand, but every one of them was computed by the code that the paper describes,
  and the verifier accepted the transcript.
\end{itemize}
Neither example has any security; real parameters are $q\approx 2^{78}$, $N=4096$ and thousands
of ballots.

\paragraph{Map of the document.}
Part~I (Sections~\ref{sec:modular}--\ref{sec:rs}) builds the background. Part~II
(Sections~\ref{sec:attack}--\ref{sec:limits}) follows the paper section by section. The boxes
labelled ``Why?'' answer the question a careful reader should ask at that point. Section~\ref{sec:simple} describes a simpler variant of the protocol that replaces the
sumchecks by Ligero's constraint tests, as a teaching aid. Part~III
(Sections~\ref{sec:formal-intro}--\ref{sec:fm-index}) explains the Rocq formalisation file by file,
from the bottom up, and how the files combine into the protocol theorems. Section~\ref{sec:glossary}
is a glossary of all symbols.

%% =====================================================================================
\part{Background}
%% =====================================================================================

%% =====================================================================================
\section{Arithmetic modulo a prime}
\label{sec:modular}
%% =====================================================================================

\paragraph{Integers modulo $q$.}
Fix a positive integer $q$. Two integers are \emph{congruent modulo $q$}, written $a\equiv b
\pmod q$, if $q$ divides $a-b$. Every integer is congruent to exactly one of $0,1,\dots,q-1$, its
\emph{residue}. The set $\Zq=\{0,1,\dots,q-1\}$ with addition and multiplication followed by
taking the residue is \emph{arithmetic modulo $q$}.

\begin{example}[$q=17$]
$9+12=21\equiv 4$, \quad $5\cdot 7=35\equiv 1$, \quad $3-8=-5\equiv 12 \pmod{17}$.
\end{example}

\paragraph{Centred representatives.} It is often convenient to represent residues by integers in
$(-q/2,q/2]$ instead of $[0,q)$. Modulo $17$, the residue $16$ is the integer $-1$ and $15$ is
$-2$. This matters for ``small'' values: the residue $16$ is \emph{small} because it stands for $-1$.
In the paper, a polynomial is \emph{ternary} if every coefficient is $-1$, $0$ or $1$, i.e.\ is one
of the residues $q-1,0,1$.

\paragraph{Inverses.} An element $a$ has an \emph{inverse} $a^{-1}$ if $a\cdot a^{-1}\equiv 1$.
When $q$ is prime, every non-zero $a$ has an inverse, and it can be computed as $a^{q-2}\bmod q$
(Fermat's little theorem says $a^{q-1}\equiv 1$).

\begin{example}[Inverses modulo 17]
\begin{center}\small
\begin{tabular}{c|cccccccccccccccc}
$a$      & 1 & 2 & 3 & 4 & 5 & 6 & 7 & 8 & 9 & 10 & 11 & 12 & 13 & 14 & 15 & 16\\\hline
$a^{-1}$ & 1 & 9 & 6 & 13 & 7 & 3 & 5 & 15 & 2 & 12 & 14 & 10 & 4 & 11 & 8 & 16
\end{tabular}
\end{center}
Check: $3\cdot 6=18\equiv 1$, and $2\cdot 9=18\equiv 1$.
\end{example}

\paragraph{Fields.} A \emph{field} is a set with addition, subtraction, multiplication and division
by non-zero elements that obey the usual rules (associativity, commutativity, distributivity). For a
prime $q$, $\Zq$ is a field. The crucial property of a field for this paper:

\begin{fact}[No zero divisors]\label{fact:nozd}
In a field, $x\cdot y=0$ implies $x=0$ or $y=0$. Consequently a non-zero polynomial of degree $d$
in one variable has at most $d$ roots.
\end{fact}

The second sentence follows from the first: if $f(a)=0$ then $f(X)=(X-a)g(X)$, and a root of $f$
other than $a$ must be a root of $g$, which has degree $d-1$. The whole paper is about the fact
that the ring used by BGV is \emph{not} a field and violates Fact~\ref{fact:nozd}.

\paragraph{Squares and non-squares.} An element $a\neq 0$ of $\Zq$ is a \emph{square}
(quadratic residue) if $a=x^2$ for some $x$. For odd prime $q$, exactly half of the non-zero
elements are squares. Modulo $17$, the squares are $\{1,2,4,8,9,13,15,16\}$, so $3$ is a
\emph{non-square}. This is needed in Section~\ref{sec:field}.

%% =====================================================================================
\section{The polynomial ring $R_q$ and why it is dangerous}
\label{sec:ring}
%% =====================================================================================

\subsection{Definition}

Let $N$ be a power of two. $\Rq=\Zq[X]/(X^N+1)$ is the set of polynomials of degree less than
$N$ with coefficients in $\Zq$,
\[
  x = x_0 + x_1X + \dots + x_{N-1}X^{N-1},
\]
added coefficient by coefficient, and multiplied as polynomials followed by the rule
$X^N=-1$ (``reduction modulo $X^N+1$''). The \emph{coefficient vector} of $x$ is
$\coeff(x)=(x_0,\dots,x_{N-1})\in\Zq^N$. For a tuple of ring elements, $\coeff$ concatenates the
vectors: $\coeff(u,v)\in\Zq^{2N}$.

\begin{example}[{Multiplication in $\Z_{17}[X]/(X^4+1)$}]
$(1+X)\cdot X^3 = X^3 + X^4 = X^3 - 1$, i.e.\ coefficient vector $(16,0,0,1)$.
The power $X^4$ ``wraps around'' to the constant position \emph{with a minus sign}. This is called
\emph{negacyclic} multiplication.
\end{example}

\subsection{Multiplication by a fixed element is a matrix}
\label{sec:negacyclic}

Fix $a\in\Rq$. The map $r\mapsto a\cdot r$ is \emph{linear} on coefficient vectors: $a(r+r')=ar+ar'$
and $a(\lambda r)=\lambda(ar)$ for $\lambda\in\Zq$. So there is an $N\times N$ matrix $A$ over $\Zq$ with
$\coeff(ar)=A\,\coeff(r)$. Column $j$ of $A$ is $\coeff(a\cdot X^j)$, which is the coefficient vector
of $a$ shifted down by $j$ places, with the entries that wrap around negated. Such a matrix is a
\emph{negacyclic matrix}.

\begin{example}[The matrix of $a=3+X+2X^3$ modulo 17]
$\coeff(a)=(3,1,0,2)$. Then
\[
A=\begin{pmatrix}3&15&0&16\\1&3&15&0\\0&1&3&15\\2&0&1&3\end{pmatrix}
\qquad (15\equiv -2,\ 16\equiv -1).
\]
Column 1 is $\coeff(aX)=\coeff(3X+X^2+2X^4)=\coeff(-2+3X+X^2)=(15,3,1,0)$.
For $r=(1,0,16,1)$, i.e.\ $r=1-X^2+X^3$ (a ternary element),
$A\,r = (2,3,12,4)$, and multiplying the polynomials directly gives the same $a\cdot r=(2,3,12,4)$.
\end{example}

\begin{why}[This is the key fact of the whole paper]
Multiplying by a \emph{public} ring element is just a public linear map. Linear maps can be
checked in \emph{any} field that contains $\Zq$. The paper uses nothing else about $\Rq$
(equation~(1) of the paper).
\end{why}

\subsection{CRT slots and zero divisors}
\label{sec:slots}

When $q\equiv 1\pmod{2N}$, the polynomial $X^N+1$ has $N$ distinct roots $\zeta_1,\dots,\zeta_N$
in $\Zq$. The map
\[
  x \;\longmapsto\; \bigl(x(\zeta_1),\,x(\zeta_2),\,\dots,\,x(\zeta_N)\bigr)\in\Zq^N
\]
is then a bijection that turns ring multiplication into \emph{component-wise} multiplication
(Chinese Remainder Theorem, CRT). The components are called \emph{CRT slots}. So $\Rq$ ``is''
$N$ independent copies of $\Zq$ glued together.

\begin{example}[{Slots of $\Z_{17}[X]/(X^4+1)$}]
$17\equiv 1\pmod 8$, and the four roots of $X^4+1$ modulo 17 are $2,8,15,9$ (indeed
$2^4=16\equiv-1$). A ring element has four slots $\bigl(x(2),x(8),x(15),x(9)\bigr)$.
\end{example}

\paragraph{Zero divisors.} Take $x$ with slots $(1,0,0,0)$ and $y$ with slots $(0,1,1,1)$. Both are
non-zero, but their product has slots $(0,0,0,0)$, i.e.\ $xy=0$.

\begin{example}[Two non-zero elements with product zero]
Interpolating the slots gives $x=(13,15,16,8)$ and $y=(5,2,1,9)$ (coefficient vectors modulo 17).
Multiplying negacyclically gives $x\cdot y=(0,0,0,0)$. Fact~\ref{fact:nozd} fails in $\Rq$.
The toy run of the paper's code shows the same phenomenon modulo 7681:
$x=(5761,6242,2766,1617)$, $y=(5761,1617,4915,6242)$, $xy=0$.
\end{example}

\begin{caution}
A polynomial identity with coefficients in $\Rq$ can therefore hold at a random point although the
two sides are different polynomials: Fact~\ref{fact:nozd}, and the Schwartz--Zippel lemma of
Section~\ref{sec:sz}, do not apply. Section~\ref{sec:attack} shows how this breaks the earlier
shuffle proof.
\end{caution}

%% =====================================================================================
\section{BGV encryption, re-randomisation and mix-nets}
\label{sec:bgv}
%% =====================================================================================

\subsection{Key generation, encryption, decryption}

Parameters: $N$ (power of two), a prime $q\equiv 1\pmod{2N}$, a small plaintext modulus $p$
(the paper uses $p=2$, so messages are bit strings of length $N$).

\begin{itemize}[leftmargin=1.3em]
\item \textbf{Key generation.} Choose a uniformly random $a\in\Rq$ and ternary $s,e$. The secret key is
  $s$; the public key is $(a,b)$ with $b=as+pe$.
\item \textbf{Encryption} of $m\in\Rq$ with coefficients in $\{0,\dots,p-1\}$: choose
  \emph{randomness} $w=(r,e_1,e_2)$, three ternary ring elements, and output
  \[
    \Enc(m;w)=(u,v)=(ar+pe_1,\;\; br+pe_2+m).
  \]
\item \textbf{Decryption} of $(u,v)$: compute $v-su$, take each coefficient as a centred integer, and
  reduce it modulo $p$.
\end{itemize}

\paragraph{Why decryption works.}
\[
  v-su = (as+pe)r+pe_2+m-s(ar+pe_1)= m + p\,(er+e_2-se_1).
\]
The term $er+e_2-se_1$ is a polynomial with small coefficients (products of ternary polynomials
have coefficients at most $N$ in absolute value). If all coefficients of $m+p(\dots)$ are smaller
than $q/2$ in absolute value, then the centred representative is the true integer, and reducing it
modulo $p$ kills $p(\dots)$ and leaves $m$. So BGV needs \emph{small} noise: this is why the
re-randomisers must be proven \emph{short} (ternary).

\begin{toyrun}[Decrypting ballot 0 by hand]
$q=7681$, $N=4$, $p=2$, secret $s=(-1,0,-1,1)$. Ciphertext $c_0=(u,v)$ with
$u=(1361,1337,263,6470)$, $v=(5244,4868,7268,1234)$.
Then $su=(5246,4870,7268,1235)$ and
\[
 v-su=(7679,7679,0,7680) \;\xrightarrow{\text{centre}}\; (-2,-2,0,-1)\;\xrightarrow{\bmod 2}\;(0,0,0,1).
\]
So $c_0$ encrypts $m_0=(0,0,0,1)$. In the same way $c_1$ gives $(-5,1,3,3)\mapsto(1,1,1,1)$.
\end{toyrun}

\subsection{Re-randomisation}

Write $\EncZ(w)=\Enc(0;w)=(ar+pe_1,\;br+pe_2)$, a fresh encryption of zero. Adding it to a ciphertext
(component-wise) gives a new ciphertext of the \emph{same} message, because
$\Enc(m;w)+\EncZ(w')=\Enc(m;w+w')$. The new ciphertext looks unrelated to the old one (this is the
semantic security of BGV, which relies on a lattice assumption; the \emph{proofs} in this paper do
not need that assumption).

\begin{caution}
$w+w'$ is no longer ternary (coefficients up to $\pm 2$). Each mix server makes the noise grow.
The paper proves statements about one server; the noise budget across servers is listed as future
work (paper, Section~10).
\end{caution}

\subsection{A mix-net and what must be proven}

A mix server receives $\tau$ ciphertexts $c_0,\dots,c_{\tau-1}$. It chooses a secret permutation
$\sigma$ of $\{0,\dots,\tau-1\}$ and secret ternary $w_0,\dots,w_{\tau-1}$, and outputs
\[
  \ch_j = c_{\sigma(j)}+\EncZ(w_{\sigma(j)}),\qquad j=0,\dots,\tau-1 .
\]
Output position $j$ holds input $\sigma(j)$, freshly re-randomised.

\begin{toyrun}[The instance used throughout]
$\tau=2$, public key $a=(7070,6955,463,750)$, $b=(1802,1013,7079,7048)$.
Inputs: $c_0$ above (encrypts $(0,0,0,1)$) and $c_1$ (encrypts $(1,1,1,1)$).
Secret: $\sigma=(1,0)$ (output 0 is input 1, output 1 is input 0) and
\[
\begin{array}{llll}
w_0:& r=(0,0,1,1) & e_1=(0,-1,0,1) & e_2=(-1,0,0,0)\\
w_1:& r=(0,1,0,0) & e_1=(1,0,1,1)  & e_2=(1,1,0,1)
\end{array}
\]
Outputs: $\ch_0=c_1+\EncZ(w_1)$ with $u=(1074,1474,152,7661)$, $v=(5280,6040,6484,7306)$, which
decrypts to $(1,1,1,1)$; and $\ch_1=c_0+\EncZ(w_0)$, which decrypts to $(0,0,0,1)$.
\end{toyrun}

\paragraph{Distributed decryption.} The secret key is shared among decryption servers. Server
$j$ holds a share $s$ (we drop the index) with public $b=as+pe$, and for each ciphertext
$(u_i,v_i)$ publishes the \emph{partial decryption} $t_i=su_i+pE_i$. The extra noise $E_i$ is large
(``flooding noise'', up to about $2^{55}$ in the parameters the paper reads from CCS 2023) so that
$t_i$ reveals nothing about $s$. Summing the partial decryptions of all servers and subtracting
from $v_i$ gives $m_i$ plus a multiple of $p$ plus noise, as above.

\subsection{The two relations}
\label{sec:relations}

A \emph{relation} is a set of pairs (statement, witness). The statement is public; the witness is
the secret that makes it true. The paper proves knowledge of a witness for two relations.

\begin{definition}[Shuffle relation $\Rmix$]
Statement $x=(a,b,(c_i)_{i<\tau},(\ch_j)_{j<\tau})$. Witness: a permutation $\sigma$ and ternary
$w_0,\dots,w_{\tau-1}\in\Rq^3$ with $\ch_j=c_{\sigma(j)}+\EncZ(w_{\sigma(j)})$ for all $j$.
\end{definition}

\begin{definition}[Decryption relation $\Rdec$]
Statement: $a,b$ and $u_i,t_i\in\Rq$ for $i<\tau$. Witness: ternary $s,e$ and $E_i\in\Rq$ with integer
coefficients in $[-B/2,B/2)$ such that (K)~$b=as+pe$ and (D)~$t_i=su_i+pE_i$ for all $i$.
\end{definition}

\begin{why}[Why must $w$ be ternary, and $E$ bounded?]
Without the size condition the relations are \emph{trivially} satisfiable for wrong outputs. Over
$\Zq$ the system $\ch_j-c_{\sigma(j)}=\EncZ(w)$ has $2N$ equations and $3N$ unknowns, so for almost
any target there are (large) solutions $w$. A cheating server could replace a vote and ``explain'' the
change with a huge $w$. Shortness is what makes $\EncZ(w)$ an encryption of zero that still
decrypts correctly. The toy test ``re-randomiser with a coefficient 2'' (Section~\ref{sec:cheat})
shows this check at work. Similarly, an unbounded $E_i$ would let a decryption server publish any
$t_i$ it likes.
\end{why}

%% =====================================================================================
\section{The extension field $\F=\mathrm{GF}(q^2)$}
\label{sec:field}
%% =====================================================================================

\paragraph{Construction.} Let $\nu$ be a non-square modulo $q$. Then $t^2-\nu$ has no root in $\Zq$,
and
\[
  \F=\Zq[t]/(t^2-\nu)=\{\,x_0+x_1t : x_0,x_1\in\Zq\,\},\qquad t^2=\nu,
\]
is a field with $q^2$ elements. Arithmetic:
\begin{align*}
(x_0+x_1t)+(y_0+y_1t)&=(x_0+y_0)+(x_1+y_1)t,\\
(x_0+x_1t)(y_0+y_1t)&=(x_0y_0+\nu x_1y_1)+(x_0y_1+x_1y_0)t,\\
(x_0+x_1t)^{-1}&=\frac{x_0-x_1t}{x_0^2-\nu x_1^2}.
\end{align*}
The denominator $x_0^2-\nu x_1^2$ (the \emph{norm}) is non-zero for $(x_0,x_1)\neq(0,0)$ exactly because
$\nu$ is not a square. $\Zq$ sits inside $\F$ as the elements with $x_1=0$.

\begin{example}[{$\mathrm{GF}(17^2)=\Z_{17}[t]/(t^2-3)$}]
$3$ is a non-square mod 17. Let $X=2+5t$, $Y=4+t$. Then
$XY=(8+3\cdot 5)+(2+20)t=23+22t\equiv 6+5t$.
Norm of $X$: $4-3\cdot 25=-71\equiv 14$, $14^{-1}=11$, so
$X^{-1}=(2-5t)\cdot 11=22-55t\equiv 5+13t$. Check: $(2+5t)(5+13t)=(10+3\cdot 65)+(26+25)t=205+51t\equiv 1+0t$.
\end{example}

\paragraph{The paper's choice.} $q=2^{78}+100663297$ and $\nu=3$ (a non-square modulo this $q$,
which we re-checked), so $|\F|=q^2\approx 2^{156}$. The toy run uses $q=7681$ and $\nu=13$, so
$|\F|=58\,997\,761\approx 2^{26}$; its elements are printed as \texttt{a+bt}.

\begin{why}[Why $q^2$ and not just $q$?]
Every cheating probability in the paper is of the form ``(small number)$/|\F|$''. With $|\F|=q\approx
2^{78}$ these would be about $2^{-65}$, too weak for 128-bit security. With $q^2\approx 2^{156}$ they
are about $2^{-140}$. Committed data (the witness) stays in $\Zq$; only the verifier's random
challenges and derived values live in $\F$.
\end{why}

\begin{why}[Why is it fine to check $\Zq$-equations in $\F$?]
$\Zq$ is a subfield of $\F$: addition and multiplication of elements of $\Zq$ give the same result in
both. So a linear equation with coefficients in $\Zq$ holds in $\F$ if and only if it holds in $\Zq$.
\end{why}

%% =====================================================================================
\section{Polynomials over a field and the Schwartz--Zippel lemma}
\label{sec:sz}
%% =====================================================================================

A polynomial in several variables $x_1,\dots,x_k$ has a \emph{total degree} (largest sum of
exponents in a monomial) and a \emph{degree in each variable}.

\begin{lemma}[Schwartz--Zippel]
Let $f$ be a non-zero polynomial over a field $\F$ of total degree $d$. If $x\in\F^k$ is uniform, then
$\Pr[f(x)=0]\le d/|\F|$.
\end{lemma}

For one variable this is Fact~\ref{fact:nozd}: at most $d$ roots out of $|\F|$ points.

\begin{example}[Comparing two lists by one random evaluation]
To test whether polynomials $f,g$ are equal, pick a random $x$ and compare $f(x)$ and $g(x)$. If
$f\neq g$, then $f-g$ is non-zero, and the test passes wrongly with probability at most
$\deg(f-g)/|\F|$. With $\deg=5$ and $|\F|=17$ that is $5/17$ (useless); with $|\F|\approx 2^{156}$ it is
negligible. All the verifier's tests in the paper are of this kind.
\end{example}

The lemma \emph{fails} over $\Rq$: $f(x)=x\cdot y$ with the zero divisor $y$ of Section~\ref{sec:slots} is
a non-zero polynomial of degree 1 in $x$ that vanishes on many more than one point.

%% =====================================================================================
\section{Tables, multilinear extensions and masks}
\label{sec:mle}
%% =====================================================================================

\subsection{Tables as functions on the Boolean cube}

A \emph{table} $T$ of length $2^m$ is a list $T[0],\dots,T[2^m-1]$. Write an index in binary,
$\text{index}=b_0+2b_1+\dots+2^{m-1}b_{m-1}$ with bits $b_i\in\{0,1\}$; variable $0$ is the
\emph{least significant} bit. Then $T$ is a function on the \emph{Boolean cube} $\cube{m}$.

\subsection{The eq polynomial}

\[
  \eq(t,x)=\prod_{i=0}^{m-1}\bigl(t_ix_i+(1-t_i)(1-x_i)\bigr).
\]
For Boolean $t,x$ each factor is $1$ if $t_i=x_i$ and $0$ otherwise, so $\eq(t,x)=1$ if $t=x$ and $0$
otherwise. For non-Boolean $t$ it is a polynomial of degree 1 in each variable. Two identities are
used constantly:
\[
  \textstyle\sum_{x\in\cube{m}}\eq(t,x)=1\ \text{ for every } t, \qquad
  \text{(each factor sums to } t_i+(1-t_i)=1\text{)}.
\]

\subsection{Multilinear extension}

\[
  \mle{T}(x)=\sum_{b\in\cube{m}}T[b]\;\eq(b,x).
\]
$\mle T$ is the unique polynomial that has degree at most 1 in each variable (``multilinear'') and
agrees with $T$ on the cube. It can be evaluated anywhere in $\F^m$.

\begin{example}[A table of length 4 over $\Z_{17}$]
$T=(3,5,7,10)$, so $T[b_0+2b_1]$: $T(0,0)=3$, $T(1,0)=5$, $T(0,1)=7$, $T(1,1)=10$.
\[
 \mle T(x_0,x_1)=3(1-x_0)(1-x_1)+5x_0(1-x_1)+7(1-x_0)x_1+10x_0x_1 .
\]
At $(x_0,x_1)=(2,3)$ the four $\eq$ weights $\eq(b,(2,3))$ for $b=(0,0),(1,0),(0,1),(1,1)$ are
$(1-2)(1-3)=2$, $2\cdot(-2)=-4\equiv 13$, $(-1)\cdot 3\equiv 14$, $2\cdot 3=6$; they sum to
$35\equiv 1$, as they must. So $\mle T(2,3)=3\cdot 2+5\cdot 13+7\cdot 14+10\cdot 6=229\equiv 8$.
\end{example}

\begin{why}[Why random evaluations of extensions are useful]
If two tables differ, their extensions are different multilinear polynomials; by Schwartz--Zippel
they differ at a random point except with probability $m/|\F|$. And $\sum_x\eq(t,x)F(x)=\mle F(t)$:
a weighted sum with random weights $\eq(t,\cdot)$ is zero for random $t$ only if the table $F$ is
zero everywhere (except with probability $m/|\F|$). The paper uses this to turn ``$F(x)=0$ for all
$x$'' into ``one sum is zero'' (sumchecks A and C, and the decryption claims).

Mini example: the table $F=(0,1,0,0)$ is not zero, and $\sum_b\eq(b,(2,3))F(b)=13\neq 0$.
\end{why}

\subsection{Splitting variables: low and high bits}

The witness tables of the paper store coefficient $y$ of ballot $i$ at index $y+Ni$. The low
$\log N$ bits are the coefficient index $y$, the high $\ell=\log\tau$ bits are the ballot index $i$.
Fixing the high variables at a point $r\in\F^\ell$ gives a table in $y$:
\[
  \mle W(y,r)=\sum_{i}\eq(i,r)\,W[y+Ni] .
\]
For $\tau=2$ this is just $(1-r)W[y]+rW[y+N]$: the two ballots ``folded'' at $r$.

\subsection{The vanishing polynomial and masked extensions}
\label{sec:masked}

$Z(x)=\prod_i x_i(1-x_i)$ is zero on every point of the cube (one factor vanishes) and non-zero at a
point $r$ unless some $r_i\in\{0,1\}$. For a random \emph{mask coefficient} $a_T\in\F$, the
\emph{masked extension} is
\[
  \msk T(x)=\mle T(x)+a_T\,Z(x).
\]
\begin{itemize}[leftmargin=1.3em]
\item On the cube $\msk T=T$, so anything proven about the table on the cube is unaffected.
\item At a point $r$ with $Z(r)\neq 0$, the value $\msk T(r)=\mle T(r)+a_TZ(r)$ is uniformly random,
  because $a_T$ is uniform and $a\mapsto \mle T(r)+aZ(r)$ is a bijection of $\F$. It reveals nothing
  about $T$.
\item The price: $\msk T$ has degree $2$ in each variable instead of $1$.
\end{itemize}

\begin{example}[Masking modulo 17]
$Z(2,3)=2\cdot(-1)\cdot 3\cdot(-2)=12$. With $a_T=4$: $\msk T(2,3)=8+4\cdot 12=56\equiv 5$, while on the
cube $\msk T(1,0)=5=T(1,0)$ is unchanged. Each of the 17 choices of $a_T$ gives a different value at
$(2,3)$; if $a_T$ is uniform, so is the value.
\end{example}

\begin{why}[Why the paper needs masks at all]
At the end of a sumcheck the prover must reveal evaluations like $\mle\sigma(r)$ of tables built from
its secret. For two ballots $\mle\sigma(r)=(1-r)\sigma_0+r\sigma_1$, from which $\sigma$ can be solved.
The masked value is uniform and reveals nothing (Section~\ref{sec:zk}).
\end{why}

%% =====================================================================================
\section{The sumcheck protocol}
\label{sec:sumcheck}
%% =====================================================================================

\subsection{The protocol}

Let $g$ be a polynomial in $m$ variables that the verifier can evaluate at \emph{one} point
(possibly with help, see below). The prover claims
\[
  \sum_{x\in\cube{m}}g(x)=S .
\]
Computing the sum directly costs $2^m$ evaluations. The sumcheck protocol of Lund, Fortnow, Karloff
and Nisan reduces it to one evaluation.

\begin{enumerate}[leftmargin=1.5em]
\item Round $j=0,1,\dots,m-1$. The \emph{running claim} is $S_j$ (initially $S_0=S$). The prover sends
  the one-variable polynomial
  \[
    p_j(X)=\sum_{x_{j+1},\dots,x_{m-1}\in\{0,1\}} g(r_0,\dots,r_{j-1},X,x_{j+1},\dots,x_{m-1}).
  \]
  The verifier checks $p_j(0)+p_j(1)=S_j$, draws a random $r_j\in\F$, and sets $S_{j+1}=p_j(r_j)$.
\item At the end the verifier checks $S_m=g(r_0,\dots,r_{m-1})$ with one evaluation.
\end{enumerate}

\paragraph{How a polynomial is sent.} If $p_j$ has degree $d_j$, it is determined by $d_j+1$
values. The paper's prover sends $p_j(0),p_j(2),p_j(3),\dots,p_j(d_j)$ and \emph{not} $p_j(1)$: the
verifier computes $p_j(1)=S_j-p_j(0)$, so the check $p_j(0)+p_j(1)=S_j$ holds by construction, and
evaluates $p_j(r_j)$ by Lagrange interpolation through the $d_j+1$ values at $0,1,\dots,d_j$:
\[
  p(r)=\sum_{i=0}^{d}p(i)\prod_{k\neq i}\frac{r-k}{i-k}.
\]

\begin{example}[A complete sumcheck over $\Z_{17}$]
Take $T=(3,5,7,10)$ from Section~\ref{sec:mle} and $g(x)=\mle T(x)^2$, so the claim is the sum of the
squares of the table: $9+25+49+100=183\equiv 13$. So $S=13$, $m=2$, degree $2$ in each variable.

\emph{Round 0} (variable $x_0$). $\mle T(X,0)=3+2X$ and $\mle T(X,1)=7+3X$, so
$p_0(X)=(3+2X)^2+(7+3X)^2=58+54X+13X^2\equiv 7+3X+13X^2$.
Prover sends $p_0(0)=7$ and $p_0(2)=14$. Verifier: $p_0(1)=13-7=6$ (indeed $7+3+13=23\equiv 6$).
Verifier draws $r_0=5$; new claim $S_1=p_0(5)=7+15+325=347\equiv 7$
(Lagrange through $(0,7),(1,6),(2,14)$ gives the same $7$).

\emph{Round 1} (variable $x_1$). $\mle T(5,Y)=13(1-Y)+22Y\equiv 13+9Y$, so
$p_1(Y)=(13+9Y)^2\equiv 16+13Y+13Y^2$. Prover sends $p_1(0)=16$, $p_1(2)=9$. Verifier:
$p_1(1)=7-16=-9\equiv 8$ (indeed $16+13+13=42\equiv 8$). Verifier draws $r_1=11$; $S_2=p_1(11)\equiv 15$.

\emph{Final check.} $\mle T(5,11)\equiv 10$ and $g(5,11)=100\equiv 15=S_2$. Accept.
\end{example}

\subsection{Why a false claim is caught}

Suppose the claim is false, $S\neq\sum g$. In round 0 the true polynomial $p_0^{\mathrm{true}}$ has
$p_0^{\mathrm{true}}(0)+p_0^{\mathrm{true}}(1)=\sum g\neq S$, so the polynomial the prover sends (which
satisfies the check with $S$) is a \emph{different} polynomial. Two different polynomials of degree
$d_0$ agree in at most $d_0$ points, so with probability at least $1-d_0/|\F|$ the new claim
$S_1=p_0(r_0)$ is again false (it is not the true sum of the remaining cube). By induction, the final
claim is false and the last check fails, except with probability at most
$\sum_jd_j/|\F|$.

\begin{example}[Cheating by one]
In the mini example, claim $S=14$ instead of $13$. If the prover keeps $p_0(0)=7$, the verifier
computes $p_0(1)=14-7=7$ instead of $6$: the prover is now committed to a different quadratic, whose
value at the random $r_0$ differs from the true $7$ except for at most 2 of the 17 values of $r_0$.
The toy run shows the same at full scale (Section~\ref{sec:toy-cheat}).
\end{example}

\paragraph{What the verifier still has to trust.} The final check needs $g(r)$. In the paper, $g$
is built from tables the prover \emph{committed} to; the prover \emph{sends} the needed evaluations
(e.g.\ $\msk u(r)$) and these become \emph{claims} that are settled later by the commitment opening
(Section~\ref{sec:opening}). This is the pattern of the whole proof: sumchecks turn claims about
entire tables into claims about single evaluations, and the opening checks those evaluations.

\subsection{Zero-knowledge sumcheck by masking}
\label{sec:zksumcheck}

The round polynomials $p_j$ depend on the secret tables. To hide them, the prover commits to a
random \emph{mask polynomial}
\[
  g(x)=a_0+\sum_{j=0}^{m-1}\sum_{t=1}^{d_j}c_{j,t}\,x_j^{\,t}
\]
(a constant plus one univariate polynomial per variable, with the same degrees $d_j$ as the sumcheck),
and also sends its sum $G=\sum_{x\in\cube m}g(x)$. The verifier draws $\rho\neq 0$ and the sumcheck is
run on $f+\rho g$ with claim $S+\rho G$. The round polynomials of $\rho g$ are uniformly random
(their non-constant coefficients are $2^{m-j-1}\rho c_{j,t}$ in round $j$), so they hide those of $f$.
At the end the verifier needs $g(r)$, which becomes one more linear claim on the committed
coefficients.

\begin{example}[The sum of a mask]
$m=2$, degree 2: $g=a_0+c_{0,1}x_0+c_{0,2}x_0^2+c_{1,1}x_1+c_{1,2}x_1^2$. Summing over the 4 points:
each $x_j^t$ is $1$ on half of them, so $G=4a_0+2(c_{0,1}+c_{0,2})+2(c_{1,1}+c_{1,2})$. With
$(a_0,c_{0,1},c_{0,2},c_{1,1},c_{1,2})=(6,2,9,1,3)$: $G=24+22+8=54\equiv 3$ modulo 17, and summing
the 4 values directly gives the same.
\end{example}

%% =====================================================================================
\section{Proofs: soundness, knowledge, zero knowledge, oracles, Fiat--Shamir}
\label{sec:proofs}
%% =====================================================================================

\paragraph{Interactive proof.} A prover $P$ and a verifier $V$ exchange messages about a public
statement $x$; at the end $V$ accepts or rejects. \emph{Public-coin}: every verifier message is fresh
uniform randomness (a \emph{challenge}).
\begin{itemize}[leftmargin=1.3em]
\item \textbf{Completeness:} an honest prover with a valid witness is always accepted.
\item \textbf{Soundness error} $\varepsilon$: if $x$ is false, every prover is accepted with probability
  at most $\varepsilon$.
\item \textbf{Knowledge soundness:} stronger. If a prover is accepted with probability above
  $\varepsilon$, a witness can be computed from it by an \emph{extractor}. In the paper the extractor is
  \emph{straight-line}: it just reads the committed rows and decodes them; no rewinding.
\item \textbf{Honest-verifier zero knowledge (HVZK):} there is a \emph{simulator} that, given only the
  statement (and the verifier's challenges), outputs transcripts with the same distribution as real
  ones. Then a transcript teaches nothing beyond ``the statement is true''. \emph{Statistical} ZK:
  the distributions are equal or negligibly close, without computational assumptions.
\end{itemize}

\paragraph{Interactive oracle proof (IOP).} Some prover messages are long. In an IOP they are
\emph{oracles}: the verifier does not read them but may query a few positions. In the paper the
oracles are the committed matrices; the verifier queries $\nq$ columns.

\paragraph{From IOP to a real proof (BCS transformation).} Ben-Sasson, Chiesa and Spooner (BCS)
replace each oracle by the root of a \emph{Merkle tree} (Section~\ref{sec:merkle}), each query answer
by the entry plus a Merkle path, and each verifier challenge by a hash of everything sent so far
(the \emph{Fiat--Shamir} transformation). The result is \emph{non-interactive}: a single string that
anyone can verify.

\paragraph{Random-oracle model (ROM).} The analysis models the hash function as a truly random
function that all parties can query. The quantum ROM (QROM) allows quantum queries.

\paragraph{Round-by-round soundness.} Fiat--Shamir lets a cheating prover \emph{retry}: it can hash
many variants of a message until it gets a lucky challenge. For a protocol with many rounds this is
dangerous unless each \emph{single} round has small error. Round-by-round (RBR) soundness, due to
Canetti et al.\ and used by Holmgren, captures exactly this: there is a notion of ``doomed'' partial
transcripts such that a false statement starts doomed, a complete doomed transcript is rejected, and
any \emph{one} challenge rescues a doomed transcript with probability at most $\epsrbr$. Then $t$
hash queries give error about $t\cdot\epsrbr$ (classical) and $O(t^2\epsrbr)$ (quantum, Chiesa,
Manohar and Spooner). Section~\ref{sec:fs} applies this.

%% =====================================================================================
\section{Hash functions and Merkle trees}
\label{sec:merkle}
%% =====================================================================================

SHA-256 maps any byte string to 32 bytes. It is used as a \emph{commitment}: hard to find two inputs
with the same output (binding), and, if the input contains a secret random \emph{salt}, the output
reveals nothing about the rest (hiding).

A \emph{Merkle tree} commits to a list of $n$ items with one 32-byte \emph{root}. Leaves are hashes of
the items; each inner node is the hash of its two children; the root is the top. To \emph{open} item
$j$, the prover reveals the item and the $\log_2 n$ sibling hashes on the path to the root (the
\emph{Merkle path}); the verifier recomputes the root.

\begin{example}[Four columns]
Columns $\mathrm{col}_0,\dots,\mathrm{col}_3$, salts $s_0,\dots,s_3$ (16 bytes each).
Leaf $\ell_j=\mathrm{SHA256}(s_j\,\|\,\mathrm{col}_j)$; $n_{01}=\mathrm{SHA256}(\ell_0\|\ell_1)$,
$n_{23}=\mathrm{SHA256}(\ell_2\|\ell_3)$, root $=\mathrm{SHA256}(n_{01}\|n_{23})$.
To open column 2: send $\mathrm{col}_2$, $s_2$, and the path $(\ell_3,\,n_{01})$.
The verifier computes $\ell_2$, then $n_{23}$, then the root, and compares.
(With the byte strings used in our script, the root begins \texttt{85d1e9de}.)
\end{example}

\paragraph{Binding (paper, Lemma 4.1).} In the ROM, a root fixes at most one value per leaf except
with probability $t^2/2^{256}$ after $t$ hash queries (finding two different inputs with the same
hash).

\paragraph{Salts.} Each leaf has a fresh 16-byte salt. Unopened leaves are then uniformly random
strings to anyone who does not know the salt, which the zero-knowledge simulator needs
(paper, Theorem 5.3: the ``hashes of unopened leaves'' are the only values that differ, with
probability $t/2^{128}$).

%% =====================================================================================
\section{Reed--Solomon codes}
\label{sec:rs}
%% =====================================================================================

\subsection{Encoding}

Let $\omega\in\Zq$ have order $n$ ($\omega^n=1$, smaller powers $\neq 1$; it exists when $n$ divides
$q-1$). A message $(x_0,\dots,x_{k-1})$ with $k\le n$ is read as the polynomial
$P(X)=\sum_i x_iX^i$ and encoded as the \emph{codeword}
\[
  \RS(x)=\bigl(P(1),P(\omega),P(\omega^2),\dots,P(\omega^{n-1})\bigr)\in\Zq^n .
\]
Position $j$ of the codeword is called \emph{column} $j$ (because many encoded rows will be stacked
into a matrix). The \emph{rate} is $k/n$.

\begin{example}[$q=17$, $n=8$, $\omega=2$, $k=3$]
Powers of $2$: $(1,2,4,8,16,15,13,9)$. A row with $C=2$ data symbols $(4,9)$ and $1$ pad $\pi=6$ is the
polynomial $P=4+9X+6X^2$, and
\[
  \RS(4,9,6)=(2,\,12,\,0,\,1,\,1,\,10,\,13,\,10).
\]
Check column 1: $P(2)=4+18+24=46\equiv 12$.
\end{example}

\subsection{Distance}

The \emph{Hamming distance} between two words is the number of positions where they differ.
Two different polynomials of degree $<k$ agree in at most $k-1$ points (Fact~\ref{fact:nozd} applied
to their difference), so two different codewords differ in at least $d=n-k+1$ columns: the code has
\emph{minimum distance} $d$.

\begin{example}
Changing only the pad from 6 to 7 gives $\RS(4,9,7)=(3,16,16,14,2,14,12,6)$, which agrees with
$\RS(4,9,6)$ in $0$ of the $8$ columns; the bound guarantees at most $k-1=2$.
\end{example}

\paragraph{Linearity.} $\RS(\lambda x+\mu y)=\lambda\RS(x)+\mu\RS(y)$, because evaluation is linear in
the coefficients. Example: $\RS(1,2,3)=(6,0,6,5,2,9,7,7)$ and
$3\RS(4,9,6)+5\RS(1,2,3)=(2,2,13,11,13,7,6,14)=\RS\bigl(3\cdot(4,9,6)+5\cdot(1,2,3)\bigr)$.
This is what lets the verifier check \emph{combinations} of committed rows by encoding a single vector.

\subsection{Close and far words; unique decoding}

A word $y\in\F^n$ is \emph{within distance $e$} of the code (``close'') if some codeword differs from
it in at most $e$ columns. If $2e<d$, that codeword is unique (two codewords within $e$ of $y$ would
be within $2e<d$ of each other): this is \emph{unique decoding}. The paper uses
\[
  e=\lfloor (d-1)/\Delta\rfloor,\qquad \Delta\in\{3,2\}\ \text{(the ``proximity regime'')},
\]
so $2e<d$ in both regimes.

\subsection{Pads make a few columns look random}
\label{sec:hiding}

In the paper a row holds $C$ data symbols followed by $\nq$ uniform pads, $k=C+\nq$. Column $j$ is
\[
  \sum_{i<C}x_i\omega^{ji}\;+\;\omega^{jC}\sum_{t<\nq}\pi_t\,\omega^{jt}.
\]
For any set $J$ of at most $\nq$ columns, the pad part restricted to $J$ is the invertible diagonal
matrix $\mathrm{diag}(\omega^{jC})_{j\in J}$ times a Vandermonde matrix $(\omega^{jt})_{j\in J,t<|J|}$
(padded with more columns if $|J|<\nq$), which has full rank because the nodes $\omega^j$ are distinct.
So uniform pads make the opened entries uniform and independent of the data
(paper, Lemma 4.4).

\begin{example}[One pad hides one column]
Column 1 of the row $(4,9,\pi)$ is $4+9\cdot 2+\pi\cdot 4=22+4\pi$. As $\pi$ runs over $0,\dots,16$,
this takes the values $5,9,13,0,4,8,12,16,3,7,11,15,2,6,10,14,1$: every residue exactly once.
\end{example}

\begin{why}[Why the number of pads equals the number of queries $\nq$]
The verifier opens $\nq$ columns. Each opened column reveals one linear combination of the row's
symbols. With $\nq$ pads, $\nq$ such combinations are still uniformly random. One more opened column
could leak data.
\end{why}

%% =====================================================================================
\part{The paper, step by step}
%% =====================================================================================

%% =====================================================================================
\section{The problem: why the earlier proof broke (paper \S1 and \S3)}
\label{sec:attack}
%% =====================================================================================

\subsection{Proving a shuffle by comparing products}

The classical idea (Neff, 2001) for showing that list $(\hat x_j)$ is a reordering of list $(x_i)$:
for a random $\alpha$, check
\[
  \prod_{i}(x_i-\alpha)\;=\;\prod_{j}(\hat x_j-\alpha).
\]
Both sides are polynomials in $\alpha$ whose roots are the list entries. Over a field, if the lists
are different multisets, the polynomials are different, have degree $\tau$, and agree at a random
$\alpha$ with probability at most $\tau/|\F|$ (Fact~\ref{fact:nozd}).

\subsection{Why it fails in $\Rq$: the slot attack}

The CCS 2023 shuffle argument tested a polynomial identity of this type with coefficients in $\Rq$.
By the CRT (Section~\ref{sec:slots}), multiplication in $\Rq$ is slot by slot, so the identity holds in
$\Rq$ if it holds \emph{in every slot separately}. In every slot the test only sees the
\emph{multiset of slot values}. A cheater can therefore exchange the values of a single slot between
two list entries: every slot keeps the same multiset, so the products agree for \emph{every} $\alpha$,
but the two entries have become different ring elements.

\begin{example}[Slot swap modulo 17]
In $\Z_{17}[X]/(X^4+1)$ take $f=(1,2,3,4)$ and $g=(5,0,7,1)$. Their slots (values at $2,8,15,9$) are
\[
  f\mapsto(15,13,11,16),\qquad g\mapsto(7,13,8,9).
\]
Exchange the first slot: $f'\mapsto(7,13,11,16)$, $g'\mapsto(15,13,8,9)$. Interpolating back,
$f'=(16,1,11,8)$ and $g'=(7,1,16,14)$. Then $\{f,g\}\neq\{f',g'\}$, yet for $\alpha=(2,0,1,5)$ (any
$\alpha$ works):
\[
  (f-\alpha)(g-\alpha)=(10,3,13,13)=(f'-\alpha)(g'-\alpha).
\]
Even the sums agree: $f+g=f'+g'=(6,2,10,5)$.
\end{example}

For ciphertexts this means: exchange one slot between two output ciphertexts. Both decrypt to
\emph{changed} votes, and the ring-level identity still holds.

\begin{toyrun}[The slot attack against the toy instance]
Outputs 0 and 1 exchange one of the four CRT slots. The attacked outputs decrypt to $(1,1,1,0)$ and
$(0,0,0,0)$ instead of $(1,1,1,1)$ and $(0,0,0,1)$. A product test computed in $\Rq$ accepted in
\textbf{20 of 20} trials. The paper's verifier rejects (Section~\ref{sec:toy-cheat}).
\end{toyrun}

\begin{why}[The fix in one sentence]
Never test an identity in $\Rq$. Turn every ciphertext into \emph{one field element} using a random
linear form that mixes all coefficients (and therefore all slots), and test identities in the field.
\end{why}

%% =====================================================================================
\section{Step 1 of the idea: the relation is linear (paper \S2.1)}
\label{sec:linear}
%% =====================================================================================

Write $\coeff(w)=(\coeff(r),\coeff(e_1),\coeff(e_2))\in\Zq^{3N}$. Since $a,b$ and $p$ are public,
\[
  \coeff(\EncZ(w))=\begin{pmatrix}\coeff(ar+pe_1)\\\coeff(br+pe_2)\end{pmatrix}
  =\underbrace{\begin{pmatrix}A & pI & 0\\ B & 0 & pI\end{pmatrix}}_{\Mpk\ \in\ \Zq^{2N\times 3N}}
  \begin{pmatrix}\coeff(r)\\\coeff(e_1)\\\coeff(e_2)\end{pmatrix},
\]
where $A$ and $B$ are the negacyclic matrices of $a$ and $b$ (Section~\ref{sec:negacyclic}) and $I$ is
the $N\times N$ identity. This is equation~(1) of the paper. Hence
\[
  \ch_j=c_{\sigma(j)}+\EncZ(w_{\sigma(j)})
  \iff
  \coeff(\ch_j)=\coeff(c_{\sigma(j)})+\Mpk\coeff(w_{\sigma(j)}),
\]
a system of linear equations over $\Zq$ in the unknown entries of the $w_i$, plus two
\emph{non-linear side conditions}: $\sigma$ is a permutation, and all entries of $w$ are in
$\{-1,0,1\}$. The proof handles the three parts separately:
\begin{center}
\begin{tabular}{@{}lll@{}}
\toprule
part & technique & section\\\midrule
linear equations & random projection + sumcheck B & \ref{sec:projection}, \ref{sec:protocol}\\
permutation & tagged logarithmic derivative + sumcheck A & \ref{sec:logder}, \ref{sec:protocol}\\
shortness & $w^3-w=0$ + sumcheck C & \ref{sec:protocol}\\
\bottomrule
\end{tabular}
\end{center}

\begin{toyrun}[The matrix of $a$]
For $a=(7070,6955,463,750)$ the negacyclic matrix is
\[
A=\begin{pmatrix}7070&6931&7218&726\\6955&7070&6931&7218\\463&6955&7070&6931\\750&463&6955&7070\end{pmatrix}
\]
(e.g.\ $6931=7681-750$: the wrapped entry is negated). For $r_0=(0,0,1,1)$: $A\,r_0=(263,6468,6320,6344)$,
equal to $a\cdot r_0$ computed in $\Rq$.
\end{toyrun}

%% =====================================================================================
\section{Step 2 of the idea: random projection (paper \S3, ``Projection'')}
\label{sec:projection}
%% =====================================================================================

\subsection{Definition}

The verifier draws $\rho=(\rho_u,\rho_v)\in\F^{2N}$ uniformly ($N$ field elements for the $u$-part of a
ciphertext and $N$ for the $v$-part). The \emph{projection} of a ciphertext $c$ is
\[
  \ip{\rho}{\coeff(c)}=\sum_{y<N}\rho_u[y]\,u_y+\sum_{y<N}\rho_v[y]\,v_y\ \in\F .
\]
The projection of the \emph{re-randomised} input $i$ is
\[
  z_i=\ip{\rho}{\coeff(c_i)+\Mpk\coeff(w_i)}
     =\underbrace{\ip{\rho}{\coeff(c_i)}}_{\zpub_i\ \text{(public)}}
     +\ip{\underbrace{\rho\Mpk}_{\rho'\in\F^{3N}}}{\coeff(w_i)} ,
\]
and the projection of output $j$ is $\zh_j=\ip{\rho}{\coeff(\ch_j)}$ (public).

\paragraph{Computing $\rho'$.} $\rho'=\rho\Mpk$ is a row vector times a matrix. With the block form
of $\Mpk$:
\[
  \rho'_1=A^{\top}\rho_u+B^{\top}\rho_v,\qquad \rho'_2=p\,\rho_u,\qquad \rho'_3=p\,\rho_v ,
\]
because $\ip{\rho_u}{A\,\coeff(r)}=\ip{A^{\top}\rho_u}{\coeff(r)}$ (the \emph{adjoint}). The verifier
computes this itself: two adjoint multiplications (cheap, like a polynomial multiplication) plus
$4\tau N$ products ``$\Zq$-element times $\F$-element'' for the $2\tau$ projections. This is linear in
the size of the statement (paper, ``The verifier's linear algebra'').

\begin{toyrun}[Projection]
$\rho_u=(975{+}2617t,\ 4116{+}4192t,\ 5301{+}7045t,\ 840{+}7242t)$,
$\rho_v=(1829{+}7330t,\ 4923{+}5091t,\ 4559{+}3445t,\ 6417{+}4690t)$.

Public projections: $\zpub=(6489{+}7426t,\ 2997{+}3299t)$, $\zh=(5055{+}1164t,\ 833{+}6567t)$.

$\rho'_r=(1313{+}7398t,\dots)$, $\rho'_{e_1}=(1950{+}5234t,\ 551{+}703t,\dots)$,
$\rho'_{e_2}=(3658{+}6979t,\dots)$. Check $\rho'_{e_1}=2\rho_u$: $2\cdot 975=1950$, $2\cdot 2617=5234$. \checkmark

Only the prover can compute $z_i=\zpub_i+\ip{\rho'}{\coeff(w_i)}$:
$z=(833{+}6567t,\ 5055{+}1164t)$. Indeed $z_0=\zh_1$ and $z_1=\zh_0$, matching $\sigma=(1,0)$.
\end{toyrun}

\subsection{Why the projection is faithful}

\begin{fact}
If $x\neq x'$ in $\Zq^{2N}$, then $\Pr_\rho[\ip{\rho}{x}=\ip{\rho}{x'}]=1/|\F|$.
\end{fact}
\begin{proof}
$x-x'$ has a non-zero entry, say at position $k$. For any fixed values of the other entries of $\rho$,
exactly one value of $\rho_k$ makes $\ip{\rho}{x-x'}=0$.
\end{proof}

So with overwhelming probability, ``$\zh_j=z_{\sigma(j)}$ for all $j$'' holds exactly when the shuffle
relation holds (for the linear part). The shuffle of $\tau$ ciphertexts, each $2N$ numbers, has become a
statement about two lists of $\tau$ field elements.

\begin{example}[The projection sees the slot swap]
With $\rho=(3,7,1,10)$ over $\Z_{17}$ (projecting just one polynomial, for illustration):
$\ip{\rho}{f}=9$, $\ip{\rho}{g}=15$, but $\ip{\rho}{f'}=10$, $\ip{\rho}{g'}=14$. The multisets $\{9,15\}$
and $\{10,14\}$ differ, so the field test catches the swap that the ring test missed. The projection
looks at \emph{coefficients}, and a slot swap changes coefficients.
\end{example}

%% =====================================================================================
\section{Step 3 of the idea: the tagged logarithmic derivative (paper \S3, \S5.2)}
\label{sec:logder}
%% =====================================================================================

\subsection{From products to sums of inverses}

The product test would need the prover to reveal the common value of the two products. That value
depends on the secret permutation (an earlier version of the design did this and, with two ballots,
revealed $\sigma$; see the paper, \S3). The \emph{logarithmic derivative} (Hab\"ock) replaces the
product by a sum of inverses. Taking the derivative of $\log\prod_i(X-x_i)$ gives
$\sum_i 1/(X-x_i)$, and
\[
  \sum_i\frac{1}{x_i-\alpha}=\sum_j\frac{1}{\hat x_j-\alpha}
\]
for random $\alpha$ is again a multiset test.

\subsection{Tags}

Each value is paired with a \emph{tag}: input $i$ with its index $i$, output $j$ with the committed
value $\sigma_j$ (the input it claims to come from). For random $\alpha,\beta$:
\[
  h_i=z_i+\beta\, i-\alpha,\qquad \hh_j=\zh_j+\beta\,\sigma_j-\alpha,\qquad
  \text{check}\ \ \sum_i\frac1{h_i}=\sum_j\frac1{\hh_j}.
\]

\begin{lemma}[Logarithmic derivative with projection; paper Lemma 5.1]\label{lem:logder}
Let $S,T$ be multisets of $\tau<q$ elements of $\Zq^{2N}\times\Zq$. If $S\neq T$, then
\[
  \sum_{(x,y)\in S}\frac{1}{\ip{\rho}{x}+\beta y-\alpha}=\sum_{(x,y)\in T}\frac{1}{\ip{\rho}{x}+\beta y-\alpha},
\]
with all denominators non-zero, holds with probability at most $(2\tau-1)/|\F|$ over uniform
$\rho\in\F^{2N}$, $\alpha,\beta\in\F$.
\end{lemma}

Here an element $(x,y)$ is (coefficient vector of a ciphertext, tag). On the left $x=\coeff(c_i)+\Mpk\coeff(w_i)$,
$y=i$; on the right $x=\coeff(\ch_j)$, $y=\sigma_j$.

\paragraph{Proof, step by step.}
\begin{enumerate}[leftmargin=1.5em]
\item Treat $\rho,\beta,\alpha$ as formal variables $R=(R_1,\dots,R_{2N}),B,A$. Each distinct pair
  $(x,y)$ gives a \emph{linear form} $L_{x,y}=\ip{R}{x}+By-A$.
\item Distinct pairs give distinct forms (the coefficients of the form \emph{are} $x$ and $y$). All forms
  have coefficient $-1$ at $A$, so no form is a constant multiple of another; distinct linear forms of
  this kind are irreducible and pairwise coprime.
\item Let $\mu_{x,y}$ = (multiplicity in $S$) $-$ (multiplicity in $T$). Since $S\neq T$, some
  $\mu_{x^*,y^*}\neq 0$. As $|\mu|\le\tau<q$, it is non-zero also as an element of $\F$ (this is
  where $\tau<q$ is needed).
\item The difference of the two sums is $\sum_{(x,y)}\mu_{x,y}/L_{x,y}$ over the $D\le 2\tau$ distinct
  pairs. Multiply by $\prod L_{x,y}$ to clear denominators: the result
  $P=\sum_{(x,y)}\mu_{x,y}\prod_{(x',y')\neq(x,y)}L_{x',y'}$ is a polynomial of degree $\le D-1\le 2\tau-1$.
\item $P\neq 0$: modulo $L_{x^*,y^*}$, every term except the one for $(x^*,y^*)$ contains the factor
  $L_{x^*,y^*}$ and vanishes, leaving $\mu_{x^*,y^*}\prod_{\text{others}}L$, which is non-zero modulo
  $L_{x^*,y^*}$ by coprimality.
\item If the sums are equal with non-zero denominators, then $P(\rho,\beta,\alpha)=0$. By Schwartz--Zippel
  this happens with probability $\le(2\tau-1)/|\F|$. \qed
\end{enumerate}

\begin{why}[What the lemma gives]
If the check passes, the multiset $\{(\coeff(c_i)+\Mpk\coeff(w_i),\,i)\}$ equals
$\{(\coeff(\ch_j),\,\sigma_j)\}$. The left tags are $0,1,\dots,\tau-1$, each exactly once (and distinct in
$\Zq$ because $\tau<q$). So the right tags $\sigma_j$ are also $0,\dots,\tau-1$ each exactly once:
$\sigma$ is a permutation. And the pair with tag $\sigma_j$ on the left is input $\sigma_j$, so
$\coeff(\ch_j)=\coeff(c_{\sigma(j)})+\Mpk\coeff(w_{\sigma(j)})$, i.e.\ the linear part of the shuffle relation.
\end{why}

\begin{example}[Log-derivative modulo 17]
$z=(5,11)$, $\zh=(11,5)$, $\sigma=(1,0)$, $\alpha=3$, $\beta=7$.
$h=(5+0-3,\ 11+7-3)=(2,15)$, $\hh=(11+7-3,\ 5+0-3)=(15,2)$.
Inverses: $u=(2^{-1},15^{-1})=(9,8)$, $\uh=(8,9)$. Sums: $9+8=17\equiv 0$ and $8+9\equiv 0$. Equal.

\emph{A changed output.} If $\zh=(5,12)$ (output 1 altered), with the same tags $\hh=(9,9)$ and
$\sum 1/\hh=2\cdot 2=4\neq 0$. Rejected.

\emph{A non-permutation.} Suppose two inputs have the same projection, $z=(5,5)$, and
$\zh=(5,5)$. With $\sigma=(1,0)$ both sums are $11$. With the non-permutation $\sigma=(0,0)$, the
right side becomes $1/2+1/2=9+9=18\equiv 1\neq 11$: the tags catch it.
\end{example}

\begin{toyrun}[Tagged values and inverses]
$\alpha=4487{+}6900t$, $\beta=5984{+}6361t$.
\[
\begin{array}{ll}
h=(4027{+}7348t,\ 6552{+}625t), & \hh=(6552{+}625t,\ 4027{+}7348t)\\
u=1/h=(6227{+}4673t,\ 128{+}7228t), & \uh=1/\hh=(128{+}7228t,\ 6227{+}4673t)\\
\textstyle\sum u=6355{+}4220t, & \textstyle\sum\uh=6355{+}4220t .
\end{array}
\]
$\hh$ holds the same two values as $h$, in the order given by $\sigma$.
\end{toyrun}

\begin{caution}[A remark of this tutorial, not of the paper]
For this particular relation the untagged check $\{\coeff(c_i)+\Mpk\coeff(w_i)\}=\{\coeff(\ch_j)\}$ would
already imply that \emph{some} permutation matches inputs to outputs. The tags make the committed table
$\sigma$ itself that permutation, so the extractor reads $\sigma$ off the commitment directly, and they
follow the tagged design the authors set out to port from Bootle et al. The paper does not discuss
an untagged variant; we mention it only so that the role of the tags is not overstated.
\end{caution}

\subsection{What the prover must show}

The verifier cannot compute $h_i$ (it depends on the secret $w_i$) and cannot see $\sigma$. So the prover
commits the inverse tables $u=1/h$ and $\uh=1/\hh$ and proves three facts about tables:
\begin{equation}\label{eq:threefacts}
  u_i\cdot h_i=1\ \ \forall i,\qquad \uh_j\cdot\hh_j=1\ \ \forall j,\qquad \sum_iu_i=\sum_j\uh_j .
\end{equation}
The first two say the committed tables really are the inverses (and that no $h_i,\hh_j$ is zero);
the third is the log-derivative identity. Plus, separately, that every entry of $W$ is ternary.

%% =====================================================================================
\section{Overview of the shuffle proof and of its cost (paper \S3)}
\label{sec:overviewproof}
%% =====================================================================================

\begin{enumerate}[leftmargin=1.5em]
\item The prover commits the witness: $W=(W_r,W_{e_1},W_{e_2})$ (three tables of length $\tau N$)
  and $\sigma$.
\item The verifier sends $\rho,\alpha,\beta$; the prover commits $u,\uh$.
\item \textbf{Sumcheck A} proves the three facts \eqref{eq:threefacts}, ending with claimed values of
  $\msk u,\msk\uh,\msk\sigma,\msk h$ at a random point $r$.
\item $h$ is \emph{never committed}: it is a linear image of $W$. \textbf{Sumcheck B} reduces the
  claimed $\msk h(r)$ to one linear claim about $W$.
\item \textbf{Sumcheck C} proves $w^3-w=0$ for all $3\tau N$ entries of $W$. Over $\Zq$ with $q>3$,
  $w^3-w=(w-1)w(w+1)=0$ iff $w\in\{-1,0,1\}$ (no zero divisors in a field!). For $w=2$: $8-2=6\neq 0$.
\item What remains: nine linear claims about committed rows, settled by \textbf{one batched opening}.
\end{enumerate}

\paragraph{Why the size grows like a square root.} The witness has $3\tau N$ symbols. Split into rows of
$C$ symbols, it is a matrix with $3\tau N/C$ rows. The opening sends a constant number of vectors of
length about $C$ (row combinations) plus $\nq$ columns, each with one entry per row, i.e.\ $\nq\cdot
3\tau N/C$ entries. Total $O(C+\nq\tau N/C)$, minimised when $C\approx\sqrt{\nq\tau N}$, giving
$O(\sqrt{\nq\tau N})$. Per ballot this is $O(\sqrt{\nq N/\tau})$: \emph{the proof per ballot shrinks as
the number of ballots grows}.

%% =====================================================================================
\section{The commitment (paper \S4.1)}
\label{sec:commit}
%% =====================================================================================

\subsection{Row layout and code length}

A row holds $C$ data symbols in $\Zq$ followed by $\nq$ uniform pads, so $k=C+\nq$. It is encoded with
the Reed--Solomon code of length $n$ on the $n$-th roots of unity of $\Zq$ (for the paper's $q$,
$q-1$ is divisible by $2^{25}$, so such roots exist for every power of two $n$ up to $2^{24}$ used by the
code). The length is $n=C/\varrho$ for a nominal rate $\varrho\in\{1/4,1/2\}$, and is doubled if the pads
make $k/n$ larger than $2\varrho$.

\subsection{Proximity parameter and number of queries}

With $d=n-k+1$, $\Delta\in\{3,2\}$ and $\lambda=128$:
\[
  e=\lfloor (d-1)/\Delta\rfloor,\qquad \nq=\min\{\,t:\ (1-e/n)^t\le 2^{-\lambda}\,\}.
\]
Note the circularity: $\nq$ depends on $e$, which depends on $k=C+\nq$. It is solved by iteration
(start from $\nq=0$, recompute until it stabilises).

\begin{example}[Recomputing the paper's Table 3, big matrix, 4096 ballots, $\Delta=3$]
$\tau N=2^{12}\cdot 2^{12}=2^{24}$, $C=2^{16}=65536$, rate $1/4$, so $n=2^{18}=262144$.
After iteration $\nq=309$, so $k=65845$, $d=n-k+1=196300$, $e=\lfloor 196299/3\rfloor=65433$,
$e/n\approx 0.2496$, and $(1-0.2496)^{309}\approx 2^{-128.01}\le 2^{-128}$ while $308$ queries would not
suffice. The number of rows is $3\cdot 2^{24}/2^{16}=768$ data rows plus $6$ mask rows (three
$\F$-rows, each stored as two $\Zq$-rows) $=774$. All match the paper's table.

For the small matrix ($C=1024$, $n=4096$): $\nq=356$, $k=1380$, $d=2717$, $e=905$.
For $\Delta=2$ at 4096 ballots ($C=32768$, $n=2^{17}$): $e=49057$, $e/n\approx 0.374$, $\nq=190$, $k=32958$.
\end{example}

\begin{why}[Why $\Delta=2$ needs fewer queries]
A larger $e$ (a larger fraction of columns where a cheater must be wrong) means each query catches
a cheater with higher probability, so fewer queries are needed for $2^{-128}$. $\Delta=2$ allows
$e\approx d/2$, but its proximity lemma needs a deeper theorem (Section~\ref{sec:opening-lemmas}).
\end{why}

\subsection{Merkle commitment over columns}

A batch of rows forms a matrix (rows = encoded rows, $n$ columns). Leaf $j$ is SHA-256 of a fresh 16-byte
salt followed by column $j$ (one entry per row, 10 bytes per $\Zq$ entry). The commitment is the Merkle
root. Committing \emph{columns} (not rows) is what allows the verifier to check a combination of rows
at a column by opening that single column.

\subsection{What is committed in the shuffle proof}

\begin{center}
\begin{tabular}{@{}lll@{}}
\toprule
matrix & content & when\\\midrule
big ($C_1$) & $W_r,W_{e_1},W_{e_2}$ (witness), 3 mask rows over $\F$ & round 1\\
small, batch 1 ($C_2$) & $\sigma$, 3 mask rows & round 1\\
small, batch 2 & $u,\uh$, and one row with the \emph{mask table} & after $\rho,\alpha,\beta$\\
\bottomrule
\end{tabular}
\end{center}
The mask table holds $a_u,a_{\uh},a_\sigma,a_h,a_{w_1},a_{w_2},a_{w_3}$ (Section~\ref{sec:masked}) and
the coefficients of the three sumcheck masks $g_A,g_B,g_C$ (Section~\ref{sec:zksumcheck}).

\begin{toyrun}[Round 1]
Tables (coefficient $y$ of ballot $i$ at index $y+4i$):
$W_r=(0,0,1,1,\,0,1,0,0)$, $W_{e_1}=(0,-1,0,1,\,1,0,1,1)$, $W_{e_2}=(-1,0,0,0,\,1,1,0,1)$, $\sigma=(1,0)$.
Big matrix: $C_1=8$ data symbols, $\nq=3$ pads, $n=32$; small matrix: $C_2=128$, $\nq=3$, $n=512$.
(Toy security $q$-bits $=1$, hence only 3 queries.) Each table fits in one row. The prover sends two
32-byte roots.
\end{toyrun}

%% =====================================================================================
\section{Linear claims and the zero-knowledge opening (paper \S4.2--4.3)}
\label{sec:opening}
%% =====================================================================================

\subsection{What a claim is}

A \emph{claim} is an equation
\[
  \sum_t \mathit{coef}_t\cdot\ip{\mathit{row}_t}{L_t}=\mathit{value}
\]
between committed rows (their first $C$ data symbols), public vectors $L_t\in\F^C$, public coefficients
and a value in $\F$. An evaluation of a multilinear extension is such a claim: if a table is split into
rows of length $C$ (low $\log C$ variables index the position in the row, high variables index the
row), then
\[
  \mle T(x_{\mathrm{low}},x_{\mathrm{high}})=\sum_{\text{rows } t}\underbrace{\eq(t,x_{\mathrm{high}})}_{\mathit{coef}_t}
  \;\ip{\mathit{row}_t}{\underbrace{(\eq(y,x_{\mathrm{low}}))_{y<C}}_{L}} .
\]
Rows that are paired with the same vector $L$ form a \emph{slot} $v$ with vector $L_v$ (not to be confused
with CRT slots). For every slot the prover committed a uniformly random \emph{mask row} $m_v\in\F^k$,
and for every matrix one more mask row $m_{\mathsf{prox}}$.

\subsection{The protocol, step by step}

\begin{enumerate}[leftmargin=1.5em]
\item \textbf{Batching.} $V$ sends $\kappa_i\in\F$ for every claim $i$. Define for each slot the combined row
  $u_v=\sum_i\kappa_i\sum_{t\in v}\mathit{coef}_{i,t}\,\mathit{row}_{i,t}$ and the batched value
  $V=\sum_i\kappa_i\,\mathit{value}_i$. If all claims hold, $\sum_v\ip{u_v}{L_v}=V$.
\item \textbf{Mask commitment value.} $P$ sends $\mu=\sum_v\ip{m_v}{L_v}$.
\item \textbf{Challenges.} $V$ sends $c\in\F$, and $\gamma_i\in\F$ for every row $i$ of every matrix.
\item \textbf{Vectors.} $P$ sends $\bar u_v=m_v+c\,u_v$ for every slot (all $k$ coordinates, including
  the pad positions), and $p=m_{\mathsf{prox}}+\sum_i\gamma_i\,\mathit{row}_i$ for every matrix.
\item \textbf{Value check.} $V$ checks $\sum_v\ip{\bar u_v}{L_v}=\mu+cV$ (inner product on the first $C$
  coordinates).
\item \textbf{Column queries.} $V$ sends $\nq$ distinct column indices per matrix; $P$ opens these columns
  with salts and Merkle paths.
\item \textbf{Consistency checks.} For every opened column $j$, $V$ encodes $\bar u_v$ and $p$ and checks
  \begin{align*}
   \RS(\bar u_v)[j]&=\mathrm{col}_j[m_v]+c\sum_i\kappa_i\sum_{t\in v}\mathit{coef}_{i,t}\,\mathrm{col}_j[\mathit{row}_{i,t}],\\
   \RS(p)[j]&=\mathrm{col}_j[m_{\mathsf{prox}}]+\sum_i\gamma_i\,\mathrm{col}_j[i].
  \end{align*}
\end{enumerate}

\begin{example}[An opening with one row and one claim, modulo 17]
Code: $n=8$, $\omega=2$, $C=2$, one pad, $k=3$. Committed row $\mathit{row}=(4,9,6)$ (data $4,9$, pad $6$),
codeword $(2,12,0,1,1,10,13,10)$. Claim: $\ip{\mathit{row}}{L}=V$ with $L=(2,3)$, so
$V=8+27=35\equiv 1$. Mask row $m=(5,1,12)$, codeword $\RS(m)=(1,4,14,16,16,0,6,0)$. (One claim, so we skip
$\kappa$ and the proximity vector.)
\begin{itemize}[leftmargin=1.2em]
\item $P$ sends $\mu=\ip{(5,1)}{(2,3)}=13$.
\item $V$ sends $c=10$.
\item $P$ sends $\bar u=m+c\cdot\mathit{row}=(5+40,\,1+90,\,12+60)\equiv(11,6,4)$.
\item Value check: $\ip{(11,6)}{(2,3)}=40\equiv 6$ and $\mu+cV=13+10=23\equiv 6$. \checkmark
\item $V$ queries column 3. Opened entries: $\mathrm{col}_3[m]=16$, $\mathrm{col}_3[\mathit{row}]=1$.
  $\RS(\bar u)=(4,5,14,9,9,15,0,15)$, so $\RS(\bar u)[3]=9$, and $16+10\cdot 1=26\equiv 9$. \checkmark
\end{itemize}
\emph{A cheater} who claims $V=2$ must pass $\ip{\bar u'}{L}=\mu+2c=33\equiv 16$, e.g.\ with
$\bar u'=(16,6,4)$. Its encoding $(9,10,2,14,14,3,5,3)$ agrees with the true combination
$\RS(m)+c\RS(\mathit{row})=(4,5,14,9,9,15,0,15)$ in \emph{none} of the 8 columns: every query catches it.
(It could instead have changed $\mu$, but $\mu$ is sent before $c$; see below.)
\end{example}

\subsection{Why each step is there}

\begin{description}[leftmargin=1em,itemsep=2pt]
\item[$\kappa$ (batching).] Proving one random combination of all claims is as good as proving each:
  if some claim is false, $\sum_i\kappa_i(\text{true}_i-\text{claimed}_i)$ is a non-zero linear form in
  $\kappa$, zero with probability $1/|\F|$ (budget line ``a false claim disappears in the batch'').
\item[$\mu$ and $c$.] Sending $u_v$ itself would leak the witness. Sending $\bar u_v=m_v+c\,u_v$ with a
  uniform $m_v$ leaks nothing (it is uniform). But then the verifier needs $\ip{m_v}{L_v}$ to subtract, so the
  prover announces $\mu$ \emph{before} $c$. If the honest combination is sent, the value check reads
  $\mu'+cV'=\mu+cV$ where $\mu',V'$ are the true values; with $\mu,\mu',V,V'$ fixed before $c$ and $V\neq V'$
  this holds for only one $c$ (budget line ``opening equation holds for a wrong $V$'').
\item[$\gamma$ and $p$ (proximity test).] The prover could commit rows that are \emph{not} codewords.
  Then ``the row'' is ambiguous and the column checks prove little. The random combination $p$ of
  \emph{all} rows of a matrix is sent and encoded; if the committed rows are far from the code, so is
  their random combination (Lemma~\ref{lem:prox}), and the column checks on $p$ fail on many columns.
\item[Column queries.] Tie the sent vectors to the commitment. A codeword that is wrong is wrong on at
  least $d-e$ columns (Lemma~\ref{lem:cons}).
\item[All $k$ coordinates of $\bar u_v$.] The verifier must re-encode $\bar u_v$, which needs the pad
  coordinates too; the inner product only uses the first $C$.
\end{description}

\subsection{The four lemmas}
\label{sec:opening-lemmas}

\begin{lemma}[Binding; paper 4.1]
In the ROM a Merkle root fixes at most one value per leaf, except with probability $t^2/2^{256}$ for $t$
oracle queries.
\end{lemma}

\begin{lemma}[Proximity; paper 4.2]\label{lem:prox}
Let $M$ be the rows of a matrix other than $m_{\mathsf{prox}}$. Suppose there is no set of $n-e$ columns
on which every row of $M$ agrees with a codeword. Then for every $m_{\mathsf{prox}}$,
$\Pr_\gamma[\,m_{\mathsf{prox}}+\sum_i\gamma_i\mathit{row}_i\text{ is within distance }e\text{ of the code over }\F\,]\le\epsprox$,
with $\epsprox=(e+2)/|\F|$ for $\Delta=3$ and $n/|\F|$ for $\Delta=2$.
\end{lemma}

In words: if the rows cannot be ``explained'' by codewords with a common set of at most $e$ bad columns,
a random combination is far from the code. Notes on this lemma:
\begin{itemize}[leftmargin=1.3em]
\item ``Agree with a codeword on a set of columns'' with a \emph{common} set for all rows is
  \emph{correlated agreement}. It is what the extractor needs: one set of $\le e$ bad columns outside
  which all rows are decodable.
\item $\Delta=2$: the proximity gap for Reed--Solomon codes at the unique-decoding radius
  (Ben-Sasson, Carmon, Ishai, Kopparty, Saraf 2020, Thm.~1.6), error $n/|\F|$.
\item $\Delta=3$: an elementary lemma about \emph{lines} $x+\gamma y$, valid for any linear code
  (Roth--Z\'emor, stated as Diamond--Posen Thm.~2.1), error $(e+1)/|\F|$. The paper extends it from lines
  to affine spaces $m_{\mathsf{prox}}+\mathrm{span}(\text{rows})$ in its appendix
  (Section~\ref{sec:appendix} below), getting $(e+2)/|\F|$.
\item Rows are in $\Zq$ but the test is over $\F$. If a $\Zq$-row agrees with an $\F$-codeword on more
  than $k$ columns, it agrees with a $\Zq$-codeword there: write the $\F$-codeword as $c_0+c_1t$; both
  $c_0,c_1$ are $\Zq$-codewords, the row equals $c_0$ on those columns and $c_1$ is zero there.
\end{itemize}

\begin{lemma}[Consistency; paper 4.3]\label{lem:cons}
Suppose the rows agree with codewords outside a set of at most $e$ columns, and some $\bar u_v$ differs
from $m_v+c\,u_v$ computed on the decoded rows. Then a uniformly random column passes the check of step 7
with probability at most $1-(d-e)/n$.
\end{lemma}
\begin{proof}
$\RS(\bar u_v)$ and the encoding of the decoded combination are distinct codewords: they differ in at
least $d$ columns. On all but at most $e$ of those, the committed entries equal the decoded ones, so the
check fails on at least $d-e$ columns.
\end{proof}
Since $d-1\ge\Delta e\ge 2e$, we have $d-e>e$, so this is at most $1-e/n$ per query; $\nq$ independent
queries all miss with probability at most $(1-e/n)^{\nq}\le 2^{-128}$.

\begin{lemma}[Hiding; paper 4.4]
For every set $J$ of at most $\nq$ columns, the entries of a row at $J$ are uniform and independent of
its data symbols.
\end{lemma}
This is Section~\ref{sec:hiding}. Together with the uniform mask rows, it makes the whole opening
simulatable.

%% =====================================================================================
\section{The shuffle protocol, message by message (paper \S5.1, Figure 1)}
\label{sec:protocol}
%% =====================================================================================

We follow the seven steps of the paper's Figure~1. Each step: what is sent, why, and the toy run.
In the toy run the prover's randomness (masks, pads, salts) is first set to \emph{zero} so that the
algebra is visible (the ``bare run''); Section~\ref{sec:zk} switches it on. Notation: $\ell=\log_2\tau$
(ballot-index variables), $m=\ell+\log_2N$ (variables of a witness table). Toy: $\ell=1$, $\log N=2$, $m=3$.

\subsection*{Step 1. Commit the witness}
$P$ commits $W$ (big matrix) and $\sigma$ (small matrix, batch 1), each with its mask rows, and sends the
two roots. The challenges that follow are therefore independent of the witness.

\subsection*{Step 2. Projection challenge}
$V$ sends $\rho\in\F^{2N}$ and $\alpha,\beta\in\F$. Both sides compute $\rho'=\rho\Mpk$,
$\zpub_i=\ip{\rho}{\coeff(c_i)}$ and $\zh_j=\ip{\rho}{\coeff(\ch_j)}$. Only $P$ computes
\[
  z_i=\zpub_i+\ip{\rho'}{\coeff(w_i)},\quad h_i=z_i+\beta i-\alpha,\quad \hh_j=\zh_j+\beta\sigma_j-\alpha,\quad
  u_i=1/h_i,\quad \uh_j=1/\hh_j .
\]
Toy values: Sections~\ref{sec:projection} and~\ref{sec:logder}.

\subsection*{Step 3. Commit the inverses and the masks}
$P$ commits $u$, $\uh$ and the mask table (small matrix, batch 2) and sends the root and the mask sums
$G_A,G_B,G_C$. (Toy bare run: $G=(0,0,0)$.)

\subsection*{Step 4. Sumcheck A: the permutation}

$V$ sends $t_c\in\F^\ell$ and $k_1,k_2,k_3,\rho_A\in\F$. The claim is
\[
  \sum_{x\in\cube{\ell}}\Bigl[\eq(t_c,x)\bigl(k_1\,\msk u\msk h+k_2\,\msk\uh\msk\hh\bigr)+k_3(\msk u-\msk\uh)+\rho_Ag_A\Bigr](x)
  \;=\;k_1+k_2+\rho_AG_A,
\]
with $\msk\hh(x)=\mle{\zh}(x)+\beta\msk\sigma(x)-\alpha$.

\paragraph{Why this is the right claim.} On the cube, masked extensions equal the tables. If the three
facts \eqref{eq:threefacts} hold, then $\sum_x\eq(t_c,x)(k_1\cdot 1+k_2\cdot 1)=k_1+k_2$ (because
$\sum_x\eq(t_c,x)=1$), $k_3\sum_x(u-\uh)=0$, and $\rho_A\sum g_A=\rho_AG_A$. Conversely, subtracting,
the claim says
\[
  \sum_x\eq(t_c,x)\bigl[k_1(uh-1)+k_2(\uh\hh-1)\bigr](x)+k_3\textstyle\sum_x(u-\uh)(x)=0 ,
\]
and for random $t_c,k_1,k_2,k_3$ this forces $uh=1$, $\uh\hh=1$ everywhere and $\sum u=\sum\uh$
(Section~\ref{sec:mle}, ``Why random evaluations''), except with probability $(\ell+1)/|\F|$.

\paragraph{Degrees.} $\eq$ has degree 1 per variable; unmasked extensions also 1, so
$\eq\cdot u\cdot h$ has degree 3. Masked extensions have degree 2 per variable, but $Z(x)$ vanishes as
soon as \emph{one} variable is Boolean. In round $j<\ell-1$ the later variables are summed over $\{0,1\}$,
so all mask terms vanish and the round polynomial has degree 3. Only in the \emph{last} round are all other
variables non-Boolean challenges; then degree $1+2+2=5$. So $P$ sends 3 values per round
($p(0),p(2),p(3)$; $p(1)$ is implied) and 5 values in the last round ($p(0),p(2),\dots,p(5)$).

\paragraph{End.} The sumcheck ends at a point $r\in\F^\ell$. $P$ sends
\[
  (U,\hat U,s,H)=(\msk u,\msk\uh,\msk\sigma,\msk h)(r).
\]
$V$ computes $\msk\hh(r)=\mle\zh(r)+\beta s-\alpha$ itself and evaluates the bracket \emph{without} the
mask term from these values; the final claim must equal that plus $\rho_Ag_A(r)$, so $V$ \emph{derives}
the value $\gamma_A=(\text{final claim}-\text{bracket})/\rho_A$ that $g_A(r)$ must take. Now five values are
owed: $U,\hat U,s$ (evaluations of committed tables), $H$ (of the \emph{uncommitted} $h$), and $g_A(r)=\gamma_A$.

\begin{toyrun}[Sumcheck A, one round ($\ell=1$, so the first round is the last: degree 5)]
$t_c=6292{+}4021t$, $k_1=6152{+}6333t$, $k_2=4804{+}3613t$, $k_3=1965{+}20t$, $\rho_A=5031{+}661t$.
Claim $k_1+k_2+\rho_A\cdot 0=3275{+}2265t$.
\begin{itemize}[leftmargin=1.2em]
\item $P$ sends $p(0)=128{+}5163t$ and $p(2),\dots,p(5)=5168{+}4781t,\ 228{+}4632t,\ 5407{+}3811t,\ 7061{+}1793t$.
\item $V$ sets $p(1)=\text{claim}-p(0)=3147{+}4783t$, draws $r=907{+}2353t$, interpolates: $p(r)=4931{+}4854t$.
\item $P$ sends $U=5474{+}7238t$, $\hat U=881{+}4663t$, $s=6775{+}5328t$, $H=6511{+}4532t$
  (bare run: these are the plain extensions, e.g.\ $U=(1-r)u_0+ru_1$).
\item $V$: $\msk\hh(r)=\mle\zh(r)+\beta s-\alpha=4068{+}3441t$, $\eq(t_c,r)=5237{+}6029t$,
  bracket $=\eq(t_c,r)(k_1UH+k_2\hat U\msk\hh(r))+k_3(U-\hat U)=4931{+}4854t=p(r)$, so $\gamma_A=0$.
\end{itemize}
\end{toyrun}

\subsection*{Step 5. Sumcheck B: $H$ is the right value of $h$}

$h$ was never committed, so $H$ is only a claim. But by definition
$h_i=\zpub_i+\sum_k\ip{\rho'_k}{W_k[\,\cdot+Ni\,]}+\beta i-\alpha$. Take the extension in the ballot index
$i$ and evaluate at $r$, using $\sum_i\eq(i,r)=1$:
\[
  \msk h(r)=\mle{\zpub}(r)+\sum_{y<N}\sum_{k=1}^3\rho'_k(y)\,\mle W_k(y,r)+\beta\,\id(r)-\alpha+a_hZ(r),
\]
where $\mle W_k(y,r)$ is the table $W_k$ folded at $r$ in the ballot index (Section~\ref{sec:mle}) and
\[
  \id(r)=\sum_{i<\ell}2^ir_i
\]
is the extension of the table $(0,1,2,\dots,\tau-1)$, which is multilinear already ($i=\sum_j2^jb_j$).
Solving for the sum gives the claim of sumcheck~B (over $y\in\cube{\log N}$, degree 2):
\[
  \sum_y\Bigl[\sum_{k=1}^3\rho'_k(y)\,\mle W_k(y,r)+\frac{a_hZ(r)}{N}+\rho_Bg_B(y)\Bigr]
  =H-\mle{\zpub}(r)-\beta\,\id(r)+\alpha+\rho_BG_B .
\]
The constant $a_hZ(r)$ is ``spread'' as $a_hZ(r)/N$ over the $N$ points so that it sums to $a_hZ(r)$.
Everything on the right is known to $V$. Degree 2: $\rho'_k(y)$ (as an extension in $y$) times
$\mle W_k(y,r)$, each of degree 1 in $y$.

At the end point $r_y$, the final claim $\mathit{cl}_B$ must equal
$\sum_k\mle{\rho'_k}(r_y)\mle W_k(r_y,r)+a_hZ(r)/N+\rho_Bg_B(r_y)$. $V$ computes $\mle{\rho'_k}(r_y)$ and
$Z(r)/N$ itself, so this is \emph{one linear claim} about the rows of $W$, the entry $a_h$ and the
coefficients of $g_B$.

\begin{toyrun}[Sumcheck B ($\log N=2$ rounds)]
Claim: $H-\mle\zpub(r)-\beta r+\alpha=4510{+}2479t$. Checking by brute force over $y=0,1,2,3$,
$\sum_k\rho'_k(y)\mle W_k(y,r)$ is $2848{+}6920t$, $6943{+}1195t$, $6482{+}5565t$, $3599{+}4161t$, summing to
$4510{+}2479t$. \checkmark{} ($\rho_B=6690{+}803t$, unused since $g_B=0$ in the bare run.)

Round 0: $p(0)=1649{+}4804t$, $p(1)=2861{+}5356t$, $p(2)=3730{+}3059t$, challenge $3683{+}94t$, new claim
$6{+}2886t$. Round 1: $p(0)=510{+}7412t$, $p(1)=7177{+}3155t$, $p(2)=1327{+}6685t$, challenge
$6674{+}7298t$, new claim $\mathit{cl}_B=5341{+}3708t$ at $r_y=(3683{+}94t,\ 6674{+}7298t)$.
(Here all three values are printed; on the wire $p(1)$ is omitted.)
\end{toyrun}

\subsection*{Step 6. Sumcheck C: the witness is ternary}

$V$ sends $t_z\in\F^m$ and $k^C_1,k^C_2,k^C_3,\rho_C\in\F$. Claim (over $x\in\cube m$, degree 4, last round 7):
\[
  \sum_x\Bigl[\eq(t_z,x)\sum_{k=1}^3k^C_k\bigl((\msk w_k)^3-\msk w_k\bigr)+\rho_Cg_C\Bigr](x)=\rho_CG_C .
\]
On the cube $(\msk w)^3-\msk w=w^3-w$, which is $0$ for ternary entries. The random $\eq(t_z,\cdot)$ and $k^C$
turn ``all $3\tau N$ values are zero'' into one sum (error $(m+1)/|\F|$). Degrees: $1+3=4$ in normal rounds,
$1+3\cdot 2=7$ in the last. At the end point $r_z$, $P$ sends $\msk w_k(r_z)$ for $k=1,2,3$, and $V$ derives
the value $\gamma_C$ that $g_C(r_z)$ must take.

\begin{toyrun}[Sumcheck C ($m=3$ rounds)]
$t_z=(5612{+}4016t,\ 5565{+}2574t,\ 1722{+}3254t)$, $k^C=(2060{+}2848t,\ 7612{+}2921t,\ 6595{+}3083t)$.
\begin{itemize}[leftmargin=1.2em]
\item Round 0: $p(0)=0$, $p(1)=0$ (all entries ternary, and the later variables are Boolean), $p(2),p(3),p(4)=
  5530{+}2379t,\ 7607{+}3239t,\ 4306{+}4510t$; challenge $5224{+}625t$; new claim $4055{+}6605t$.
\item Round 1: $p(0)=2763{+}1802t$, $p(1)=1292{+}4803t$, three more values; challenge $5923{+}2788t$;
  new claim $941{+}3043t$.
\item Round 2 (last, degree 7): $p(0)=297{+}5614t$, $p(1)=644{+}5110t$, six more values;
  challenge $734{+}4566t$; new claim $3656{+}6497t$.
\item $P$ sends $\mle w_k(r_z)=(555{+}1978t,\ 7531{+}6763t,\ 891{+}2913t)$; $V$ computes
  $\eq(t_z,r_z)\sum_kk^C_k(w^3-w)=3656{+}6497t$, equal to the last claim, so $\gamma_C=0$.
\end{itemize}
\end{toyrun}

\subsection*{Step 7. One opening for nine claims}

\begin{center}
\begin{tabular}{@{}rlll@{}}
\toprule
\# & claim & about & toy value (bare)\\\midrule
1 & $\msk u(r)=U$ & row $u$, $a_u$ & $5474{+}7238t$\\
2 & $\msk\uh(r)=\hat U$ & row $\uh$, $a_{\uh}$ & $881{+}4663t$\\
3 & $\msk\sigma(r)=s$ & row $\sigma$, $a_\sigma$ & $6775{+}5328t$\\
4 & final value of B $=\mathit{cl}_B$ & rows of $W$ at $(r_y,r)$, $a_h$, $g_B$ & $5341{+}3708t$\\
5--7 & $\msk w_k(r_z)$ & rows of $W$ at $r_z$, $a_{w_k}$ & $555{+}1978t$, $7531{+}6763t$, $891{+}2913t$\\
8 & $g_A(r)=\gamma_A$ & coefficients of $g_A$ & $0$\\
9 & $g_C(r_z)=\gamma_C$ & coefficients of $g_C$ & $0$\\
\bottomrule
\end{tabular}
\end{center}
Every claim is linear in committed data: e.g.\ $\msk u(r)=\mle u(r)+a_uZ(r)$, where $\mle u(r)$ is a
combination of the entries of row $u$ and $a_u$ is an entry of the mask-table row.

\begin{toyrun}[The opening]
$\mu=0$ (bare), $c=742{+}3266t$. Slot vectors: two for the big matrix ($W$ at $(r_y,r)$ and at $r_z$; 11
entries each $=C_1+\nq$), two for the small matrix (131 entries), one proximity vector per matrix.
Queried columns: big $\{9,30,14\}$, small $\{124,467,480\}$. Column 9 of the big matrix has 9 entries
(rows $W_r,W_{e_1},W_{e_2}$ and six mask rows): $(3078,3182,1414,0,0,0,0,0,0)$. Encoding of the proximity
vector at position 9: $6991{+}4583t$; mask entries $+\sum_i\gamma_i\,\mathrm{col}_9[i]=6991{+}4583t$. \checkmark

\textbf{Verdict:} \texttt{mix\_verify} accepts.
\end{toyrun}

\begin{recap}[Where each fact is proven]
permutation and linear part $\Leftarrow$ log-derivative lemma $\Leftarrow$ three facts \eqref{eq:threefacts}
$\Leftarrow$ sumcheck A (+ B for $H$) $\Leftarrow$ claims 1--4, 8 $\Leftarrow$ opening.
Shortness $\Leftarrow$ sumcheck C $\Leftarrow$ claims 5--7, 9 $\Leftarrow$ opening.
Opening $\Leftarrow$ batching, value check, proximity test, column queries $\Leftarrow$ Merkle binding.
\end{recap}

%% =====================================================================================
\section{Knowledge soundness (paper \S5.3, Theorem 5.2)}
\label{sec:sound}
%% =====================================================================================

\begin{theorem}[Knowledge soundness, interactive; paper Theorem 5.2]
Consider the protocol with the matrices as oracles, and $\tau<q$. Decode every committed row, at the
time its batch is committed, to the unique codeword within distance $e$ if there is one. If the decoded
$(\sigma,W)$ is not a witness, the verifier accepts with probability at most
\[
  \frac{2\tau+9\ell+7\log N+9}{|\F|}+2\epsprox+\Bigl(1-\frac en\Bigr)^{\nq}.
\]
\end{theorem}

\paragraph{How to read it.} ``Decode at the time the batch is committed'' defines the extracted witness
\emph{before} the challenges that depend on it, so the prover cannot adapt the witness to the challenges.
The theorem says: either the decoded rows are a valid witness (so the extractor succeeds), or the
prover is caught except with the stated probability.

\subsection{The error budget, line by line}

Each line is an event under which one verifier message lets a false claim survive.

\begin{center}\small
\begin{tabular}{@{}p{0.43\textwidth}lp{0.3\textwidth}@{}}
\toprule
event & probability & where it comes from\\\midrule
log-derivative identity holds for two different multisets & $(2\tau-1)/|\F|$ & Lemma~\ref{lem:logder}\\
claim of A true although an inverse or the sum is wrong & $(\ell+1)/|\F|$ & a non-zero polynomial of degree $\ell+1$ in $(t_c,k,\rho_A)$\\
sumcheck A accepts a false claim & $(3\ell+2)/|\F|$ & degrees $3(\ell-1)+5$\\
$H$ wrong although the claim of B is true & $1/|\F|$ & $(\msk h(r)-H)+\rho_B(G'_B-G_B)=0$ for one $\rho_B$\\
sumcheck B accepts a false claim & $2\log N/|\F|$ & $\log N$ rounds of degree 2\\
claim of C true although $W$ is not ternary & $(m+1)/|\F|$ & degree $m+1$ in $(t_z,k^C,\rho_C)$\\
sumcheck C accepts a false claim & $(4m+3)/|\F|$ & degrees $4(m-1)+7$\\
a false claim disappears in the batch & $1/|\F|$ & linear in $\kappa$\\
opening equation holds for a wrong $V$ & $1/|\F|$ & linear in $c$\\
a far matrix yields a close combination & $\epsprox$ per matrix & Lemma~\ref{lem:prox}, 2 matrices\\
all column queries miss & $(1-e/n)^{\nq}$ & Lemma~\ref{lem:cons}\\
\bottomrule
\end{tabular}
\end{center}

\paragraph{Adding up.} $(2\tau-1)+(\ell+1)+(3\ell+2)+1+2\log N+(m+1)+(4m+3)+1+1=2\tau+4\ell+2\log N+5m+9$,
and with $m=\ell+\log N$ this is $2\tau+9\ell+7\log N+9$. \checkmark

\begin{example}[The bound at real parameters]
$\tau=4096$, $\ell=\log N=12$, $m=24$, $|\F|\approx 2^{156}$:
\begin{itemize}[leftmargin=1.2em]
\item $(2\cdot 4096+108+84+9)/2^{156}=8393/2^{156}\approx 2^{-143}$;
\item $\Delta=3$: $\epsprox=(e+2)/|\F|=65435/2^{156}\approx 2^{-140}$, twice that $\approx 2^{-139}$;
\item queries: $(1-e/n)^{\nq}\approx 2^{-128.01}$.
\end{itemize}
The total is about $2^{-128}$, dominated by the column queries, by design.
\end{example}

\subsection{The proof, walked backwards}

Assume none of the events happens and $V$ accepts. We show the decoded rows form a witness.
\begin{enumerate}[leftmargin=1.5em]
\item \textbf{Opening.} If a matrix had no decoding with $\le e$ bad columns, $p$ would be far from the code
  (Lemma~\ref{lem:prox}). The vector $p$ that $P$ sends is encoded by $V$, so it is a codeword; it differs
  from the committed combination on more than $e$ columns, and some query would hit one. So both matrices
  decode. By Lemma~\ref{lem:cons} every $\bar u_v$ equals $m_v+c\,u_v$ on the decoded rows. The value check
  reads $\mu'+cV'=\mu+cV$ with $\mu',V'$ computed from the decoded rows and fixed before $c$; so $V'=V$;
  and since batching did not hide a false claim, \emph{all nine claims hold for the decoded rows}.
\item \textbf{Sumcheck B.} Claim 4 says the final value of B is the true summand at $r_y$. Since sumcheck B
  did not accept a false claim, its initial claim is true:
  $(\msk h(r)-H)+\rho_B(G'_B-G_B)=0$ ($G'_B$ = true sum of the committed $g_B$). Both differences were fixed
  before $\rho_B$, so (no event) $H=\msk h(r)$.
\item \textbf{Sumcheck A.} Now $U,\hat U,s,H$ and $g_A(r)$ are all correct, so the last check of A is the true
  summand at $r$; the claimed sum is therefore true; so $uh=1$, $\uh\hh=1$ on the cube and $\sum u=\sum\uh$,
  i.e.\ $\sum_i1/(z_i+\beta i-\alpha)=\sum_j1/(\zh_j+\beta\sigma_j-\alpha)$ with no zero denominator.
\item \textbf{Multisets.} By Lemma~\ref{lem:logder} the tagged multisets are equal; $\sigma$ is a permutation and
  $\ch_j=c_{\sigma(j)}+\EncZ(w_{\sigma(j)})$ (Section~\ref{sec:logder}, ``What the lemma gives'').
\item \textbf{Sumcheck C.} As for A: $w^3=w$ for every entry, so every entry is in $\{-1,0,1\}$. \qed
\end{enumerate}
The extractor is \emph{straight-line}: read the committed leaves off the prover's hash queries and run
Reed--Solomon decoding on the first message. No rewinding.

\subsection{The cheaters in action}
\label{sec:cheat}\label{sec:toy-cheat}

\begin{toyrun}[False statements are rejected]
\begin{center}\small
\begin{tabular}{@{}lll@{}}
\toprule
 & (a) one plaintext bit of output 0 flipped & (b) slot attack\\\midrule
outputs decrypt to & $(0,1,1,1)$, $(0,0,0,1)$ & $(1,1,1,0)$, $(0,0,0,0)$\\
product check in $\Rq$ accepts & 0 of 20 trials & \textbf{20 of 20 trials}\\
$\sum 1/h$ (inputs) & $6355{+}4220t$ & $6355{+}4220t$\\
$\sum 1/\hh$ (outputs) & $3169{+}2227t$ & $3408{+}4838t$\\
prover run on the false statement & rejected (opening) & rejected (opening)\\
honest proof replayed on it & rejected (opening) & rejected (opening)\\
\bottomrule
\end{tabular}
\end{center}
And a re-randomiser with coefficient $2$ (outputs computed consistently, so the linear part and the
permutation are fine): $2^3-2=6\neq 0$, the sum in sumcheck C is not $0$, rejected at the opening.
A changed entry in an opened column: rejected (Merkle path).
\end{toyrun}

\begin{why}[Why are all rejections reported ``at the opening''?]
Sumcheck messages always look consistent, because $p(1)$ is derived from the running claim. A false
claim is \emph{pushed} through the sumchecks until it becomes a false claimed evaluation of a committed
table (e.g.\ $g_A(r)\neq\gamma_A$), and only the opening can expose that.
\end{why}

%% =====================================================================================
\section{Zero knowledge (paper \S5.4, Theorem 5.3)}
\label{sec:zk}
%% =====================================================================================

\paragraph{Good challenges.} The simulator works when: no $h_i$ or $\hh_j$ is $0$ (inverses exist),
$Z(r)\neq 0$ and $Z(r_z)\neq 0$ (masks are effective), $\rho_A,\rho_B,\rho_C\neq 0$ (sumcheck masks are
effective). Bad challenges occur with probability at most $(2\tau+2\ell+2m+3)/|\F|$: each $h_i$, $\hh_j$ is
zero for one $\alpha$ ($2\tau$); $Z(r)=0$ iff some $r_i\in\{0,1\}$ ($2\ell$); same for $r_z$ ($2m$); three $\rho$'s.

\begin{theorem}[Honest-verifier zero knowledge; paper Theorem 5.3]
There is a simulator that, given a true statement and good challenges, outputs a transcript in which every
value other than the hashes of unopened leaves has exactly the distribution of a real transcript. The
hashes differ from real ones with probability at most $t/2^{128}$ for a distinguisher with $t$ queries.
\end{theorem}

\paragraph{The simulator and why each step is exact.}
\begin{enumerate}[leftmargin=1.5em]
\item Draw $U,\hat U,s,H$ and $\msk w_k(r_z)$ uniformly. \emph{Real:} each is $\mle T(\text{pt})+a_TZ(\text{pt})$
  with its own uniform $a_T$ used at one point only and $Z(\text{pt})\neq 0$; so each is uniform and
  independent (Section~\ref{sec:masked}).
\item Draw $G_A,G_B,G_C$ uniformly and every sumcheck round message uniformly among those consistent with
  the running claim. \emph{Real:} in round $j$ the non-constant coefficients of the mask's round polynomial
  are $2^{m-j-1}c_{j,t}$, and $G=2^ma_0+(\text{function of the }c_{j,t})$. So the map from the mask's
  coefficients to $(G,\text{round messages})$ is a triangular affine bijection; with $\rho\neq 0$ the
  messages are uniform subject to the verifier's checks.
\item Derive $\mathit{cl}_B,\gamma_A,\gamma_C$ and $V$ with the verifier's own formulas (same in both worlds).
\item Draw all $\bar u_v$ and $p$ uniformly; set $\mu=\sum_v\ip{\bar u_v}{L_v}-cV$. \emph{Real:} $m_v$ and
  $m_{\mathsf{prox}}$ are uniform and used once, so $\bar u_v$, $p$ are uniform, and $\mu$ is then fixed by the
  value check.
\item In the opened columns, draw data-row entries uniformly (\emph{real:} Lemma 4.4) and solve the mask-row
  entries from the step-7 checks (\emph{real:} determined by the same checks).
\item Choose salts, fill unopened leaves with random hashes, build the trees (\emph{real:} salted hashes of
  unopened leaves look random).
\end{enumerate}

\begin{toyrun}[Bare, real (masked) and simulated transcripts]
Same statement and challenges. The simulator is \emph{not} given $\sigma$ or $w$.
\begin{center}\small
\begin{tabular}{@{}llll@{}}
\toprule
value & bare (no randomness) & real prover & simulator\\\midrule
$G_A$ & $0$ & $499{+}2211t$ & $3792{+}5023t$\\
first message of A, $p(0)$ & $128{+}5163t$ & $997{+}2287t$ & $335{+}5965t$\\
$U$ & $5474{+}7238t$ & $513{+}176t$ & $7098{+}5543t$\\
$s$ (was $\mle\sigma(r)$) & $6775{+}5328t$ & $6669{+}5028t$ & $4118{+}3798t$\\
$H$ & $6511{+}4532t$ & $6541{+}5163t$ & $7358{+}4953t$\\
$\mu$ & $0$ & $4815{+}7202t$ & $4275{+}5353t$\\
\bottomrule
\end{tabular}
\end{center}
All three are accepted. In the bare run $s=\mle\sigma(r)=(1-r)\sigma_0+r\sigma_1$ would reveal $\sigma$ to anyone
who knows $r$; in the real run it is uniform.
\end{toyrun}

\begin{why}[Isn't the simulator a forgery?]
No. It needs the challenges \emph{before} writing its messages. In the non-interactive proof the challenges are
hashes of those messages (Section~\ref{sec:fs}), so this order is impossible for a real prover.
\end{why}

%% =====================================================================================
\section{Non-interactive proofs (paper \S6)}
\label{sec:fs}
%% =====================================================================================

\paragraph{The transformation.} Oracles become Merkle roots, and every challenge is a hash of the
parameters, the statement and all prover messages so far. In the implementation (\texttt{docs/SPEC.md}~\S4)
the transcript is a hash chain. The state starts as $H(\texttt{"FP-BGV-v2|mix"})$. Absorbing a prover message
$\mathit{data}$ with a label, and drawing a challenge called $\mathit{name}$, update it as
\begin{align*}
  \mathit{state}&\leftarrow H\bigl(\mathit{state}\,\|\,\texttt{"|A|"}\,\|\,\mathit{label}\,\|\,\texttt{"|"}\,\|\,H(\mathit{data})\bigr),\\
  \mathit{state}&\leftarrow H\bigl(\mathit{state}\,\|\,\texttt{"|C|"}\,\|\,\mathit{name}\bigr),
\end{align*}
and the challenge is read from the blocks $H(\mathit{state}\,\|\,\texttt{"|S|"}\,\|\,\mathit{counter})$:
$\Zq$-elements by rejection sampling (10 bytes masked to 79 bits, retried if $\ge q$), $\F$-elements as two
$\Zq$-elements, column indices the same way, repeated until distinct. Because every challenge depends on
everything absorbed before it, the prover cannot choose a message after seeing the challenge it produces.

\paragraph{Doomed transcripts.} A partial transcript is \emph{doomed} if a committed row has no decoding and
the proximity test can still catch it, or if the claim currently being reduced is false for the decoded rows.
\begin{itemize}[leftmargin=1.3em]
\item The empty transcript is doomed if the decoded first message is not a witness.
\item A complete doomed transcript is rejected.
\item A verifier message rescues a doomed transcript only through its event in the budget table. The
  challenges $c$ and $\gamma$ form one message; their two events are mutually exclusive.
\end{itemize}

\begin{proposition}[Round-by-round soundness, sketch; paper Proposition 6.1]
\[
  \epsrbr=\max\Bigl\{\frac{\max(2\tau-1,\,m+1,\,7)}{|\F|},\ \epsprox,\ \Bigl(1-\frac en\Bigr)^{\nq}\Bigr\}\le 2^{-128}.
\]
\end{proposition}
Each term is the worst single message: the projection challenge ($2\tau-1$), a batching challenge such as
$(t_z,k^C,\rho_C)$ ($m+1$), one sumcheck round of degree at most 7, the proximity challenge, the queries.
The authors label it a sketch because the state function is not written out round by round.

\begin{theorem}[Non-interactive security; paper Theorem 6.2]
In the ROM the non-interactive proof has knowledge error at most $t\,\epsrbr+O(t^2/2^{256})$ against $t$
classical queries, and $O(t^2\epsrbr+t^3/2^{256})$ against $t$ quantum queries. It is statistically zero
knowledge.
\end{theorem}

\begin{caution}[What this does and does not say]
With $\epsrbr\approx 2^{-128}$, the quantum bound $t^2\epsrbr$ is below 1 only for $t<2^{64}$. The quantum
statement assumes the theorem of Chiesa--Manohar--Spooner applies to this exact hashing layout (salted
leaves, hash-chain transcript, rejection sampling), which the authors flag as not yet checked.
\end{caution}

%% =====================================================================================
\section{The decryption proof (paper \S7)}
\label{sec:dec}
%% =====================================================================================
\section{Implementation and evaluation (paper \S8--9)}
\label{sec:eval}
%% =====================================================================================

\paragraph{The code.} The verifiers and provers of both proofs are written in Rocq and extracted to
OCaml. Their algebra is proven correct (Part~\ref{part:formal} explains how); the parsers, the
Fiat--Shamir transcript replay, the Merkle checks and the OCaml drivers are tested only. SHA-256 is
HACL*'s verified implementation. A $\Zq$-element takes 10 bytes on the wire ($q<2^{79}$).

\paragraph{Measurements} (extracted code, one core of an Apple M5~Max, $N=4096$, $q\approx2^{78}$,
$\lambda=128$, proximity regime $\Delta=3$, rate $1/4$; memory is the peak resident size).
\begin{center}\small
\begin{tabular}{@{}rrrrrr@{}}
\toprule
ballots & prove (s/ballot) & verify (s) & proof (MB) & prover (GB) & verifier (GB)\\\midrule
16 & 0.43 & 1.0 & 0.90 & 0.46 & 0.20\\
256 & 0.30 & 5.6 & 2.19 & 5.9 & 2.1\\
1024 & 0.29 & 17.6 & 3.78 & 19.8 & 8.3\\
\bottomrule
\end{tabular}
\end{center}

\paragraph{Reading the numbers.}
\begin{itemize}[leftmargin=1.3em]
\item \emph{Square-root growth:} 64 times more ballots ($16\to1024$) make the proof only $4.2$ times larger
  ($0.90\to3.78$\,MB), as predicted by $O(\sqrt{\nq\tau N})$.
\item \emph{Time per ballot falls} as the number of ballots grows, since part of the work is per proof.
\item \emph{Why the code is slow:} it computes with lists of arbitrary-precision integers, because it is
  written to be proven, not to be fast. Three changes made it usable: the adjoint by three NTTs instead
  of $N^2$ products (proven equal), tail-recursive list functions at extraction, and NTT encoding of the
  slot and proximity vectors (proven equal).
\item \emph{Decryption proof:} 16 ballots with 28 base-4 digits take $62.5$\,s to prove and $2.5$\,s to
  verify; the proof has $1.93$\,MB.
\end{itemize}

\paragraph{Tests.} On small parameters in both proximity regimes, the extracted verifiers accept the
extracted provers' proofs and reject bit-flipped proofs, truncated proofs, and (for the shuffle) the
statement in which two outputs exchange one CRT slot.

%% =====================================================================================
\section{Related work, limitations and open items (paper \S10--11)}
\label{sec:limits}
%% =====================================================================================

\paragraph{Related work.} Lattice shuffles (Costa--Mart\'inez--Morillo, Strand, Aranha et al., Farzaliyev
et al., Herranz--Mart\'inez--S\'anchez with an arithmetic circuit over $\Zq$ and no implementation), the attack
and repair of Bootle et al.\ (tags with invertible differences, still lattice-based), Abdolmaleki et al.,
Orthus. Verified mix-nets: Haines, Gor\'e and Sharma verified a discrete-logarithm shuffle verifier in Coq.
Tools: Ligero/Brakedown commitment, proximity results of Ligero, Diamond--Posen and BCIKS, masked sumcheck of
Chiesa--Forbes--Spooner and Libra, logarithmic derivative of Hab\"ock. The claimed novelty is the combination
``random projection of ring elements + tagged logarithmic derivative'', separating the linear part from the
permutation and shortness.

\paragraph{Limitations stated by the paper.}
\begin{itemize}[leftmargin=1.3em]
\item Machine-checked: both soundness theorems, both zero-knowledge theorems, the appendix, the completeness
  of both provers and their refinement of the zero-knowledge model. Not machine-checked: the line lemma and
  the proximity gap (cited, entered as hypotheses), the round-by-round proposition (on paper; its transitions
  are mechanised), the BCS transformation with its Merkle trees, and the step from counts to adversary
  probabilities.
\item Scope: one mix server and one decryption server; no mix-net-level theorem (noise growth over several
  servers, privacy); no analysis of distributed key generation; $\tau$ must be a power of two (padding with
  trivial ciphertexts proposed).
\item Implementation: slow (lists of big integers); no domain separation in hashing; not constant-time;
  unverified drivers. Not for production.
\item Estimated (not built) size reductions of 14--20\% and 35--41\%.
\end{itemize}

%% =====================================================================================

\paragraph{Related work.} Lattice shuffles (Costa--Mart\'inez--Morillo, Strand, Aranha et al., Farzaliyev
et al., Herranz--Mart\'inez--S\'anchez with an arithmetic circuit over $\Zq$ and no implementation), the attack
and repair of Bootle et al.\ (tags with invertible differences, still lattice-based), Abdolmaleki et al.,
Orthus. Tools: Ligero/Brakedown commitment, proximity results of Ligero, Diamond--Posen and BCIKS, masked
sumcheck of Chiesa--Forbes--Spooner and Libra, logarithmic derivative of Hab\"ock. The claimed novelty is the
combination ``random projection of ring elements + tagged logarithmic derivative'', separating the linear
part from the permutation and shortness.

\paragraph{Limitations stated by the paper.}
\begin{itemize}[leftmargin=1.3em]
\item Proofs are pen-and-paper; the round-by-round proposition is a sketch; the lines-to-affine extension is
  the authors' own derivation.
\item Scope: one mix server and one decryption server; no mix-net-level theorem (noise growth over several
  servers, privacy); no analysis of distributed key generation; $\tau$ must be a power of two (padding with
  trivial ciphertexts proposed).
\item Implementation: non-cryptographic seeded randomness in the native prover, no domain separation in
  hashing, not constant-time; macOS build untested. Not for production.
\item Estimated (not built) size reductions of 14--20\% and 35--41\%.
\end{itemize}

%% =====================================================================================
\section{The appendix: from lines to affine spaces}
\label{sec:appendix}
%% =====================================================================================
\section{A simpler instantiation: the projection with Ligero's constraint tests}
\label{sec:simple}
%% =====================================================================================

The paper's proof has four ingredients: the field projection, the tagged logarithmic derivative,
three masked sumchecks, and a Ligero-style opening. After the projection, the statement is a set of
equations over a field, and any proof system for such equations can be used (paper, Section~3,
``A general template''). This section describes the variant that is \emph{simplest to explain}: it
keeps the projection and the logarithmic derivative, drops the sumchecks entirely, and proves every
equation with the two constraint tests of the original Ligero proof system (Ames, Hazay, Ishai and
Venkitasubramaniam, CCS 2017). The only tool the reader needs, besides the projection, is the one
already used for the commitment: a random combination of committed rows, checked at a few opened
columns.

\begin{caution}[Status]
This variant is \emph{not} what the paper implements, proves or formalises. It is a teaching
device. Its proof size grows as $O(\sqrt{\cdot})$, like the paper's, but it commits more rows, so it
is somewhat larger; no size has been computed for it. Details such as the exact masking polynomials
differ between presentations of Ligero; we give one consistent version.
\end{caution}

\subsection{What has to be proven after the projection}

Recall Sections~\ref{sec:projection} and~\ref{sec:logder}. The prover commits to the witness $W$
(the re-randomisers $w_0,\dots,w_{\tau-1}$, $3\tau N$ entries in $\{-1,0,1\}$) and to the table
$\sigma$. Then the verifier draws $\rho\in\F^{2N}$ and $\alpha,\beta\in\F$, and both sides compute
$\rho'=\rho\,\Mpk$, $z^{\mathsf{pub}}_i=\ip{\rho}{\coeff(c_i)}$ and $\zh_j=\ip{\rho}{\coeff(\ch_j)}$.
The prover then has to convince the verifier of five facts:
\begin{align*}
  &\text{(L1)}\quad h_i = z^{\mathsf{pub}}_i+\ip{\rho'}{\coeff(w_i)}+\beta i-\alpha &&\text{for all } i<\tau,\\
  &\text{(Q1)}\quad u_i\,h_i=1 &&\text{for all } i,\\
  &\text{(Q2)}\quad \uh_j\,\hh_j=1,\ \text{ where }\hh_j=\zh_j+\beta\sigma_j-\alpha &&\text{for all } j,\\
  &\text{(L2)}\quad \textstyle\sum_iu_i=\sum_j\uh_j,\\
  &\text{(Q3)}\quad w^3=w &&\text{for every entry } w \text{ of } W.
\end{align*}
(L1) and (L2) are \emph{linear} in the committed values; (Q1) and (Q2) are \emph{quadratic}; (Q3)
is cubic, and we make it quadratic with one auxiliary table: commit to $W_2=W\odot W$ (entry-wise
square) and prove the two quadratic facts
\[
  \text{(Q3a)}\quad W_2=W\odot W,\qquad \text{(Q3b)}\quad W_2\odot W=W .
\]
Together they say $w^3=w$, which in a field means $w\in\{-1,0,1\}$ (Section~\ref{sec:fm-base},
\texttt{cube\_ternary}). Then, exactly as in the paper, (L1), (Q1), (Q2), (L2) give
$\sum_i1/h_i=\sum_j1/\hh_j$ with no zero denominator, and the logarithmic-derivative lemma turns this
into ``$\sigma$ is a permutation and $\ch_j=c_{\sigma(j)}+\EncZ(w_{\sigma(j)})$'', except with
probability $(2\tau-1)/|\F|$.

\begin{why}[Why is $\hh$ not committed?]
$\hh_j=\zh_j+\beta\sigma_j-\alpha$ is an \emph{entry-wise} affine function of the committed table
$\sigma$ with public coefficients. Because the encoding below is linear, the verifier can compute
the encoding of $\hh$ at any column from the encoding of $\sigma$ at that column and the public
values. So (Q2) can use $\hh$ as if it were committed. This does not work for $h$: in (L1), $h_i$
is an inner product over the $3N$ coefficients of ballot $i$, not an entry-wise function of one
committed entry. So $h$ is committed, and (L1) is proven as a linear constraint.
\end{why}

All of this happens over $\F=\mathrm{GF}(q^2)$: the values $h,u,\uh$ depend on the challenges and live
in $\F$, while $W$, $W_2$ and $\sigma$ have entries in $\Zq\subset\F$. Committing an element of $\F$
costs two $\Zq$ numbers.

\subsection{Committing rows by their values}

The commitment is the one of Section~\ref{sec:commit}, with one change that the quadratic test
needs. Fix distinct points of $\F$:
\begin{itemize}[leftmargin=1.3em]
\item $\ell$ \emph{data points} $\zeta_1,\dots,\zeta_\ell$;
\item $t$ \emph{hiding points} $\zeta_{\ell+1},\dots,\zeta_{\ell+t}$;
\item $n$ \emph{code points} $\eta_1,\dots,\eta_n$, disjoint from the $\zeta$'s.
\end{itemize}
A row of data $x=(x_1,\dots,x_\ell)$ is committed as follows. The prover picks $t$ random values
$x_{\ell+1},\dots,x_{\ell+t}$ and takes the unique polynomial $P_x$ of degree $<k=\ell+t$ with
$P_x(\zeta_c)=x_c$ for all $c\le\ell+t$. The codeword is $(P_x(\eta_1),\dots,P_x(\eta_n))$. As in
the paper, many rows form a matrix, each column is hashed with a salt, and the Merkle root is the
commitment.

The data are now the \emph{values} of the row polynomial at the data points, not its coefficients
as in the paper. The reason is that values multiply entry by entry:
\[
  (P_x\cdot P_y)(\zeta_c)=x_c\,y_c .
\]
So ``$x\odot y=z$'' becomes ``$P_x P_y-P_z$ vanishes at every data point'', a statement about one
polynomial. With coefficients this would fail, since the product of two polynomials has coefficients
that are convolutions, not entry-wise products.

\begin{why}[Why are the code points disjoint from the data points?]
Opening column $j$ reveals $P_x(\eta_j)$ for every row. If $\eta_j$ were a data point, this would be
a data entry. With disjoint points and $t$ random hiding values, any $t$ opened columns are uniform
and independent of the data, by the same argument as Lemma~4.4 of the paper
(Section~\ref{sec:hiding}): the map from the $t$ hiding values to the $t$ opened values is a
bijection.
\end{why}

\subsection{The three tests}

All three tests have the same shape. The verifier sends random coefficients; the prover sends one
polynomial, which is a random combination of row polynomials (plus a random mask committed in
advance); the verifier checks a property of that polynomial and, at the opened columns, that its
values agree with the same combination of the opened entries.

\paragraph{The proximity test.} As in Section~\ref{sec:opening}: for random $\gamma$, the prover sends
$p=m_{\mathsf{prox}}+\sum_i\gamma_iP_i$ (as a polynomial of degree $<k$), and the verifier checks at
the opened columns that $p(\eta_j)=m_{\mathsf{prox}}(\eta_j)+\sum_i\gamma_i\,\mathrm{col}_j[i]$. By the
proximity lemma, if this passes, every committed row is close to a codeword, so it \emph{decodes} to
a unique polynomial $P_i$ of degree $<k$. Soundness is about these decoded polynomials.

\paragraph{The linear test.} The statement is a system $A\,x=b$, where $x$ lists all committed data
entries and $A$, $b$ are public. Write row $i$'s entries as $x_{i,c}=P_i(\zeta_c)$.
\begin{enumerate}[leftmargin=1.5em,itemsep=1pt]
\item The verifier sends a random vector $r$, one coordinate per equation. The combined equation is
  $r^{\top}Ax=r^{\top}b$, that is, $\sum_i\sum_{c\le\ell}a_{i,c}\,P_i(\zeta_c)=r^{\top}b$, where
  $a_{i,c}$ are the entries of $r^{\top}A$.
\item For each row $i$, let $R_i$ be the polynomial of degree $<k$ with $R_i(\zeta_c)=a_{i,c}$ for
  $c\le\ell$ and $R_i(\zeta_c)=0$ at the hiding points. The prover sends
  \[
    q=m_{\mathsf{lin}}+\sum_iR_i\,P_i \qquad(\deg q<2k-1),
  \]
  where $m_{\mathsf{lin}}$ is a random polynomial committed in advance with
  $\sum_{c\le\ell+t}m_{\mathsf{lin}}(\zeta_c)=0$.
\item The verifier checks $\sum_{c\le\ell+t}q(\zeta_c)=r^{\top}b$, and at every opened column
  $q(\eta_j)=m_{\mathsf{lin}}(\eta_j)+\sum_iR_i(\eta_j)\,\mathrm{col}_j[i]$.
\end{enumerate}
\emph{Why it is sound.} If the decoded data violate $Ax=b$, then for a random $r$ the combined
equation fails, except with probability $1/|\F|$. The honest combination $Q=m_{\mathsf{lin}}+\sum_iR_iP_i$
then has the wrong sum, so the prover must send some $q\ne Q$. Both have degree $<2k-1$, so they
agree on at most $2k-2$ points, and the column check catches the difference unless every opened
column falls among these few points (or among the at most $e$ columns where the committed rows differ
from their decodings).

\paragraph{The quadratic test.} The statement is $x^{(s)}\odot y^{(s)}=z^{(s)}$ for several triples
of rows $s$ (one triple per committed row of the products).
\begin{enumerate}[leftmargin=1.5em,itemsep=1pt]
\item The verifier sends random coefficients $r_s$.
\item The prover sends
  \[
    p_0=m_{\mathsf{quad}}+\sum_sr_s\bigl(P_{x^{(s)}}P_{y^{(s)}}-P_{z^{(s)}}\bigr)\qquad(\deg p_0\le 2k-2),
  \]
  where $m_{\mathsf{quad}}$ is a random polynomial committed in advance that vanishes at all data
  points.
\item The verifier checks $p_0(\zeta_c)=0$ for every data point $c\le\ell$, and at every opened column
  $p_0(\eta_j)=m_{\mathsf{quad}}(\eta_j)+\sum_sr_s\bigl(\mathrm{col}_j[x^{(s)}]\cdot\mathrm{col}_j[y^{(s)}]-\mathrm{col}_j[z^{(s)}]\bigr)$.
\end{enumerate}
\emph{Why it is sound.} If some product is wrong at some data point $\zeta_c$, then for random $r$ the
honest combination $P_0$ does not vanish at $\zeta_c$, except with probability $1/|\F|$. The prover
must send a $p_0$ that vanishes there, so $p_0\ne P_0$; both have degree $\le2k-2$, so they agree on
at most $2k-2$ code points, and the column check catches them. This needs $n$ comfortably larger than
$2k$; for instance $n=4k$ (rate about $1/4$) makes each opened column catch a cheater with probability
about $1/2$.

\emph{Why the masks.} Without $m_{\mathsf{lin}}$ and $m_{\mathsf{quad}}$, the polynomials $q$ and $p_0$
would reveal combinations of the hidden values. A mask is a uniformly random polynomial with exactly
the property the verifier checks (zero sum, or zero at the data points), so the sent polynomial is
uniform among those that pass, and reveals nothing beyond the fact being proven. This plays the role
of the slot masks $m_v$ of the paper's opening.

\begin{example}[The quadratic test over $\Z_{17}$]
Take one row pair for (Q3b), $W_2\odot W=W$, with $\ell=2$ data entries, $t=1$ hiding point and
$n=8$ code points: data points $\zeta_1=1$, $\zeta_2=2$, hiding point $\zeta_3=3$, code points
$\eta=4,5,\dots,11$.

\emph{Honest prover.} $w=(1,16)$ (that is, $(1,-1)$) and $W_2=w\odot w=(1,1)$. With hiding values $5$
and $9$:
\[
  P_w=11+3X+4X^2,\qquad P_{w_2}=9+5X+4X^2 .
\]
(Check: $P_w(1)=18\equiv1$, $P_w(2)=33\equiv16$, $P_w(3)=56\equiv5$.) The codewords at $\eta=4..11$ are
\[
  P_w:\ (2,7,3,7,2,5,16,1),\qquad P_{w_2}:\ (8,15,13,2,16,4,0,4).
\]
With $r=1$ and no mask, $P_0=P_{w_2}P_w-P_w=3+11X+6X^2+15X^3+16X^4$. It vanishes at the data points:
$P_0(1)=P_0(2)=0$. At the code points it takes the values $(14,13,2,7,13,15,1,3)$, which is exactly
$\mathrm{col}_j[w_2]\cdot\mathrm{col}_j[w]-\mathrm{col}_j[w]$ column by column; for $\eta=4$:
$8\cdot2-2=14$. Every check passes.

\emph{Cheating prover.} Now $w=(2,16)$, which is not ternary, and $W_2=(4,1)$. The true
$P_0=16X+12X^2+12X^4$ has $P_0(1)=6\ne0$. To pass the first check, the prover must send some
polynomial of degree $\le4$ that vanishes at $1$ and $2$; it differs from $P_0$, and two different
polynomials of degree $\le4$ agree on at most $4$ of the $8$ code points. So each opened column
exposes the cheat with probability at least $1/2$, and $t$ columns let it through with probability
at most $2^{-t}$.
\end{example}

\subsection{The whole protocol}

\begin{enumerate}[leftmargin=1.6em,itemsep=2pt]
\item[\textbf{1.}] The prover commits to the rows of $W$, $W_2=W\odot W$ and $\sigma$, and to the mask
  rows $m_{\mathsf{prox}}$, $m_{\mathsf{lin}}$, $m_{\mathsf{quad}}$. It sends the root.
\item[\textbf{2.}] The verifier sends $\rho$, $\alpha$, $\beta$. Both sides compute $\rho'$,
  $z^{\mathsf{pub}}$ and $\zh$.
\item[\textbf{3.}] The prover computes $h$, $u=1/h$ and $\uh=1/\hh$, commits to their rows (a second
  batch, over $\F$), and sends the second root. In (Q1) the right-hand side is the public all-ones
  row, whose polynomial is the constant $1$. If some $h_i$ or $\hh_j$ is zero it aborts (this
  happens with probability $2\tau/|\F|$).
\item[\textbf{4.}] The verifier sends the coefficients of the three tests: $\gamma$ for proximity,
  $r$ for the linear test (covering (L1) and (L2)), and $r_s$ for the quadratic test (covering (Q1),
  (Q2), (Q3a), (Q3b)).
\item[\textbf{5.}] The prover sends the three polynomials $p$, $q$, $p_0$.
\item[\textbf{6.}] The verifier checks the sum condition of $q$ and the vanishing condition of $p_0$,
  then sends $\nq$ column indices. The prover opens these columns of both batches, with salts and
  Merkle paths.
\item[\textbf{7.}] The verifier checks the Merkle paths and the three column equations. For (Q2) it
  computes the encoding of $\hh$ at column $j$ itself, as
  $\mathrm{enc}(\zh)_j+\beta\,\mathrm{col}_j[\sigma]-\alpha$.
\end{enumerate}

\paragraph{What the verifier computes on its own.} The linear test needs $r^{\top}A$. For (L1), the
equation of ballot $i$ involves the $3N$ entries of $w_i$ with coefficients $\rho'$, so the
coefficient of entry $(i,y,k)$ in $r^{\top}A$ is $r_i\,\rho'_k(y)$. The verifier computes $\rho'$ once
(two adjoint multiplications, as in the paper) and then $3\tau N$ products. Its work stays linear
in the size of the statement.

\subsection{Why it is sound and zero-knowledge}

\paragraph{Soundness, step by step.} Read the checks backwards, as in Section~\ref{sec:sound}.
The proximity test makes every committed row decodable. The linear test makes (L1) and (L2) hold on
the decoded data, except with probability about $1/|\F|$ plus the column-query error. The quadratic
test does the same for (Q1), (Q2), (Q3a), (Q3b). From (Q3a) and (Q3b), every entry of $W$ is
ternary. From (L1), (Q1), (Q2) and (L2), the sums of inverses agree with no zero denominator, and
the logarithmic-derivative lemma (with the projection inside it) gives the shuffle relation, except
with probability $(2\tau-1)/|\F|$. The challenges $\rho,\alpha,\beta$ are drawn after $W$ and $\sigma$
are committed, which is what that lemma needs. The total error is a sum of a few terms of the form
$O(\tau)/|\F|$ and the query term, which is $2^{-128}$-small for the paper's parameters.

\paragraph{Zero knowledge.} Every value the verifier sees is uniform subject to the checks: the
opened columns, by the hiding values of every row; the three polynomials, by their masks. No value
that depends on $\sigma$ is ever sent in the clear. (Ligero's paper states the precise conditions:
the number of hiding values per row must be at least the number of opened columns, and the masks
must have the degrees given above.) In particular, the sums $\sum u$ and $\sum\uh$
are never revealed; only their equality is proven, inside the linear test. (Revealing a product
$\prod h_i$, as the earlier grand-product design did, would leak $\sigma$.)

\subsection{Comparison with the paper's construction}

\begin{center}\small
\begin{tabular}{@{}p{0.27\textwidth}p{0.33\textwidth}p{0.33\textwidth}@{}}
\toprule
 & this section & the paper\\
\midrule
reduces the equations by & Ligero's linear and quadratic tests & three masked sumchecks\\
committed tables & $W$, $W_2$, $\sigma$, $h$, $u$, $\uh$ & $W$, $\sigma$, $u$, $\uh$, mask table\\
(L1), the linear map & one linear test, $h$ committed & sumcheck B, $h$ never committed\\
ternary check & two quadratic tests via $W_2$ & sumcheck C on $w^3-w$\\
row encoding & values at data points & coefficients\\
proof size & $O(\sqrt{\cdot})$, somewhat larger (more rows) & $O(\sqrt{\nq\,\tau N})$\\
ideas to explain & projection, tags, RS rows, random combinations & the same, plus multilinear extensions, sumcheck, masks\\
formalised & no & yes\\
\bottomrule
\end{tabular}
\end{center}

The paper's sumchecks are worth their extra explanation for one reason above all: sumcheck B proves
the linear map (L1) without committing $h$, and sumcheck C proves ternarity without committing $W_2$,
a table as large as $W$ itself. Since $W$ dominates the proof, avoiding a second table of that size
matters. For explaining \emph{why} the construction works, however, the version in this section
contains everything essential: once the ciphertexts are projected into a field, a shuffle is nothing
more than a few linear and quadratic equations on committed rows.

%% =====================================================================================
\part{The formalisation, file by file}
\label{part:formal}
%% =====================================================================================

%% =====================================================================================
\section{How to read the formal development}
\label{sec:formal-intro}
%% =====================================================================================

Parts~I and~II explained the mathematics of the paper. This part explains how that mathematics is
written down and checked by a computer. We go through the Rocq files of the repository
(directory \texttt{src/}) from the bottom up: each file uses only files described before it. For
every file we say what mathematical object it defines, state its main theorems in ordinary
mathematical notation, sketch why they are true, and say which later files use them. The last
section of this part puts the pieces together and shows how they give the two protocol theorems,
the executable verifiers and provers, and the zero-knowledge results.

\subsection{What Rocq is, in three paragraphs}

Rocq (formerly called Coq) is a \emph{proof assistant}: a program in which one writes
definitions, statements and proofs, and which checks every proof step by step. A definition is a
function or a predicate written in a small programming language. A statement is a formula. A proof
is a script of \emph{tactics} (``rewrite with this equation'', ``do induction on this list'',
``this follows from that lemma'') that Rocq turns into a proof term and checks with a small trusted
\emph{kernel}. If the kernel accepts a proof, the statement follows from the definitions and from
the axioms in use, and from nothing else.

Two features matter for reading the files. First, a file often opens a \texttt{Section} with
\texttt{Variable} and \texttt{Hypothesis} declarations. Inside the section these are fixed objects
and assumptions; when the section is closed, every theorem becomes universally quantified over the
variables it uses and gets the hypotheses it uses as extra premises. So ``\texttt{Variable F},
\texttt{Hypothesis Hf : field F}, \texttt{Theorem T : P}'' means ``for every field $F$, $P$
holds''. Second, Rocq can \emph{extract} a definition to OCaml: the function becomes an OCaml
program that runs at normal speed. The proofs about it carry over, provided that the extraction
mechanism and the OCaml compiler are correct; these are part of the trusted base.

The development uses \emph{no axioms}. The script \texttt{audit.sh} asks Rocq, for each of $335$
main theorems, which axioms the proof depends on (\texttt{Print Assumptions}). The answer is
``Closed under the global context'' for every one of them: they follow from the definitions alone.
What a theorem assumes beyond that is visible in its statement, as an explicit hypothesis. The most
important such hypothesis is the proximity lemma for Reed--Solomon codes, which is cited from the
literature and not proven (Section~\ref{sec:fm-codes}).

\subsection{How probability is expressed}

The protocol theorems bound probabilities, such as ``a cheating prover is accepted with probability
at most $\varepsilon$''. The development expresses every such bound as a \emph{count}. A uniform
challenge is drawn from a finite list without duplicates: the list \texttt{lf} of all field
elements, or the list $\texttt{all\_lists lf } m$ of all $|\F|^m$ vectors of length $m$. An event is
a predicate on challenges, and ``$\Pr[E]\le B/|\Omega|$'' becomes ``at most $B$ elements of the
list $\Omega$ satisfy $E$''. Bounds are kept multiplied out in natural numbers. For example, the
sumcheck bound $\Pr\le\sum_j d_j/|\F|$ over challenges in $\F^n$ is stated as
\[
  \#\{r\in\F^n : \text{the verifier is fooled at } r\}\;\le\;\Bigl(\sum_j d_j\Bigr)\,|\F|^{n-1}.
\]
This avoids real numbers and measure theory entirely. The cost is that ``probability'' never
appears as a word in the files; it appears as a length of a filtered list or as the predicate
\texttt{card\_le} of \texttt{Card.v} (Section~\ref{sec:fm-counting}).

\subsection{The layers}

The files form a hierarchy. Listed from the bottom:
\begin{enumerate}[leftmargin=1.6em,itemsep=1pt]
\item \textbf{General libraries} (\texttt{src/Algebra}, \texttt{src/Utility},
  \texttt{src/Probability}): the algebraic hierarchy up to fields, modular arithmetic and Fermat's
  little theorem, SHA-256 (Section~\ref{sec:fm-lib}).
\item \textbf{Counting and polynomials over any field} (\texttt{src/Crypto}, first files):
  counting, root bounds, multilinear extensions, the negacyclic ring, BGV
  (Section~\ref{sec:fm-base}).
\item \textbf{Sumcheck} (Section~\ref{sec:fm-sumcheck}).
\item \textbf{Codes and the opening} (Section~\ref{sec:fm-codes}).
\item \textbf{Soundness of the shuffle proof} (Section~\ref{sec:fm-mix}).
\item \textbf{Soundness of the decryption proof} (Section~\ref{sec:fm-dec}).
\item \textbf{Zero knowledge over any field} (Section~\ref{sec:fm-zk}).
\item \textbf{The concrete field $\mathrm{GF}(q^2)$} (Section~\ref{sec:fm-concrete}).
\item \textbf{Executable verifiers} (\texttt{src/Examples/Exec*}, Section~\ref{sec:fm-exec}).
\item \textbf{Executable provers and their completeness} (Section~\ref{sec:fm-prove}).
\item \textbf{Zero knowledge at $\mathrm{GF}(q^2)$ and for the executable provers}
  (Section~\ref{sec:fm-zkexec}).
\end{enumerate}
Section~\ref{sec:fm-together} shows how the layers combine, and Section~\ref{sec:fm-index} is an
index of all files.

Layers 2 to 7 are \emph{generic}: they work over an arbitrary field $F$ with decidable equality
(given by a field structure \texttt{Hf : field ...} and a function \texttt{Fdec} that decides
$x=y$). Nothing in them uses that $F=\mathrm{GF}(q^2)$. The concrete field enters in layer 8, and
the executable code in layers 9 to 11.

\begin{caution}[Names]
Rocq names are written in \texttt{typewriter} font with the file in which they live, for instance
\texttt{LogDer.v: logder\_count}. Mathematical statements are paraphrases of the Rocq statements
that keep every hypothesis; where the Rocq statement is more general than needed, we say so.
\end{caution}

%% =====================================================================================
\section{Layer 1: general libraries}
\label{sec:fm-lib}
%% =====================================================================================

These directories come from an earlier, more general library for sigma protocols. Only a few of
their results are used by the mix-net development; we describe all files briefly and say which
ones matter.

\subsection{\texttt{src/Algebra}: from monoids to fields}

\paragraph{\texttt{Hierarchy.v}.} Defines the algebraic structures as Rocq \emph{classes}
(collections of operations and laws that Rocq finds automatically). Each structure carries its own
equality and the requirement that the operations respect it.
\begin{itemize}[leftmargin=1.3em]
\item A \texttt{monoid} is a set with an associative operation and a two-sided identity.
\item A \texttt{group} adds inverses: $x^{-1}x=xx^{-1}=e$. A \texttt{commutative\_group} adds $xy=yx$.
\item A \texttt{ring} $(T,0,1,-,+,\cdot)$: $(T,+,0,-)$ is a commutative group, $(T,\cdot,1)$ is a
  monoid, multiplication distributes over addition on both sides, and $x-y=x+(-y)$. A
  \texttt{commutative\_ring} has $xy=yx$.
\item An \texttt{integral\_domain} is a commutative ring with $0\ne1$ and no zero divisors
  ($xy=0\Rightarrow x=0\lor y=0$).
\item A \texttt{field} is a commutative ring with $0\neq1$, $x^{-1}x=1$ for $x\ne0$, and
  $x/y=xy^{-1}$. The value of $0^{-1}$ is left unspecified.
\item A \texttt{vector\_space} over a field (written as a right action $v\cdot r$) satisfies the
  usual axioms.
\end{itemize}
Every protocol file is stated ``for every \texttt{field}''; this file is where that word is defined.

\paragraph{\texttt{Monoid.v}, \texttt{Group.v}, \texttt{Ring.v}, \texttt{Integral\_domain.v}.} The
standard identities: cancellation ($zx=zy\iff x=y$ in a group), $(xy)^{-1}=y^{-1}x^{-1}$,
uniqueness of inverses, $0x=0$, $x(-y)=-(xy)$, $x-y=0\iff x=y$. \texttt{Ring.v} also connects an
abstract commutative ring to Rocq's \texttt{ring} tactic, which proves polynomial identities
automatically, and \texttt{Integral\_domain.v} does the same for Rocq's integral-domain tactics.

\paragraph{\texttt{Field.v}.} The important export is \texttt{field\_theory\_for\_stdlib\_tactic}:
every abstract field is registered with Rocq's \texttt{field} tactic, which proves identities of
rational expressions given that the denominators are non-zero. Almost every protocol file starts
with \texttt{Add Field}, using this theorem; this is why long algebraic manipulations in the proofs
often take one line. The file also proves $xx^{-1}=1$ for $x\ne0$, uniqueness of inverses, and that
a product of non-zero elements is non-zero, so every field with decidable equality is an integral
domain.

\paragraph{\texttt{Vector\_space.v}, \texttt{Product\_space.v}, \texttt{Vector\_product\_space.v}.}
Vector spaces read multiplicatively (a group with exponents in a field, as in discrete-logarithm
protocols), and products of vector spaces. They are not used by the mix-net development.

\subsection{\texttt{src/Utility}: arithmetic modulo a prime, and SHA-256}

\paragraph{\texttt{Util.v}.} Infrastructure for fixed-length vectors and finite index types
(conversions, take and drop, pairs of distinct positions) and a few modular-arithmetic facts. It is
used by the other Utility files.

\paragraph{\texttt{Functions.v}.} Fast modular exponentiation by square-and-multiply, proven equal to
the slow repeated multiplication (\texttt{npow\_mod\_exp\_unary\_binary\_eqv}) and to Rocq's
\texttt{Zpow\_mod} (\texttt{zmod\_nmod}). It is used to compute inverses as $u^{p-2}$.

\paragraph{\texttt{Fermat.v}.} A self-contained proof of Fermat's little theorem: the binomial
theorem, $p\mid\binom pk$ for $0<k<p$, the ``freshman's dream'' $(x+1)^p\equiv x^p+1\pmod p$, and
finally \texttt{fermat\_little\_ZZ}: for a prime $p$ and $0<a<p$, $a^{p-1}\equiv1\pmod p$. This
justifies $a^{-1}=a^{p-2}$, and \texttt{ConcreteField.v} uses it to show that $3$ is not a square
modulo $q$.

\paragraph{\texttt{Zpstar.v}.} Concrete instances. \texttt{Zpfield} is the field
$\mathbb{Z}_p$ for a prime $p$: elements are integers $0\le v<p$, operations reduce modulo $p$, the
inverse is $v^{p-2}$, and \texttt{zp\_field} proves the field axioms (the inverse law from Fermat's
theorem). This is the field $\Zq$ of the implementation. The file also defines the multiplicative
group $\mathbb{Z}_p^*$ and a Schnorr subgroup, which the mix-net does not use.

\paragraph{\texttt{SchnorrOrder.v}, \texttt{Insert.v}.} Order of elements in a Schnorr group, and
inserting an element into a vector. Not used here.

\paragraph{\texttt{Sha256.v}.} A Rocq definition of SHA-256 following FIPS 180-4, with proofs about
the padding and four NIST test vectors checked by computation. Nothing deeper is proven about
SHA-256: in the security analysis the hash is a random oracle, so its internals do not matter.

\subsection{\texttt{src/Probability}}

\texttt{Prob.v} (non-negative rationals) and \texttt{Distr.v} (finite distributions as weighted
lists, and the fact that a bijection of the coins does not change a uniform distribution) come from
the same earlier library. \emph{No file of the mix-net development imports them}: as explained in
Section~\ref{sec:formal-intro}, probabilities here are counts, and the zero-knowledge statements
are permutations of lists.

%% =====================================================================================
\section{Layer 2: counting and polynomials over any field}
\label{sec:fm-base}
%% =====================================================================================

From here on, unless said otherwise, $F$ is an arbitrary field with decidable equality, and
$\mathtt{lf}$ is a list of elements of $F$ without duplicates; when a theorem needs ``uniform over
all of $F$'', it assumes that $\mathtt{lf}$ contains every element of $F$, and then $|\mathtt{lf}|=|F|$.

\subsection{\texttt{Counting.v}: probability as counting}
\label{sec:fm-counting}

This file contains no field at all; it is pure combinatorics on finite lists, and every later
probability bound rests on it.

\paragraph{Definitions.} $\mathtt{all\_lists}\ \mathtt{lf}\ m$ is the list of all lists of length
$m$ with entries in $\mathtt{lf}$, that is $\mathtt{lf}^m$. $\mathtt{fall}(p,t)=p(p-1)\cdots(p-t+1)$
is the falling factorial, the number of duplicate-free $t$-tuples from $p$ elements.

\paragraph{Theorems.}
\begin{itemize}[leftmargin=1.3em]
\item \texttt{all\_lists\_length}: $|\mathtt{lf}^m|=|\mathtt{lf}|^m$, and \texttt{all\_lists\_NoDup}:
  the list has no duplicates if $\mathtt{lf}$ has none. So $\mathtt{lf}^m$ is a valid uniform
  sample space.
\item \texttt{count\_swap} (double counting): for a relation $f$ between two lists,
  $\sum_{a}\#\{b: f(a,b)\}=\sum_b\#\{a:f(a,b)\}$.
\item \texttt{shift\_count} (averaging). Let $X$ be a duplicate-free list and $(\varphi_t)_{t}$ a
  family of injective maps $X\to X$. If every $x\in X$ satisfies $\#\{t:\mathit{bad}(\varphi_t x)\}\le d$,
  then $\#\{x\in X:\mathit{bad}(x)\}\cdot(\text{number of } t)\le d\cdot|X|$.

  In words: if every point is bad for few of its ``translates'', then few points are bad. This is
  how a bound along one line becomes a bound on the whole space.
\item \texttt{all\_pass\_count}: the number of $t$-tuples whose entries all pass a test is
  (number of passing elements)$^t$. This is independence of repeated draws.
\item \texttt{distinct\_pass\_count}: the number of duplicate-free $t$-tuples whose entries all pass
  is $\mathtt{fall}(P,t)$, where $P$ is the number of passing elements; and
  \texttt{distinct\_query\_bound} turns this into $\Pr[\text{all } t \text{ distinct queries pass}]
  \le (P/n)^t$. This is the counting behind the column queries of the opening.
\end{itemize}

\paragraph{Proof idea of \texttt{shift\_count}.} Sum the per-point bound over $X$: the total is at
most $d|X|$. Swap the two sums with \texttt{count\_swap}. For each fixed $t$, the map $\varphi_t$
permutes $X$ (an injective map of a finite set to itself is onto), so it hits exactly as many bad
points as there are; the swapped sum is $\#\mathit{bad}\cdot\#t$.

\subsection{\texttt{Card.v}: events that are not decidable, and composing layers}

Some events are statements that a program cannot test (``there is a witness such that \dots'').
For them the file defines
\[
  \mathtt{card\_le}\ \Omega\ E\ B\quad:\Longleftrightarrow\quad
  \text{every duplicate-free list of elements of } \Omega \text{ that all satisfy } E \text{ has length } \le B,
\]
which is ``$\#\{\omega\in\Omega:E(\omega)\}\le B$'' without needing to decide $E$.

\paragraph{The calculus.}
\begin{itemize}[leftmargin=1.3em]
\item \texttt{card\_le\_mono} (a smaller event has a smaller count), \texttt{card\_le\_or} (union
  bound: $B_1+B_2$), \texttt{card\_le\_of\_filter} (from an ordinary count of a decidable event).
\item \texttt{card\_le\_cons}, \texttt{card\_le\_head}: on $\mathtt{lf}^{m+1}$, if for every first
  coordinate the rest of the event has count $B$, the whole has count $B\cdot|\mathtt{lf}|$; an event
  that depends only on the first coordinate and has count $B$ there has count $B|\mathtt{lf}|^m$.
\item \texttt{card\_le\_prod\_cyl} and \texttt{card\_le\_prod\_fiber}: the same on a product
  $\Omega_A\times\Omega_B$ (``cylinder'' events and events bounded fibre by fibre).
\item \texttt{card\_join} (composing two protocol layers). Let $\Omega_m$ be the challenges of a
  first layer and $\Omega_o$ those of a second. Suppose
  (i) the first layer is sound: $\mathtt{card\_le}\ \Omega_m\ \mathit{Ideal}\ B_m$;
  (ii) whenever a ``link'' condition $D(w)$ holds and the second layer accepts, the ideal checks of
  the first layer hold;
  (iii) whenever $D(w)$ fails, the second layer accepts at most $B_o$ of its challenges.
  Then the combined acceptance has count at most $B_m|\Omega_o|+B_o|\Omega_m|$.

  This is exactly the shape of ``sumcheck layer + opening'': $D$ says that the values the prover
  claims are the true values of the committed rows.
\item \texttt{qspace}$(n,t)$ is the list of duplicate-free $t$-tuples of columns $<n$, with
  $|\texttt{qspace}(n,t)|=\mathtt{fall}(n,t)$ (\texttt{qspace\_length}). \texttt{query\_card}: if at
  most $n-e-1$ columns pass a test, then at most $\mathtt{fall}(n-e-1,t)$ query tuples pass on all
  their columns.
\end{itemize}

\subsection{\texttt{PolyRoots.v}: polynomials and their roots}

A polynomial is its list of coefficients, constant term first; $\mathtt{eval}\ p\ x$ evaluates it
by Horner's rule. $\mathtt{pzero}\ p$ means that all coefficients are zero.

\begin{theorem}[\texttt{roots\_bound}]
If $p$ is not the zero list, and $r_1,\dots,r_s$ are distinct roots of $p$, then $s<|p|$, where
$|p|$ is the number of coefficients. In other words, a non-zero polynomial of degree $\le d$ has at
most $d$ roots.
\end{theorem}
\emph{Proof idea.} Synthetic division (\texttt{divide}, \texttt{divide\_spec}): $p(x)=(x-r)q(x)+p(r)$
with $q$ one coefficient shorter. At a root, $q$ is non-zero and has the other roots; induct.

This one fact is behind almost every soundness bound: the Schwartz--Zippel steps of the sumcheck,
the distance of Reed--Solomon codes, the injectivity of evaluation, and the logarithmic-derivative
lemma. The file also has \texttt{pigeonhole\_notin}: if fewer than $|\mathtt{lf}|$ elements are
``bad'', some element of $\mathtt{lf}$ is good.

\subsection{\texttt{LogDer.v}: the logarithmic-derivative lemma}

This file proves Lemma~5.1 of the paper (``tagged logarithmic derivative'', Section~\ref{sec:logder}
of this tutorial), in counting form and without multivariate polynomials. It also defines the
list operations used everywhere: $\mathtt{fsum}$ (sum of a list), $\mathtt{dot}\ a\ p=\sum_i a_ip_i$
(truncated to the shorter list), $\mathtt{padd}$, $\mathtt{pscale}$ (coefficient-wise addition and
scaling), and $\mathtt{nF}\ k=1+\dots+1$ ($k$ times), the image of the natural number $k$ in $F$.

\paragraph{Setting.} A challenge is $p=(\alpha,c)\in F^{n+1}$. Each item of the two multisets is a
vector $a=(-1,a')$ with $|a'|=n$, so that $\ip{a}{p}=\ip{a'}{c}-\alpha$. In the shuffle proof
$a'=\coeff(\text{ciphertext})$ followed by the tag, and $c=\rho$ followed by $\beta$.

\begin{theorem}[\texttt{logder\_count}]
Assume $\mathtt{lf}$ enumerates $F$; $S$ and $T$ are lists of $\tau$ items of the form above;
$1,2,\dots,\tau$ are non-zero in $F$; and $S$ is not a permutation of $T$. Then
\[
  \#\Bigl\{p\in F^{n+1} :\ \text{all } \ip{a}{p}\ne 0 \text{ and }
  \sum_{a\in S}\frac{1}{\ip{a}{p}}=\sum_{a\in T}\frac{1}{\ip{a}{p}}\Bigr\}\;\le\;(2\tau-1)\,|F|^n,
\]
that is, the probability is at most $(2\tau-1)/|F|$.
\end{theorem}

\emph{Proof idea.} Let $D$ be the distinct items and $\mu_a$ the net multiplicity of $a$ (its
number of occurrences in $S$ minus in $T$); some $\mu_s\ne 0$, and it stays non-zero in $F$ because
$1,\dots,\tau\ne0$ (\texttt{nF\_inj}). Clearing denominators, the event implies $\Phi(p)=0$ for
\[
  \Phi(p)=\sum_{a\in D}\mu_a\prod_{b\in D,\,b\ne a}\ip{b}{p}
\]
(\texttt{SS} in the file; \texttt{SS\_rational} is the clearing of denominators). Restricted to a
line $p+t\,\mathit{dir}$, $\Phi$ is a polynomial in $t$ with at most $|D|$ coefficients whose top
coefficient is $\Phi(\mathit{dir})$ (\texttt{eval\_SSp}, \texttt{nth\_SSp\_top}). Choose $\mathit{dir}$
with $\ip{s}{\mathit{dir}}=0$ and $\ip{b}{\mathit{dir}}\ne0$ for all other $b$
(\texttt{forms\_common\_nonzero}); then $\Phi(\mathit{dir})=\mu_s\prod_{b\ne s}\ip{b}{\mathit{dir}}\ne0$
(\texttt{SS\_special}). So on every line in this direction $\Phi$ has at most $|D|-1\le 2\tau-1$ roots
(\texttt{roots\_bound}), and \texttt{shift\_count} averages over all parallel lines.

The file is imported by almost every later file (for $\mathtt{fsum}$, $\mathtt{dot}$, $\mathtt{nF}$);
the theorem itself is used in \texttt{MixMultiset.v}.

\subsection{\texttt{PolyFun.v}: ``this function is a polynomial of degree $\le d$''}

$\mathtt{is\_poly}\ d\ f$ means that some coefficient list $p$ with at most $d+1$ entries satisfies
$\mathtt{eval}\ p\ X=f(X)$ for all $X$. The file proves closure rules: constants and the identity
are polynomials; sums and differences have degree $\max(d_1,d_2)$; products have degree $d_1+d_2$
(\texttt{is\_poly\_mul}, via the product of coefficient lists \texttt{pmul}); finite sums keep the
degree; and if $\varphi$ has degree $\le D$ and $f$ degree $\le d$ then $\varphi\circ f$ has degree
$\le Dd$ (\texttt{is\_poly\_peval}). The sumcheck soundness theorem needs a degree bound for every
round function, and these rules let later files read the degree off the formula of a summand.

\subsection{\texttt{Mle.v}: the Boolean cube and multilinear extensions}

\paragraph{Definitions.} $\mathtt{cube}\ m=\{0,1\}^m$ (as the list $\mathtt{all\_lists}\ [0;1]\ m$),
$\mathtt{eqp}(b,x)=\prod_i\bigl(b_ix_i+(1-b_i)(1-x_i)\bigr)$, and for a table $T$ (any function from
lists to $F$; only its values on the cube matter)
\[
  \mathtt{mle}\ m\ T\ x=\sum_{b\in\{0,1\}^m}T(b)\,\mathtt{eqp}(b,x).
\]
\paragraph{Theorems.} \texttt{mle\_S} splits off the first variable:
$\mathtt{mle}_{m+1}T(x_0,x')=(1-x_0)\,\mathtt{mle}_m T(0,\cdot)(x')+x_0\,\mathtt{mle}_mT(1,\cdot)(x')$.
\texttt{mle\_cube}: on the cube, the extension equals the table. \texttt{mle\_const\_one}:
$\sum_b\mathtt{eqp}(b,x)=1$.
\begin{theorem}[\texttt{mle\_zero\_count}]
If $T(b)\ne0$ for some $b\in\{0,1\}^m$, then $\#\{x\in F^m:\mathtt{mle}\ m\ T\ x=0\}\le m\,|F|^{m-1}$.
\end{theorem}
\emph{Proof idea.} Induction on $m$. One of the halves $T(0,\cdot)$, $T(1,\cdot)$ is non-zero on the
smaller cube. For each tail $x'$, either that half's extension vanishes at $x'$ (counted by
induction), or the expression in $x_0$ is a non-zero affine function, which has at most one root
(\texttt{affine\_roots}).

This theorem justifies all ``batching'' challenges: if $\sum_x\eq(t,x)T(x)=0$ at a random $t$,
then $T$ is zero on the cube, except with probability $m/|F|$.

\subsection{\texttt{CubeIndex.v}: cube points and indices}

The implementation stores a table as an array indexed by integers; the proofs see a function on
the cube. $\mathtt{pt}\ l\ i$ is the point of $\{0,1\}^l$ whose coordinates are the bits of $i$,
least significant first. \texttt{pt\_perm}: $i\mapsto\mathtt{pt}\ l\ i$ for $i<2^l$ lists the cube
exactly once, so $\sum_{b\in\{0,1\}^l}f(b)=\sum_{i<2^l}f(\mathtt{pt}\ l\ i)$ (\texttt{cube\_sum\_seq}).
$\mathtt{idr}(x)=\sum_i2^ix_i$ is the tag function $\mathrm{id}$ of the paper; \texttt{idr\_pt}:
$\mathtt{idr}(\mathtt{pt}\ l\ i)=i$, and \texttt{mle\_idr}: $\mathrm{id}$ is its own multilinear
extension. The file also proves that $\mathtt{mle}$ is linear (\texttt{mle\_add}, \texttt{mle\_scale},
\texttt{mle\_fsum}) and the ``prefix'' rule
$\mathtt{mle}_{n+l}T(y,r)=\mathtt{mle}_l\bigl(T(y,\cdot)\bigr)(r)$ for $y$ on the cube
(\texttt{mle\_prefix\_cube}), which gives the table $\mle W(y,r)$ of sumcheck~B.

\subsection{\texttt{Negacyclic.v}: the ring $F[X]/(X^N+1)$}

Elements are coefficient lists of length $N$. $\mathtt{xmul}\ r$ is multiplication by $X$: shift up
and negate the coefficient that wraps around, because $X^N=-1$. $\mathtt{nmul}\ a\ r=\sum_ia_iX^ir$
is the product $a\cdot r$, and $\mathtt{nadj}\ a\ \rho$ is its transpose.
\begin{itemize}[leftmargin=1.3em]
\item \texttt{nmul\_adjoint}: $\ip{\rho}{a\cdot r}=\ip{\mathtt{nadj}\ a\ \rho}{r}$. This is the
  ``verifier's linear algebra'' of the paper: the verifier moves the public factor $a$ from the
  secret $r$ onto its own challenge $\rho$.
\item \texttt{nmul\_eval}: if $\zeta^N=-1$ then $(a\cdot r)(\zeta)=a(\zeta)\,r(\zeta)$. Evaluation at a
  root of $X^N+1$ is a ring homomorphism; these are the CRT slots.
\item \texttt{poly\_eq\_of\_evals}: two lists of length $N$ that agree at $N$ distinct points are
  equal (by \texttt{roots\_bound} applied to their difference).
\end{itemize}

\subsection{\texttt{NttInverse.v}: the adjoint through evaluations}

This file justifies computing the adjoint with three number-theoretic transforms instead of $N^2$
products.
\begin{itemize}[leftmargin=1.3em]
\item \texttt{eval\_nadj}: if $z^N=-1$ and $zw=1$, then $(\mathtt{nadj}\ a\ \rho)(z)=a(w)\,\rho(z)$, that
  is, the adjoint is $a(X^{-1})\rho(X)$ at every root $z$ of $X^N+1$.
\item \texttt{orth}: if $\omega$ has order exactly $n$ and $\omega\omega'=1$, then
  $\sum_{j<n}(\omega^k\omega'^i)^j$ is $n$ when $k=i$ and $0$ otherwise.
\item \texttt{ntt\_inverse}: under the same hypotheses, for $|r|=n$ and $i<n$,
  $\sum_{j<n}r(\psi\omega^j)\,\omega'^{ij}=n\,\psi^i r_i$. So the coefficients can be recovered from
  the values at the points $\psi\omega^j$.
\end{itemize}
\emph{Proof idea.} Expand $r(\psi\omega^j)=\sum_kr_k\psi^k\omega^{jk}$, swap the sums, and apply
\texttt{orth}; only $k=i$ survives. The executable version is in \texttt{ExecNadj.v}
(Section~\ref{sec:fm-exec}).

\subsection{\texttt{Bgv.v}: ciphertexts, projection and the shuffle relation}

With a public key $(a,b)$ and plaintext modulus $p$ (all over $F$):
$\mathtt{enc0}(r,e_1,e_2)=(a\cdot r+pe_1,\ b\cdot r+pe_2)$; $\mathtt{proj}\ \rho_u\ \rho_v\ (u,v)=
\ip{\rho_u}{u}+\ip{\rho_v}{v}$ is the projection of the paper; and the three blocks of
$\rho'=\rho\,\Mpk$ are $\rho'_1=\mathtt{nadj}\ a\ \rho_u+\mathtt{nadj}\ b\ \rho_v$, $\rho'_2=p\rho_u$,
$\rho'_3=p\rho_v$.
\begin{itemize}[leftmargin=1.3em]
\item \texttt{proj\_rerand}: $\mathtt{proj}(c+\mathtt{enc0}(w))=\mathtt{proj}(c)+\ip{\rho'}{w}$. This is
  equation~(1) of the paper composed with the projection.
\item \texttt{cube\_ternary}: $x^3=x$ if and only if $x\in\{-1,0,1\}$, in every field, because
  $x^3-x=x(x-1)(x+1)$ and a field has no zero divisors.
\item $\mathtt{Rmix}$ is the shuffle relation of the paper: lengths $\tau$, $\sigma$ a permutation of
  $\{0,\dots,\tau-1\}$, every $w_i$ ternary of length $N$, and $\hat c_j=c_{\sigma(j)}+\mathtt{enc0}(w_{\sigma(j)})$.
\end{itemize}

\subsection{\texttt{SlotAttack.v}: the attack, as a theorem}

This file is not used by any other file; it records why the paper leaves the ring.
\texttt{slot\_swap\_products}: if $X^N+1$ has $N$ distinct roots and, at every root $\zeta$, the pair
$(f'(\zeta),g'(\zeta))$ is either $(f(\zeta),g(\zeta))$ or $(g(\zeta),f(\zeta))$, then
$(f-\alpha)(g-\alpha)=(f'-\alpha)(g'-\alpha)$ in the ring, for every $\alpha$. So a grand-product test
over the ring cannot see that slots were swapped. The file then checks a concrete instance over
$\mathbb{Z}_{17}$ with $N=4$ (roots $2,8,15,9$): the ring test passes for all $\alpha$
(\texttt{slot\_attack\_ring\_test}), the lists differ (\texttt{slot\_attack\_lists\_differ}), and the
projection with $\rho=(3,7,1,10)$ tells them apart (\texttt{slot\_attack\_projection\_sees}).

%% =====================================================================================
\section{Layer 3: the sumcheck protocol}
\label{sec:fm-sumcheck}
%% =====================================================================================

Section~\ref{sec:sumcheck} explained the sumcheck protocol. Here it is formalised in four files.

\subsection{\texttt{Sumcheck.v}: the protocol and one round}

Fix a summand $g:F^m\to F$ (a function on lists). Write
\[
  \mathtt{tsum}(rs,k)=\sum_{b\in\{0,1\}^k}g(rs,b),\qquad
  \mathtt{rnd}(rs,k,X)=\sum_{b\in\{0,1\}^k}g(rs,X,b),
\]
the true claim once a prefix $rs$ of challenges is fixed, and the true round function.
$\mathtt{sc\_verify}(\mathit{msgs},\mathit{chals},\mathit{claim})$ is the verifier: for each round
polynomial $p$ (a coefficient list) and challenge $r$ it checks $p(0)+p(1)=\mathit{claim}$ and
continues with $\mathit{claim}:=p(r)$; it returns the final claim.
\begin{theorem}[\texttt{round\_sound}]
Suppose the true round function has a representation $p^*$ with at most $d+1$ coefficients, the
prover's polynomial $p$ also has at most $d+1$ coefficients, $p(0)+p(1)=\mathit{claim}$, and the claim
is false: $\mathit{claim}\ne\mathtt{tsum}(rs,k+1)$. Then at most $d$ challenges $r$ make the next claim
true, that is $p(r)=\mathtt{tsum}(rs\|r,k)$.
\end{theorem}
\emph{Proof.} $p-p^*$ is not zero (otherwise the claim would be true), and every bad $r$ is a root
of it; apply \texttt{roots\_bound}. \texttt{sc\_verify\_honest} is completeness: if every message is
the true round function, the verifier ends with $g(rs,\mathit{chals})$.

\subsection{\texttt{SumcheckMask.v}: the zero-knowledge mask}

The mask of Section~\ref{sec:zksumcheck} is
$\mathtt{gmask}(a_0,cs)(x)=a_0+\sum_jx_j\,c_j(x_j)$, where $c_j$ is a polynomial (so the $j$-th
summand has no constant term). Define $W([])=0$ and $W(c::cs)=2^{|cs|}c(1)+2W(cs)$.
\begin{itemize}[leftmargin=1.3em]
\item \texttt{mask\_tsum}: $\sum_{b\in\{0,1\}^{|cs|}}\mathtt{gmask}(a_0,cs)(b)=2^{|cs|}a_0+W(cs)$. This
  is the value $G$ the prover announces.
\item \texttt{mask\_round\_poly}: in round $j$ the round polynomial of the mask has constant term
  $2^{k}(a_0+\text{earlier terms})+W(\text{later }c)$ and other coefficients $2^{k}c_j$, where $k$ is
  the number of later variables. So the round messages of the mask are an affine image of its
  coefficients, which is the heart of the zero-knowledge argument (Section~\ref{sec:fm-zk}).
\item \texttt{p2\_nonzero}: if $2\ne0$ in $F$ then $2^k\ne0$. The paper had not stated that the
  argument needs $2\ne0$; the formalisation found this.
\end{itemize}

\subsection{\texttt{SumcheckSound.v}: soundness against an adaptive prover}

The prover is a function \texttt{strat} from the challenges so far to the next message (it may
choose every message adaptively). \texttt{Hcert} says that every true round function is a
polynomial of degree at most $\deg(j)$ in round $j$; $\mathtt{dsum}(i,k)=\deg(i)+\dots+\deg(i+k-1)$.
$\mathtt{sc\_run}$ replays the verifier against \texttt{strat}.
\begin{theorem}[\texttt{sc\_sound}]
If the starting claim is false, then the number of challenge vectors $cs\in F^k$ for which the run
accepts and ends with the true value $g(pre,cs)$ is at most $\mathtt{dsum}(|pre|,k)\cdot|F|^{k-1}$.
That is, the probability is at most $\sum_j\deg(j)/|F|$.
\end{theorem}
\emph{Proof idea.} Induction on $k$. In the first round, at most $\deg$ challenges make the claim
true (\texttt{round\_sound}); they contribute $\deg\cdot|F|^{k-1}$ vectors (\texttt{card\_le\_head}).
For all other first challenges the next claim is still false, and the induction hypothesis bounds
each fibre (\texttt{card\_le\_cons}).

The file also states the per-round fact in the round-by-round form used in Section~\ref{sec:fs}:
\texttt{sc\_round\_rbr}: from a ``doomed'' state (false claim), whatever message the prover sends, at
most $\deg$ challenges lead to a state that is not doomed; and \texttt{sc\_doomed\_final}: a doomed
state with no variables left is rejected by the final check.

\subsection{\texttt{MaskedExt.v}: masked extensions and their degrees}

$Z(x)=\prod_ix_i(1-x_i)$ and $\msk T(x)=\mathtt{mstar}\ m\ T\ a\ x=\mle T(x)+aZ(x)$.
\texttt{mstar\_cube\_pos}: on the cube (with $m\ge1$) $\msk T=T$, because $Z$ vanishes there.
\texttt{mstar\_is\_poly}: as a function of one variable, with the later variables Boolean, $\msk T$
has degree $1$; in the last variable it has degree $2$. This is the ``Degrees'' remark of the
paper: the mask terms change only the last round's degree.

\subsection{\texttt{Interp.v}: from values at nodes to coefficients}

The implementation sends each round polynomial as its values at $0,2,3,\dots,d$; the value at $1$
is implied by the claim. The soundness theorems talk about coefficient lists. This file connects
the two by Lagrange interpolation, under the hypothesis that $1,\dots,D$ are non-zero in $F$ (so the
nodes are distinct).
\begin{itemize}[leftmargin=1.3em]
\item $\mathtt{interpF}(vals,x)=\sum_i vals_i\prod_{j\ne i}(x-j)/(i-j)$; $\mathtt{lagr}$ is the same
  as a coefficient list (\texttt{lagr\_eval}), with at most $|vals|$ coefficients, and it takes the
  value $vals_k$ at node $k$ (\texttt{lagr\_nodes}).
\item \texttt{esc\_verify} is the verifier on values: it rebuilds the value at $1$ as
  $\mathit{claim}-v_0$ and moves to $\mathtt{interpF}(\dots)(r)$.
\item \texttt{esc\_sound}: if \texttt{esc\_verify} accepts, then the ordinary \texttt{sc\_verify}
  accepts the interpolated coefficient lists with the same final claim. So every soundness theorem
  applies to the format the implementation uses.
\item \texttt{interp\_unique}: a polynomial with at most $L$ coefficients is determined by its values
  at $L$ distinct nodes; this lets the zero-knowledge files identify the model's interpolant with
  the true round function.
\end{itemize}

%% =====================================================================================
\section{Layer 4: Reed--Solomon codes and the opening}
\label{sec:fm-codes}
%% =====================================================================================

This layer formalises Sections~\ref{sec:commit}--\ref{sec:opening-lemmas} of this tutorial: the
row code, the proximity and consistency lemmas, the opening protocol and its soundness, and the
way the verifier compiles claims into slots.

\subsection{\texttt{Proximity.v}: the proximity lemma for any linear code}

A \emph{word} is a function from column indices to $F$; only columns $j<n$ matter. The file works
with an abstract code: a predicate \texttt{code} closed under addition and scaling, with
\emph{distance} $d$ in the sense that two codewords that agree outside fewer than $d$ columns are
equal. Let $2e<d$ and $e+1<|F|$. A word $x$ is \emph{close} if it agrees with some codeword outside
a set $E$ of at most $e$ columns (\texttt{agree\_out}, \texttt{close}).

The \emph{line lemma} (Section~\ref{sec:appendix}; Diamond--Posen Thm.~2.1) is
\textbf{a hypothesis} (\texttt{LineLemma}): if more than $e+1$ points $x+ty$ are close, then $x$ and
$y$ agree with codewords outside one common set of at most $e$ columns. From it the file proves the
extension to affine spaces of the appendix.
\begin{theorem}[\texttt{prox\_affine}]
Let $\pi(\gamma)=m+\sum_i\gamma_i\,\mathit{row}_i$. If more than an $(e+2)/|F|$ fraction of the
$\gamma\in F^R$ make $\pi(\gamma)$ close, then there is one set $E$ with $|E|\le e$ and codewords
$c_m,c_1,\dots,c_R$ such that $m$ agrees with $c_m$ and every $\mathit{row}_i$ with $c_i$ outside
$E$.
\end{theorem}
The proof follows the three steps of the appendix: every point is close (\texttt{all\_close},
averaging over lines with \texttt{count\_swap}), the decoding is affine (\texttt{affine\_decode}),
and the joint error support has at most $e$ columns (via \texttt{affine\_common\_nonzero}).

The same file proves the consistency lemma (Lemma~4.3 of the paper):
\texttt{consistency}: if $c_u\ne c_{dec}$ are codewords and $x$ agrees with $c_{dec}$ outside $E$,
then $\#\{j<n:c_u(j)=x(j)\}+d\le n+|E|$; and \texttt{far\_few\_pass}: a codeword agrees with a word
that is not close on fewer than $n-e$ columns.

\subsection{\texttt{ReedSolomon.v}: the row code}

$\mathtt{rs\_enc}(p)(j)=p(\omega^j)$, where the powers $\omega^0,\dots,\omega^{n-1}$ are distinct.
$\mathtt{rs\_code}(c)$: $c$ is the encoding of a polynomial with at most $k\le n$ coefficients.
\begin{itemize}[leftmargin=1.3em]
\item \texttt{rs\_dist}: the code has distance $n-k+1$ (a non-zero polynomial with $k$ coefficients has
  fewer than $k$ roots).
\item \texttt{rs\_prox\_affine}: \texttt{prox\_affine} for this code, still assuming the line lemma
  for it.
\item \texttt{rs\_hiding} (Lemma~4.4, hiding): for $|J|$ distinct opened columns and a row
  ``data followed by $|J|$ pads'', the map from pads to the opened entries is a bijection of
  $F^{|J|}$, whatever the data. So uniform pads make the opened entries uniform and independent of
  the data. \emph{Proof:} $(\mathit{data}\|\mathit{pads})(x)=\mathit{data}(x)+x^{|\mathit{data}|}\mathit{pads}(x)$, and a
  polynomial with $|J|$ coefficients is determined by its values at $|J|$ distinct points.
\end{itemize}

\subsection{\texttt{RsDecode.v}: the extractor's decoder}

The soundness theorem speaks of ``the decoded rows''. This file makes decoding a function:
$\mathtt{dec}(x)$ is the first message $p\in F^k$ (in the enumeration of all of them) whose encoding
is within distance $e$ of $x$, if any. This is not meant to run (it searches $|F|^k$ messages); it is
a mathematical definition.
\texttt{dec\_agree}: if $x$ agrees with a codeword outside $E$, $|E|\le e$, then $x$ agrees with the
encoding of $\mathtt{dec}(x)$ outside the same $E$ (unique decoding, from $2e<n-k+1$).
\texttt{dec\_many}: given a proximity statement for the code (a hypothesis \texttt{Prox}), if many
combinations are close, then every row decodes with one common error set. These are exactly the
facts the opening soundness needs about the extractor.

\subsection{\texttt{Opening.v}: the algebra of the opening}

\paragraph{Definitions.} For decoded rows $P_i$:
$\mathtt{msum}(P,cb)=\sum_{(i,w)\in cb}wP_i$ (a combination of rows),
$\mathtt{vsum}(P,cb,L)=\sum w\ip{P_i}{L}$, and $\mathtt{dm}=\mathtt{msum}(P,\text{mask rows})+c\,\mathtt{msum}(P,\text{claim rows})$,
the combination a slot vector should be. A \emph{term} is $(v,i,a)$: slot $v$, row $i$, coefficient
$a$; a claim is a list of terms, with value $\sum a\ip{P_i}{L_v}$.
\paragraph{Theorems.}
\begin{itemize}[leftmargin=1.3em]
\item \texttt{linear\_form\_zero\_count} and \texttt{batch\_count}: if the claimed values differ from
  the true values, at most $|F|^{K-1}$ of the $|F|^K$ batching vectors $\kappa$ give
  $\ip{\kappa}{\mathit{vals}}=\ip{\kappa}{\mathit{true}}$ (the line ``a false claim disappears in the
  batch'').
\item \texttt{c\_count}: if $V\ne V'$, at most one $c$ satisfies $\mu+cV=\mu'+cV'$ (the line ``opening
  equation holds for a wrong $V$'').
\item \texttt{opening\_values}: if every sent slot vector equals the decoded combination, the value
  check reads $\mu+cV=\sum\mathtt{vsum}(\text{masks})+c\cdot T$, where $T$ is the true batched value.
\item \texttt{batch\_sum}: summing the slot contributions over all slots gives
  $\ip{\kappa}{(\text{value of claim } i)_i}$.
\item \texttt{slot\_dichotomy} (Lemma~4.3 for one slot): either the sent vector equals the decoded
  combination, or its encoding agrees with the committed combination on few columns.
\end{itemize}

\subsection{\texttt{EqBatch.v}: batching with eq}

The sumchecks batch several statements ``$T_i=0$ on the cube'' and ``$e_j=0$'' into one claim
\[
  \sum_i\kappa_i\,\mle{T_i}(t)+\sum_jx_je_j=0
\]
with random $t\in F^m$ and random coefficients $(\kappa,x)$.
\begin{theorem}[\texttt{eq\_batch\_card}]
If some $T_i$ is non-zero somewhere on the cube, or some $e_j\ne0$, then at most
$(m+1)|F|^{m+n'-1}$ of the $|F|^{m+n'}$ challenges satisfy the batched claim ($n'$ is the number of
coefficients). The probability is at most $(m+1)/|F|$.
\end{theorem}
\emph{Proof idea.} Either $\mle{T_i}(t)=0$ (at most $m|F|^{m-1}$ values of $t$, by
\texttt{mle\_zero\_count}), or the coefficient vector is non-zero and a random linear combination
vanishes with probability $1/|F|$. \texttt{affine\_card}: a non-zero affine function $a+xb$ has at
most one root; this is the line ``$H$ wrong although the claim of B is true''.

\subsection{\texttt{MatQueries.v}: column queries on one matrix}

$\mathtt{macc}(J)$ says: the proximity vector and every slot vector have at most $k$ coefficients, and at
every queried column $j\in J$ their encodings equal the corresponding combinations of the committed
entries. \texttt{mat\_consistent}: if all slot vectors encode to the decoded combinations, they equal
them. \texttt{mat\_bad\_queries}: if the $\gamma$-combination is far from the code, or the matrix
decodes but some slot vector is inconsistent, then at most $\mathtt{fall}(n-e-1,\nq)$ of the
$\mathtt{fall}(n,\nq)$ query tuples pass. This is the last line of the error budget,
$(1-e/n)^{\nq}$.

\subsection{\texttt{OpenSound.v}: soundness of the opening}

This file puts the pieces together for two committed matrices (big B and small S). The challenges
are $\kappa\in F^K$, then $(c,\gamma_B,\gamma_S)$, then the two query tuples. The acceptance
predicate $\mathtt{AccO}$ is the value check plus $\mathtt{macc}$ for both matrices. The main hypothesis
\texttt{Hbatch} says that, for every $\kappa$, the slot contributions of the decoded rows sum to
$\ip{\kappa}{\mathit{true}}$; the decoding hypotheses say that a matrix with many close combinations
decodes.
\begin{theorem}[\texttt{open\_sound}]
If the claimed values differ from the true values, the number of accepting challenge tuples is at
most \texttt{open\_bound}, which as a fraction of all tuples is
\[
  \frac1{|F|}+\frac1{|F|}+\frac{p_B}{|F|}+\frac{p_S}{|F|}
  +\max_{M\in\{B,S\}}\frac{\mathtt{fall}(n_M-e_M-1,n_{q,M})}{\mathtt{fall}(n_M,n_{q,M})},
\]
where $p_M$ is the proximity constant of matrix $M$.
\end{theorem}
\emph{Proof idea.} A union bound over five events: $\kappa$ hides the error, $c$ hits the wrong
value, one of the two matrices is far but its combination is close, or the queries miss. For the
queries only one matrix needs to be caught, which gives the maximum instead of a sum (a finding of
the formalisation).

\subsection{\texttt{OpenSoundDec.v}, \texttt{SlotKinds.v}, \texttt{OpenClaims.v}}

\texttt{OpenSoundDec.v} instantiates the decoding hypotheses of \texttt{open\_sound} with the decoder
of \texttt{RsDecode.v}. What remains assumed is one \textbf{proximity statement per code}
(\texttt{ProxB}, \texttt{ProxS}); for $\Delta=3$ it is \texttt{rs\_prox\_affine} with $p=e+2$, which rests on
the line lemma, and for $\Delta=2$ it is the proximity gap of BCIKS20 with $p=n$.

\texttt{SlotKinds.v} formalises the two kinds of slot (Section~\ref{sec:opening}). A general claim is a
list of terms $(\mathit{slot},\mathit{row},\mathit{coef},\mathit{vector})$ and has value
$\mathtt{gtval}=\sum\mathit{coef}\cdot\ip{P_{\mathit{row}}}{L}$, where $L$ is the slot's fixed vector
(tensor slot) or the term's own vector (functional slot). The compilation \texttt{cbk}, \texttt{Lk}
builds, for each slot, the rows with their batched weights and the batched vector. The key identity
\begin{theorem}[\texttt{batch\_mixed}]
If every term of a functional slot uses that slot's single row, then
$\sum_v\mathtt{vsum}(P,\mathtt{cbk}(\kappa,v),\mathtt{Lk}(\kappa,v))=\ip{\kappa}{(\mathtt{gtval}(cl))_{cl}}$.
\end{theorem}
is exactly the hypothesis \texttt{Hbatch}. \texttt{OpenClaims.v} combines the three files:
\texttt{open\_sound\_claims} is opening soundness stated for a list of claims as the verifier compiles
them, with side conditions on the claims (which rows each slot may use). Both protocol theorems
apply this theorem.

\subsection{\texttt{LayoutLemmas.v} and \texttt{RowLayout.v}: evaluations are linear claims}

A table of $2^{lo+hi}$ entries is stored in $2^{hi}$ rows of $2^{lo}$ entries. \texttt{mle\_rows}:
\[
  \mle T(x_{lo},x_{hi})=\sum_{t<2^{hi}}\eq(t,x_{hi})\,\ip{\mathit{row}_t}{(\eq(y,x_{lo}))_{y}},
\]
so an evaluation of a multilinear extension is a linear claim whose vector is the $\eq$ vector of
the low variables and whose coefficients are the $\eq$ values of the high variables.
\texttt{gmask\_functional}: a sumcheck mask evaluated at a point is the inner product of its
coefficients with $(1,r_1,r_1^2,\dots)$, a functional slot.

\texttt{RowLayout.v} turns this into the implementation's layout. $\mathtt{rdtab}$ reads a table off
consecutive rows (\texttt{rdtab\_spec}); \texttt{tensor\_claim} states that the tensor terms
\texttt{tabt} have value $\mathit{coef}\cdot\mle T(x)$; a field element of $\mathrm{GF}(q^2)$ is stored as
two entries $(x_0,x_1)$ read as $x_0+tx_1$ (\texttt{ment}); and \texttt{mask\_row\_gmask} reads a
sumcheck mask off the mask-table row. These are used by the layouts of both proofs.

%% =====================================================================================
\section{Layer 5: soundness of the shuffle proof}
\label{sec:fm-mix}
%% =====================================================================================

These files prove Theorem~5.2 of the paper (Section~\ref{sec:sound} of this tutorial) for the protocol
as implemented. The order follows the proof: first the multiset step, then the three summands and
their sums, then each sumcheck, then the whole layer, and finally the layout and the opening.

\subsection{\texttt{MixMultiset.v}: from the logarithmic derivative to a shuffle}

A ciphertext $c=(u,v)$ with tag $\mathit{tag}$ becomes the form $(-1,u,v,\mathit{tag})$, and a
challenge $(\alpha,\rho_u,\rho_v,\beta)$ becomes a point; their inner product is
$\mathtt{proj}(c)+\beta\,\mathit{tag}-\alpha$, the quantity $h$ of the paper (\texttt{dot\_form}).
With the re-randomised inputs $rr_i$ tagged by $i$ and the outputs $co_j$ tagged by the committed
$\sigma_j$:
\begin{itemize}[leftmargin=1.3em]
\item \texttt{perm\_witness}: if the two lists of forms are permutations of each other, then there is
  a permutation $\pi$ with $co_j=rr_{\pi(j)}$ and $\sigma_j=\pi(j)$ (\texttt{is\_shuffle}). The tags
  $0,\dots,\tau-1$ are distinct in $F$ because $1,\dots,\tau\ne0$, so the matching is a permutation.
\item \texttt{multiset\_step}: if there is no such shuffle, the logarithmic-derivative identity holds
  for at most $(2\tau-1)|F|^{2N+1}$ of the $|F|^{2N+2}$ challenges. This is \texttt{logder\_count}
  applied to the contrapositive of \texttt{perm\_witness}, and it is the first line of the error
  budget.
\end{itemize}

\subsection{\texttt{MixSummands.v}: the three summands and their degrees}

With masked extensions $\msk u,\msk{\uh},\msk\sigma,\msk h$ and
$\msk{\hh}=\mle{\zh}+\beta\msk\sigma-\alpha$:
\begin{align*}
  g_A(x)&=\eq(x,t_c)\bigl[k_1\msk u\msk h+k_2\msk{\uh}\msk{\hh}\bigr]+k_3(\msk u-\msk{\uh})+\rho_A\,\mathtt{gmask}_A(x),\\
  g_B(y)&=\textstyle\sum_{k=1}^3\mle{\rho'_k}(y)\,\mle{W_k}(y,r)+c_B+\rho_B\,\mathtt{gmask}_B(y),\\
  g_C(x)&=\eq(x,t_z)\textstyle\sum_{k}k^C_k\bigl((\msk{w_k})^3-\msk{w_k}\bigr)+\rho_C\,\mathtt{gmask}_C(x).
\end{align*}
The theorems \texttt{certA}, \texttt{certB}, \texttt{certC} give the degree of every round function:
$3$ for $g_A$ ($5$ in the last round), $2$ for $g_B$, and $4$ for $g_C$ ($7$ in the last round). They
supply the hypothesis \texttt{Hcert} of \texttt{sc\_sound}, and the sums of these degrees are the
terms $3\ell+2$, $2\log N$ and $4m+3$ of the error budget.

\subsection{\texttt{MixIdentities.v}: true sum minus claimed sum}

Each sumcheck's error must be written in the form that \texttt{eq\_batch\_card} counts. For A:
\[
  \sum_x g_A(x)-(k_1+k_2+\rho_AG)=
  k_1\mle{(uh-1)}(t_c)+k_2\mle{(\uh\hh-1)}(t_c)+k_3\textstyle\sum(u-\uh)+\rho_A(G_{\mathrm{true}}-G)
\]
(\texttt{A\_claim\_diff}). For C the tables are $w_k^3-w_k$ (\texttt{C\_claim\_diff}). For B the
difference is the affine expression $(\msk h(r)-H)+\rho_B(G_{\mathrm{true}}-G)$ (\texttt{B\_claim\_diff}),
using $Nc_B=a_hZ(r)$ and the linearity of $\mathtt{mle}$ to swap it with the sum over coefficients.

\subsection{\texttt{MixSoundC.v}, \texttt{MixSoundA.v}, \texttt{MixSoundB.v}: the three sumchecks}

Each file fixes everything that happens before its challenge as an arbitrary finite set of
``earlier challenges'' and the prover's messages as functions of them.
\begin{itemize}[leftmargin=1.3em]
\item \texttt{soundC}: if some entry of some $W_k$ is not in $\{-1,0,1\}$, then
  $(t_z,k^C,\rho_C)$ and the rounds of C accept with probability at most $((m+1)+(4m+3))/|F|$.
  \emph{Proof:} $w^3-w\ne0$ for a non-ternary $w$ (\texttt{cube\_ternary}), so \texttt{eq\_batch\_card}
  with \texttt{C\_claim\_diff} bounds the challenges for which the claim is true; otherwise
  \texttt{sc\_sound} with \texttt{certC} bounds the rounds.
\item \texttt{soundA}: if there is no shuffle, the event ``A accepts and the claimed $H$ is correct''
  has probability at most $((2\tau-1)+(\ell+1)+(3\ell+2))/|F|$. \emph{Proof:} if $uh=1$, $\uh\hh=1$
  and $\sum u=\sum\uh$ on the cube, then the logarithmic-derivative identity holds, which
  \texttt{multiset\_step} bounds; otherwise the batched claim is false except with probability
  $(\ell+1)/|F|$, and \texttt{sc\_sound} handles the rounds.
\item \texttt{soundB}: the event ``$H\ne\msk h(r)$ and B accepts'' has probability at most
  $(1+2\log N)/|F|$, by \texttt{affine\_card} and \texttt{sc\_sound}.
\end{itemize}
Whether $H$ is correct is left to sumcheck~B: \texttt{soundA} covers the case $H=\msk h(r)$ and
\texttt{soundB} the case $H\ne\msk h(r)$.

\subsection{\texttt{MixBridge.v}: from tables to ciphertexts}

The sumchecks speak about tables; the relation speaks about ciphertexts. With the witness tables
$W_1,W_2,W_3$ read as the re-randomisers $(r_i,e_{1,i},e_{2,i})$ of ballot $i$, and the public tables
$\rho'_k$, $z^{\mathsf{pub}}$, $\zh$ equal to the projections, \texttt{bridge\_h} and
\texttt{bridge\_hh} show that the table $h$ at ballot $i$ is the inner product of the form of
$c_i+\EncZ(w_i)$ with the challenge point, and likewise for $\hh$. The key step is
\texttt{proj\_rerand} of \texttt{Bgv.v}.

\subsection{\texttt{MixSound.v}: the sumcheck layer as a whole}

\texttt{is\_witness} is ``$(\sigma,W)$ is a shuffle and every entry of $W$ is ternary''.
\begin{theorem}[\texttt{mix\_sound}]
If the decoded tables are not a witness, the three sumchecks with ideal checks of the claimed
values accept at most a $\mathtt{mix\_K}/|F|$ fraction of the challenge tuples, where
$\mathtt{mix\_K}=2\tau+4\ell+2\log N+5m+7$ (\texttt{mix\_bound\_fraction}, \texttt{mix\_K\_value}).
\end{theorem}
\emph{Proof.} If some entry is non-ternary, \texttt{soundC}. Otherwise there is no shuffle, and
\texttt{soundA} and \texttt{soundB} cover the two cases for $H$.

The file then adds an abstract opening. $\mathtt{vals9}$ are the nine values the verifier derives
from the transcript, and $\mathtt{tvI}$ their ideal values (the masked evaluations, the value of
$g_B$, and the two mask values). \texttt{link}: if the three runs succeed and $\mathtt{vals9}=\mathtt{tvI}$,
the ideal checks hold. With any opening that is sound for these nine values (hypothesis
\texttt{Hopen}), \texttt{card\_join} gives
\begin{theorem}[\texttt{mix\_full}]
If the decoded tables are not a witness, the number of accepting pairs (layer challenges, opening
challenges) is at most $\mathtt{mix\_bound}\cdot|O_o|+B_o\cdot|O_m|$.
\end{theorem}

\subsection{\texttt{MixLayout.v} and \texttt{MixFinal.v}: the real layout and the real opening}

\texttt{MixLayout.v} fixes the row layout of the implementation: $W_k$ in rows $k\cdot rw+t$ of the big
matrix; $\sigma$, the two halves of $u$ and $\uh$, and the mask table in the small matrix; the
offsets of the masks inside the mask table. It writes the nine claims $\mathtt{cl1},\dots,\mathtt{cl9}$
as lists of terms and proves
\begin{theorem}[\texttt{layout\_identity}]
On any rows $P$, the values of the nine claims are exactly the nine ideal values $\mathtt{tvI}$ computed
from the tables read off $P$.
\end{theorem}
\texttt{MixFinal.v} instantiates everything. The extractor decodes every row with \texttt{RsDecode}
and reads the tables with the layout; rows of the second batch may depend on the first challenge,
and \texttt{W\_fixed} shows the extracted witness does not. \texttt{Hopen\_F} discharges the abstract
opening hypothesis by \texttt{open\_sound\_claims} and \texttt{layout\_identity}.
\begin{theorem}[\texttt{mix\_theorem\_5\_2}]
Assume a proximity statement for each of the two codes, the code parameters ($k\le n$, distinct
evaluation points, $2e<n-k+1$), and $1,\dots,\tau$, $N$ non-zero in $F$. If the extracted
$(\sigma,W)$ is not a witness for the statement, then the number of accepting challenge tuples of
the protocol is at most $\mathtt{mix\_bound}\cdot|O_o|+B_o\cdot|O_m|$, that is, a fraction
\[
  \frac{2\tau+9\ell+7\log N+9}{|F|}+2\epsprox+\Bigl(1-\frac en\Bigr)^{\nq}
\]
where the query term is exactly $\max_M\mathtt{fall}(n_M-e_M-1,n_{q,M})/\mathtt{fall}(n_M,n_{q,M})$, which is
at most $(1-e/n)^{\nq}$.
\end{theorem}

%% =====================================================================================
\section{Layer 6: soundness of the decryption proof}
\label{sec:fm-dec}
%% =====================================================================================

The decryption proof (Section~\ref{sec:dec}) uses the same opening and two sumchecks of one common
shape, a \emph{zero test}.

\subsection{\texttt{ZeroTest.v}: a masked sumcheck that a polynomial vanishes on tables}

Given a polynomial $\varphi$ of degree at most $D$ and tables $T_1,\dots,T_n$, the statement is
$\varphi(T_i(b))=0$ for all $i$ and all cube points $b$. Three instances occur: $\varphi=w^3-w$ for
the shuffle's shortness, $\varphi=\prod_{v<\mathit{base}}(X-v)$ for the digits of the decryption
proof, and again $w^3-w$ for the key share $s,e$. The summand is
\[
  g_Z(x)=\eq(x,t)\sum_ik_i\,\varphi\bigl(\msk{T_i}(x)\bigr)+\rho\,\mathtt{gmask}(x),
\]
with round degrees $1+D$ and $1+2D$ in the last round (\texttt{degZ}, \texttt{certZ}).
\begin{theorem}[\texttt{soundZ}]
If $\varphi(T(b))\ne0$ for some table $T$ and some cube point $b$, the challenges $(t,k,\rho)$ and
the rounds accept with the correct claimed evaluations for at most a
$\bigl((m+1)+\sum_i\mathtt{degZ}(i)\bigr)/|F|$ fraction of all challenges.
\end{theorem}
The proof is the same as for \texttt{soundC}: \texttt{eq\_batch\_card} via \texttt{Z\_claim\_diff}, then
\texttt{sc\_sound}.

\subsection{\texttt{DecSound.v}: the sumcheck layer of the decryption proof}

\paragraph{Range and ternarity.} $\mathtt{rangep}(k)=\prod_{v<k}(X-v)$; \texttt{rangep\_root}: its roots are
exactly $0,1,\dots,k-1$ (as field elements). \texttt{phi3\_tern}: $w^3-w\ne0$ for non-ternary $w$.
\texttt{dval\_bound}: digits below $\mathit{base}$ give a number below $\mathit{base}^{k}$.

\paragraph{The relation.} With $N=2^{n_1}$, $\tau=2^{\ell}$, digit tables $\delta_d$ on $m=n_1+\ell$
variables and tables $s,e$ on $n_1$ variables, the noise coefficient of ballot $i$ at position $x$
is $E_{i,x}=\sum_d\mathit{base}^d\delta_d(x,i)-\mathit{off}$. The defects
\[
  D_{i,x}=t_i[x]-(u_i\cdot s)[x]-pE_{i,x},\qquad K_x=b[x]-(a\cdot s)[x]-pe[x]
\]
use the negacyclic product. \texttt{is\_dec\_witness}: all digits are roots of $\mathtt{rangep}$, $s,e$
are ternary, and all $D_{i,x}$ and $K_x$ vanish.
\begin{itemize}[leftmargin=1.3em]
\item \texttt{range\_exact}: if all digits are roots of $\mathtt{rangep}(\mathit{base})$, then every
  $E_{i,x}=v-\mathit{off}$ for a natural number $v<\mathit{base}^{k}$. This is the exact range proof.
\item \texttt{dec\_sound}: if the committed tables are not a witness, the compression challenges
  $\theta$ and the two zero tests accept (with ideal checks) at most a $\mathtt{dec\_K}/|F|$
  fraction, with $\mathtt{dec\_K}=(\mathit{base}+3)m+6n_1+\mathit{base}+5$ (\texttt{dec\_K\_value}).
  \emph{Proof:} a case split. An out-of-range digit or a non-ternary $s,e$ is caught by \texttt{soundZ};
  a non-zero defect is caught by the compression, because a table that is non-zero on the cube has
  a multilinear extension that vanishes at a random point with probability at most $m/|F|$
  (\texttt{mle\_zero\_card}).
\item \texttt{dec\_full}: as \texttt{mix\_full}, with an abstract opening for the $kd+6$ claimed values.
\end{itemize}

\subsection{\texttt{DecLayout.v}: the claims of the decryption proof}

The layout: digit table $d$ in rows $d\cdot rw+t$ of the big matrix; $s$, $e$ and the mask table in
rows $SO$, $SO+1$, $SO+2$ of the small matrix. Slots 1 and 2 are tensor slots, slots 3, 4, 5 are
functional slots on the rows of $s$, $e$ and the mask table. The $kd+6$ claims are: the compressed
relation (D), the compressed relation (K), one masked evaluation per digit table, the masked
evaluations of $s$ and $e$, and the two sumcheck masks.

The key step is how the verifier handles the secret product $u_i\cdot s$: the adjoint identity
\texttt{nmul\_adjoint} moves $u_i$ onto the public side, so
$\sum_i\eq(\theta,i)\,\ip{\eq(\theta_1)}{u_i\cdot s}=\ip{\sum_i\eq(\theta,i)\,\mathtt{nadj}(u_i,\eq(\theta_1))}{s}$
is a linear claim on the row of $s$ (\texttt{us\_val}).
\begin{itemize}[leftmargin=1.3em]
\item \texttt{claimD}, \texttt{claimK}: the public value minus the claim's value on the rows is the
  multilinear extension of the defect table of (D), respectively (K), at the challenge point.
\item \texttt{layout\_identityD}: the values of all $kd+6$ claims on any rows are the ideal values
  of \texttt{DecSound.v}.
\end{itemize}

\subsection{\texttt{DecFinal.v}: Theorem 7.1}

As \texttt{MixFinal.v}: the extractor decodes the rows, the opening hypothesis is discharged with
\texttt{open\_sound\_claims} and \texttt{layout\_identityD}, and
\begin{theorem}[\texttt{dec\_theorem\_7\_1}]
Assume a proximity statement for each code and the code parameters. If the extracted tables are not
a witness for $\Rdec$, the number of accepting challenge tuples is at most
$\mathtt{dec\_bound}\cdot|O_o|+B_o\cdot|O_m|$, a fraction
$\frac{(\mathit{base}+3)m+6\log N+\mathit{base}+7}{|F|}+2\epsprox+(1-e/n)^{\nq}$.
\end{theorem}

%% =====================================================================================
\section{Layer 7: zero knowledge over any field}
\label{sec:fm-zk}
%% =====================================================================================

\subsection{How zero knowledge is stated}

The prover's coins are uniform over a finite list $R$, and the simulator's coins are uniform over a
finite list $V$. ``The honest transcript and the simulated transcript have the same distribution''
is stated as
\[
  \mathtt{Permutation}\bigl(\mathtt{map}\ \mathit{hon}\ R,\ \mathtt{map}\ \mathit{sim}\ V\bigr):
\]
the list of honest transcripts over all prover coins is a reordering of the list of simulated
transcripts over all simulator coins. Two uniform distributions whose outcome lists are
permutations of each other are equal, so this is exact equality of distributions, with no
probabilities computed at all. The challenges are fixed (and ``good'': no division by zero, $Z\ne0$
at the evaluation points).

The proofs all have the same structure. One finds a map $f$ from prover coins to the \emph{free
values} of the transcript (the values the simulator draws) such that
(i) $f$ is a bijection from $R$ onto $V$, and (ii) the honest transcript is the simulator applied to
$f$ of the coins: $\mathit{hon}(x)=\mathit{sim}(f(x))$. This is \texttt{zk\_perm}
(\texttt{ZkTranscript.v}).

\subsection{\texttt{ZkBlocks.v}: elementary bijections}

\begin{itemize}[leftmargin=1.3em]
\item \texttt{mstar\_inj}, \texttt{mstar\_surj}: if $Z(r)\ne0$, then $a\mapsto\msk T(r)=\mle T(r)+aZ(r)$ is a
  bijection of $F$. A masked evaluation is uniform.
\item \texttt{sim\_msg}, \texttt{check\_head}: given the running claim and the non-constant coefficients of
  a round polynomial, the constant coefficient is determined (when $2\ne0$), because
  $p(0)+p(1)=2p_0+\sum_{i\ge1}p_i$. \texttt{sim\_run\_complete}: an accepted sequence of messages is
  determined by the claims, the challenges and the non-constant coefficients.
\item \texttt{affine\_list\_inj}/\texttt{surj}, \texttt{shift\_list\_inj}/\texttt{surj}: $c\mapsto t+sc$ ($s\ne0$) and
  $c\mapsto m+c$ are bijections of vectors.
\item \texttt{Zp\_zero\_card}: $Z(r)=0$ for at most $2m|F|^{m-1}$ points $r\in F^m$; this bounds the
  probability of bad challenges.
\end{itemize}

\subsection{\texttt{ZkSumcheck.v}: one masked sumcheck}

\texttt{hmsg} is the honest round-$j$ message of the masked sumcheck: the unmasked round polynomial
$pf(\text{prefix})$ plus $\rho$ times the mask's round polynomial (\texttt{hmsg\_eval}). Its
non-constant coefficients are $\mathrm{tl}(pf)+\rho2^{m-1-j}c_j$ (\texttt{hmsgs\_tails}), an affine image
of the mask coefficients. \texttt{mview} maps the mask $(a_0,cs)$ to
$(2^{|cs|}a_0+W(cs),\ \text{tails})$.
\texttt{mview\_inj}, \texttt{mview\_surj}: if $\rho\ne0$ and $2\ne0$, \texttt{mview} is a bijection: recover
the $c_j$ from the tails (dividing by $\rho2^{m-1-j}$), then $a_0=(G-W(cs))/2^{|cs|}$.

\subsection{\texttt{ZkMix.v} and \texttt{ZkOpen.v}: the free-value bijections}

\texttt{ZkMix.v}: for the whole sumcheck layer (masked evaluations and three masked sumchecks), the map
from the prover's layer coins (one mask coefficient per masked table, three sumcheck masks) to the
free values is a bijection (\texttt{mix\_view\_inj}, \texttt{mix\_view\_surj}). It is \emph{triangular}:
the mask coefficients are recovered from the evaluations first, which fixes the unmasked round
polynomials; then each sumcheck mask is recovered.

\texttt{ZkOpen.v}: for the opening, the prover's coins are the pads of the data rows and one uniform
mask per masked vector. The view is the data entries at the opened columns, and each mask plus a
shift that depends on the pads. \texttt{oview\_inj}, \texttt{oview\_surj}: this is a bijection. The pads
are recovered from the opened entries by \texttt{rs\_hiding}, then the masks by subtracting the
shifts.

\subsection{\texttt{ZkTranscript.v}: transcripts and simulators}

\begin{itemize}[leftmargin=1.3em]
\item \texttt{zk\_perm} (see above).
\item One sumcheck: $\mathtt{honT}(a_0,cs)=(G,\text{honest messages})$ and
  $\mathtt{simT}(cl)(G,\text{tails})=(G,\ \text{the messages rebuilt from the tails and the claim } cl+\rho G)$.
  \texttt{sc\_complete} (completeness of the masked sumcheck) and \texttt{sc\_honest\_sim}:
  $\mathtt{honT}(x)=\mathtt{simT}(\text{true sum})(\mathtt{mview}(x))$. \texttt{sc\_zk}: the permutation statement.
\item The layer with three sumchecks (\texttt{honL}, \texttt{simL}, \texttt{layer\_zk}) and with two
  (\texttt{honL2}, \texttt{simL2}, \texttt{layer2\_zk}, used for decryption). The simulator needs the true
  unmasked sum of each sumcheck; the hypothesis \texttt{HcX} says that the verifier can compute it
  from the masked evaluations. This is where the witness enters: for a valid witness the true sums
  are known values (often $0$).
\end{itemize}

\subsection{\texttt{ZkOpenTranscript.v}: the opening of one matrix}

The honest opening sends $\mu=\sum_v\ip{m_v}{L_v}$, the masked vectors $m_v+\mathit{shift}_v$, the
opened data entries, and the opened entries of the mask rows. The simulator draws the opened data
entries and the masked vectors uniformly, sets $\mu=\sum_v\ip{z_v}{L_v}-cv$ and solves the mask
entries from the check equations. Two hypotheses connect it to the claims: \texttt{Hmu} (the shifts
pair with the $L_v$ to the batched target value $cv$) and \texttt{Hsh} (the shifts at the opened
columns are computable from the opened data entries). \texttt{open\_honest\_sim} and
\texttt{open\_zk}.

\subsection{\texttt{ZkJoin.v}: joining the layer and the opening}

\begin{theorem}[\texttt{zk\_join}]
Let a first part have coins $R_1$, view map $f_1$ and honest/simulated transcripts $\mathit{hon}_1$,
$\mathit{sim}_1$, and a second part have coins $R_2$ and a view map $f_2(x,\cdot)$ that may depend
on the first coins $x$, with $\mathit{sim}_2$ allowed to use the first free values. If both parts
satisfy the hypotheses of \texttt{zk\_perm} (the second for every $x$), then the joint honest
transcripts over $R_1\times R_2$ are a permutation of the joint simulated transcripts over
$V_1\times V_2$.
\end{theorem}
This is needed because the mask table is randomness of the sumcheck layer \emph{and} data of a
committed row, so the opening's view depends on the layer's coins, and the opening's simulator
needs $cv$, which it computes from the layer's transcript.

The file then treats the opening over a subfield: pads over $K=\Zq$, masks over $F$ (sections
\texttt{Open2} and \texttt{OpenMix}), with one proximity vector per matrix whose shift depends on the
pads and on the slot masks. \texttt{open3\_zk} and \texttt{open3\_parts} are the version that is
instantiated later. In the simulator, the proximity mask row is solved last, from the slot mask
rows already simulated, an ordering detail that the formalisation made explicit.

\subsection{\texttt{ZkDecLayer.v}: the decryption layer}

The two zero-test sumchecks of the decryption proof, generic over $F$: the summands $f_G,f_T$ (zero
tests without masks), the Lagrange interpolants $pf_G,pf_T$ of their round functions
(\texttt{dec\_HpfG}, \texttt{dec\_HpfT}, using \texttt{certZ} for the degrees and the hypothesis that
$1,\dots,\max(2\mathit{base}+2,8)$ are non-zero), and the true sums $0$ for a witness
(\texttt{dec\_HcG\_true}, \texttt{dec\_HcT\_true}). \texttt{dec\_layer\_zk} is \texttt{layer2\_zk} for these.

%% =====================================================================================
\section{Layer 8: the concrete field $\mathrm{GF}(q^2)$}
\label{sec:fm-concrete}
%% =====================================================================================

\subsection{\texttt{ConcreteField.v}}

\begin{itemize}[leftmargin=1.3em]
\item \texttt{Q\_value}: $q=2^{78}+100663297=302231454903657394339841$.
\item \texttt{Q\_prime}: $q$ is prime, by a Pocklington certificate checked by computation
  (CoqPrime). The certificate uses the factorisation
  $q-1=2^{25}\cdot5\cdot7\cdot11\cdot1187\cdot19709623201$.
\item $\Zq$ is \texttt{Zpfield} of \texttt{Zpstar.v} at this $q$.
\item \texttt{three\_nonres}: no $x\in\Zq$ has $x^2=3$. \emph{Proof:} Euler's criterion. By Fermat,
  $x^{q-1}=1$; if $x^2=3$ then $3^{(q-1)/2}=1$, but computing $3^{(q-1)/2}\bmod q$ shows it is not.
\item For any field $K$ and non-square $\nu$, $K[t]/(t^2-\nu)$ (pairs $(a,b)$ with
  $(a,b)(c,d)=(ac+\nu bd,\ ad+bc)$ and inverse $(a,-b)/(a^2-\nu b^2)$) is a field
  (\texttt{ext\_field}). $\mathrm{GF}(q^2)$ is this field at $K=\Zq$, $\nu=3$ (\texttt{Fq2\_field}).
\item \texttt{Fq2\_char}: $k\cdot1\ne0$ for $0<k<q$. \texttt{Fq2\_two}: $2\ne0$.
\item \texttt{lf2}: the list of all $q^2$ elements, without duplicates (\texttt{lf2\_all},
  \texttt{lf2\_NoDup}). It is never computed: proofs only use these two facts about it.
\end{itemize}

\subsection{\texttt{ConcreteInstance.v}}

\texttt{mix\_theorem\_5\_2\_Fq2} and \texttt{dec\_theorem\_7\_1\_Fq2}: the two soundness theorems with the
field, its decidable equality, its axioms and its enumeration filled in. The characteristic
conditions $1,\dots,\tau\ne0$ and $N\ne0$ are discharged from $\tau<q$ and $N<q$
(\texttt{char\_tau}, \texttt{char\_N}). What remains a hypothesis: the instance and layout parameters,
the Reed--Solomon parameters, and one proximity statement per code. \texttt{lf2\_length}: $|F|=q^2$.

%% =====================================================================================
\section{Layer 9: the executable verifiers}
\label{sec:fm-exec}
%% =====================================================================================

Everything so far is mathematics over an abstract field. The files \texttt{src/Examples/Exec*.v}
contain programs that compute with integers, are extracted to OCaml, and run on real proofs. The
link between the two is one map and a family of \emph{refinement lemmas}.

\subsection{\texttt{ExecField.v}: field arithmetic on pairs of integers}

An element of $\mathrm{GF}(q^2)$ is represented by a pair of integers \texttt{xf}$\,=(a,b)$, meaning
$a+bt$. The operations \texttt{xadd}, \texttt{xsub}, \texttt{xneg} work componentwise modulo $q$;
$\mathtt{xmul}((a,b),(c,d))=(ac+3bd,\ ad+bc)\bmod q$; \texttt{xinv} uses $(a,b)^{-1}=(a,-b)/(a^2-3b^2)$
with the inverse of the norm computed as a power $(\cdot)^{q-2}$. The map
\[
  \mathtt{to\_F}(a,b)=(a\bmod q,\ b\bmod q)\in\mathrm{GF}(q^2)
\]
sends a representation to the field element.
\begin{itemize}[leftmargin=1.3em]
\item \texttt{to\_F\_add}, \texttt{to\_F\_sub}, \texttt{to\_F\_neg}, \texttt{to\_F\_mul}, \texttt{to\_F\_inv}: with no
  hypotheses, $\mathtt{to\_F}(\mathtt{xmul}\ x\ y)=\mathtt{to\_F}(x)\cdot\mathtt{to\_F}(y)$, and so on. These are
  the refinement lemmas: a computation on representations, read through \texttt{to\_F}, is the
  corresponding computation in the field.
\item \texttt{xeqb\_spec}: on \emph{reduced} pairs (both components in $[0,q)$), the boolean equality
  test agrees with equality in the field. Every operation returns reduced pairs.
\item \texttt{to\_F\_eq\_eval}, \texttt{to\_F\_zval}: the executable $\eq(r,s)$ and $Z(r)$ refine the
  abstract ones.
\end{itemize}

\subsection{\texttt{ExecSumcheck.v}: the sumcheck verifier}

\texttt{xeq\_table} computes the table $(\eq(\mathtt{pt}(i),r))_i$ (variable $0$ is the least significant
bit), \texttt{xinterp} is Lagrange interpolation, and \texttt{xsc\_verify} is the sumcheck verifier on
messages given as values at $0,2,\dots,d$.
\texttt{xsc\_sound}: if every message has at most $D<q$ entries and \texttt{xsc\_verify} accepts with final
claim $c$, then the abstract \texttt{sc\_verify} accepts the interpolated coefficient lists with final
claim $\mathtt{to\_F}(c)$. So the soundness theorem of \texttt{SumcheckSound.v} applies to what the code
checks.

\subsection{\texttt{ExecFFT.v} and \texttt{ExecNtt.v}: encoding rows}

\texttt{zenc\_horner} encodes a row by evaluating it at $\omega^j$ with Horner's rule (the
specification); \texttt{zfft} is a radix-2 number-theoretic transform; \texttt{zenc} uses the transform
if a run-time check succeeds ($n$ a power of two and $\omega^n=1$) and Horner otherwise.
\texttt{xfft\_zenc\_spec}: $\mathtt{zenc}=\mathtt{zenc\_horner}$, with no hypothesis. The proof splits a
polynomial into even and odd parts, $p(x)=p_e(x^2)+xp_o(x^2)$. \texttt{ExecNtt.v} restates the same
facts for the prover's side.

\subsection{\texttt{ExecHash.v} and \texttt{ExecMerkle.v}: hashing}

\texttt{ExecHash.v} defines SHA-256, the serialisation of field elements (10 bytes per $\Zq$ entry),
the Fiat--Shamir challenger (absorb a labelled message; draw field elements by rejection sampling;
draw distinct column indices) and Merkle path recomputation. \emph{Nothing is proven about these}
beyond four SHA-256 test vectors: they are outside the model of the soundness theorems. At
extraction, SHA-256 is replaced by HACL*'s verified implementation. \texttt{ExecMerkle.v} builds Merkle
trees and proves that the verifier's path check accepts every path the tree produces
(\texttt{root\_from\_path}), which is used for completeness.

\subsection{\texttt{ExecClaims.v}, \texttt{ExecOpen.v}: claims and opening checks}

\texttt{ExecClaims.v} builds the nine claims of the shuffle verifier as executable term lists and
proves \texttt{to\_F\_claims9}: mapped to the field they are exactly \texttt{MixLayout.claims9}.
\texttt{ExecOpen.v} implements the opening checks as a boolean function \texttt{xaccO} (the value
check and, for each matrix, the proximity and slot checks at the opened columns, with encodings
computed by the transform). \texttt{xaccO\_sound}: if \texttt{xaccO} returns \texttt{true}, then
\texttt{OpenSound.AccO} holds in the field, with the committed words replaced by the opened entries.

\subsection{\texttt{ExecNadj.v}: the fast adjoint}

The verifier computes $\rho'_1=\mathtt{nadj}(a,\rho_u)+\mathtt{nadj}(b,\rho_v)$, which takes $N^2$
products when done directly. \texttt{xnadj} computes it with three transforms and a pointwise
product, using $\mathtt{nadj}(a,\rho)(z)=a(z^{-1})\rho(z)$ at the roots $z=\psi\omega^j$ of $X^N+1$ and
the inverse transform (\texttt{NttInverse.v}). It checks the needed facts ($\psi^N=-1$, the order of
$\omega=\psi^2$, $N\cdot N^{-1}=1$) at run time and otherwise falls back to the direct computation.
\texttt{to\_F\_nadj}: with no hypotheses, the result is $\mathtt{nadj}$ in the field.

\subsection{\texttt{ExecPublic.v}, \texttt{ExecDerive.v}: the verifier's own values}

\texttt{ExecPublic.v} defines the public tables of the shuffle verifier (the blocks of $\rho'$ and the
projections of the input and output ciphertexts) and proves \texttt{Hver\_tables}: they satisfy the
hypothesis \texttt{Hver} of \texttt{mix\_full}. So in the final theorem these tables are no longer an
assumption. \texttt{to\_F\_mle}: the executable $\mathtt{xmle}(\mathit{tbl},r)=\ip{\mathit{tbl}}{\mathtt{xeq\_table}(r)}$ is
the multilinear extension.

\texttt{ExecDerive.v} implements the sumcheck layer of the verifier: shape checks, the three
sumchecks, the checks $\rho_A,\rho_C\ne0$, and the nine derived values.
\texttt{xderive\_sound}: if it returns values, they are $\mathtt{vals9}$ in the field, and the
precondition \texttt{PreAcc} of \texttt{mix\_full} holds.

\subsection{\texttt{ExecVerify.v}: the shuffle verifier}

The verifier is split in two.
\begin{itemize}[leftmargin=1.3em]
\item The \emph{front end} \texttt{xmix\_front} parses the bytes, replays the Fiat--Shamir transcript
  (drawing every challenge from the hash) and checks the Merkle paths; it produces a record with
  the challenges, the prover's messages and the opened entries.
\item The \emph{back end} \texttt{xmix\_back} does the algebra: length checks, \texttt{xderive}, the
  claims, and \texttt{xaccO}.
\end{itemize}
$\mathtt{xmix\_verify}(\mathit{proof})$ is ``front end, then back end''.
\begin{theorem}[\texttt{xmix\_back\_sound}]
If $\mathtt{xmix\_back}(f)=\mathtt{true}$, then there are values $\mathit{vals}$ with
$\mathtt{to\_F}(\mathit{vals})=\mathtt{vals9}$, the precondition \texttt{PreAcc} holds, the record's claims
are \texttt{claims9} built from the verifier's public tables, and \texttt{AccO} holds with the opened
entries as committed words.
\end{theorem}
These are exactly the events whose number \texttt{mix\_theorem\_5\_2} bounds. The front end is not
proven; as the paper says, transferring the theorems to the hashed proof is the subject of the
round-by-round analysis.

\subsection{\texttt{ExecDec.v}: the decryption verifier}

The same structure for the decryption proof: the $kd+6$ claims (\texttt{to\_F\_claimsD}), the sumcheck
layer (\texttt{xdec\_derive\_sound}), the back end (\texttt{xdec\_back\_sound}, with the events of
\texttt{dec\_theorem\_7\_1}), the front end, and the byte-level entry point
\texttt{xdec\_verify\_bytes}. \texttt{xrangep} is $\prod_{v<k}(z-v)$, and \texttt{to\_F\_rangep} shows it
refines $\mathtt{rangep}$.

\subsection{What is extracted, and what is trusted}

\texttt{src/Extraction/*/Extraction.v} extract \texttt{xmix\_verify}, \texttt{xdec\_verify\_bytes} and the
provers to OCaml. Integers become arbitrary-precision integers (Zarith); ten list functions
(\texttt{map}, \texttt{app}, \texttt{seq}, \dots) and four natural-number operations are replaced by their
OCaml counterparts for speed; SHA-256 is replaced by HACL*. The OCaml drivers in
\texttt{src/Executable} read files and compute the code parameters. Proven: the back ends (soundness
direction) and the provers' completeness (Section~\ref{sec:fm-prove}). Not proven and trusted: the
parsers, the transcript replay and the Merkle checks (tested only), the drivers, the extraction
mechanism, the OCaml compiler and Zarith, the replaced functions, and HACL*.

%% =====================================================================================
\section{Layer 10: the executable provers and their completeness}
\label{sec:fm-prove}
%% =====================================================================================

\emph{Completeness} means: for a valid witness and any random coins, if the prover returns a proof,
the verifier accepts it. The provers return no proof only when they would divide by zero (a
denominator $h_i$ or $\hh_j$ is zero, or $\rho_A$, $\rho_C$, $\rho_G$, $\rho_T$ is zero); how rarely
this happens is not part of the theorems.

\subsection{\texttt{ExecProveBase.v}: shared prover pieces}

\begin{itemize}[leftmargin=1.3em]
\item Coins come as two streams: integers modulo $q$ and 16-byte salts. $\mathtt{cut}(k,s)$ and
  $\mathtt{cuts}(ks,s)$ cut consecutive pieces of given lengths off a stream; every random row, pad,
  salt and mask is cut this way, in the order the prover draws them.
\item $\mathtt{padrow}(C,\mathit{row},\mathit{pads})$: a data row reduced modulo $q$, padded with zeros
  to $C$ entries, followed by its random pads. $\mathtt{split\_rows}(C,l)$ splits a table into rows of
  $C$ entries.
\item \texttt{commit}: encode every row with \texttt{zenc}, take the columns, hash each column with its
  salt, and build the Merkle tree.
\item The folding sumcheck prover \texttt{sc\_loop}. The tables are lists of length $2^k$; a term is
  a coefficient times a product of tables. In each round the prover computes the round polynomial's
  values at $0,2,\dots,d$ from pairs of adjacent table entries, adds $\rho$ times the mask's round
  value, absorbs the message, draws the challenge $r$, and \emph{folds} every table:
  $t'_i=t_{2i}+r(t_{2i+1}-t_{2i})$. In the last round it returns the final values, adding
  $a\,Z(\cdot)$ for masked tables.
\item \texttt{sc\_loop\_chals}: the verifier's transcript replay draws exactly the challenges the
  prover drew. Both make the same hash calls on the same data.
\end{itemize}

\subsection{\texttt{ExecFold.v}: completeness of the folding sumcheck}

\texttt{mle\_fold}: folding a table at $r_0$ and then evaluating its extension at $rs$ is evaluating
the original extension at $(r_0,rs)$. In words, folding fixes the first free variable.
\begin{theorem}[\texttt{sc\_loop\_complete}]
If the claim is the true sum of the terms plus $\rho$ times the mask's sum, the round degrees are
large enough for the terms (with masked tables counting degree $2$ in the last round), and
\texttt{sc\_loop} returns messages, point and final values, then \texttt{xsc\_verify} accepts them with
final claim equal to the terms at the final values plus $\rho$ times the mask at the point, and every
final value is the masked extension of its table at the point:
$\mathit{fin}_k=\mle{T_k}(pt)+a_kZ(pt)$.
\end{theorem}
\emph{Proof idea.} Each message is a polynomial of degree at most $d$ (by the closure rules of
\texttt{PolyFun}) whose values at $0$ and $1$ sum to the claim; interpolation at the nodes recovers it
(\texttt{interp\_step}), so the verifier's next claim is its value at $r$.

\texttt{ExecProveSC.v} is an earlier, stand-alone sumcheck prover that sums over the cube directly,
with its own completeness theorem \texttt{xsc\_complete}; no other file uses it.

\subsection{\texttt{ExecRows.v}: tables laid out in rows}

\texttt{rdtab\_list}: a table that the prover splits into rows and pads (with arbitrary pads) is read
back by the verifier's \texttt{rdtab} as itself. \texttt{mle\_list}: the multilinear extension of such
a table is \texttt{xmle} of the list. \texttt{ment\_row} and \texttt{maskCs\_row}: the mask-table row,
read as field elements, is the prover's mask table.

\subsection{\texttt{ExecProve.v}: the shuffle prover}

\texttt{coins\_of} cuts the coin streams into the record \texttt{coins} (random rows, pads, salts, and
the mask table). \texttt{xmix\_prove} is the prover of the paper's Figure~1 with Fiat--Shamir: commit
$W$ and $\sigma$; draw $\rho,\alpha,\beta$; compute $h$, $\hh$ and their inverses; commit $u$, $\uh$ and
the mask table; run the three folding sumchecks; run the opening (slot vectors
$m_v+c\cdot u_v$, proximity vectors, $\mu$); open the drawn columns; serialise.
\begin{itemize}[leftmargin=1.3em]
\item \texttt{prove\_transcript}: every successful run has a \emph{trace} $T$ (all its challenges and
  sumcheck outputs) with $\mathtt{ptrans}(P,T)$: the proof parts are built from $T$, every challenge was
  drawn from the challenger, the three sumcheck runs produced $T$'s data, and the verifier's
  transcript replay reproduces the same challenges.
\item \texttt{prove\_wf}: the output has the shapes the parser expects. \texttt{prove\_front}: the
  verifier's front end accepts and builds the record of $T$.
\end{itemize}

\subsection{\texttt{ExecSerial.v}, \texttt{ExecFront.v}, \texttt{ExecShape.v}}

\texttt{parse\_ser}: the parser reads back exactly what the serialiser writes. \texttt{merkle\_ok\_prover}:
the Merkle check accepts the prover's openings. \texttt{get\_F\_shape}, \texttt{get\_I\_shape},
\texttt{drawn\_nodup}: the challenger returns the right number of reduced field elements and of
distinct column indices in range. \texttt{ExecShape.v} collects shape facts (lengths, reducedness,
which rows each claim uses). These are completeness facts: they do not show that the parser rejects
malformed input.

\subsection{\texttt{ExecBack.v}, \texttt{ExecBackD.v}, \texttt{ExecBackO.v}: the back end accepts}

\texttt{valid} is the shuffle relation over the field for the given statement and witness.
\begin{itemize}[leftmargin=1.3em]
\item \texttt{ExecBack.v}: the true sums. \texttt{hh\_eq\_h}: $\hh_j=h_{\sigma(j)}$ for a valid witness
  (from \texttt{proj\_rerand}). \texttt{sumA}: the summand of A sums to $k_1+k_2$, since $uh=1$,
  $\uh\hh=1$ and $\sum_i1/h_i=\sum_j1/\hh_j$ (the second sum is a reordering of the first).
  \texttt{sumC}: the summand of C sums to $0$ for ternary entries. \texttt{sumB}: the sum of B.
\item \texttt{ExecBackD.v}: \texttt{derive\_complete}: the verifier's sumcheck layer accepts and
  returns the nine target values (by \texttt{sc\_loop\_complete} and the true sums).
\item \texttt{ExecBackO.v}: \texttt{layout\_targets} (the claims on the honest rows have the target
  values, by \texttt{layout\_identity}), \texttt{batch\_honest} (by \texttt{batch\_mixed}),
  \texttt{XeB}, \texttt{XeS} (the opened entries are the codewords of the honest rows at the opened
  columns), \texttt{enc\_comb} (encoding is linear), and \texttt{open\_complete}: the opening checks
  pass.
\end{itemize}

\subsection{\texttt{ExecComplete.v}: completeness of the shuffle proof}

\begin{theorem}[\texttt{xmix\_complete}]
Assume the code parameters, the coin and witness shapes, the statement shapes ($N+\tau<q$) and a
valid witness. If $\mathtt{xmix\_prove}$ returns a proof, then $\mathtt{xmix\_verify}$ accepts it.
\end{theorem}
\emph{Proof:} parse the bytes back, replay the transcript, the sumcheck layer accepts, the opening
accepts.

\subsection{\texttt{ExecDecProve.v}, \texttt{ExecDecComplete.v}: the decryption prover}

\texttt{ExecDecProve.v} defines the decryption prover \texttt{xdec\_prove} and its coin record
\texttt{dcoins}. The noise is written in digits,
$\mathtt{digit}(d,x)=\lfloor((x+\mathit{off})\bmod q)/\mathit{base}^d\rfloor\bmod\mathit{base}$. The range
check $\prod_{v<\mathit{base}}(D-v)$ is expanded into its coefficients (\texttt{xrcoefs}), since the
sumcheck prover works with sums of products of tables; the values are the same at every point.

\texttt{ExecDecComplete.v} proves, in the same steps as for the shuffle,
\begin{theorem}[\texttt{xdec\_complete}]
Assume the code parameters, the coin and witness shapes, $1\le n_1$, $2\mathit{base}+1\le\mathtt{DMAX}$,
and \texttt{dvalid}: $s,e$ ternary, $(E+\mathit{off})\bmod q<\mathit{base}^{kd}$, $t_i=su_i+pE_i$ and
$b=as+pe$. If $\mathtt{xdec\_prove}$ returns a proof, $\mathtt{xdec\_verify\_bytes}$ accepts it.
\end{theorem}
Steps: \texttt{dparse\_ser} (parsing), \texttt{dprove\_transcript} (the trace relation \texttt{dptrans}),
\texttt{sumG\_zero} (the range sum is $0$ for \emph{any} noise, because every digit is reduced modulo
$\mathit{base}$), \texttt{sumT\_zero}, \texttt{dscG}, \texttt{dscT}, \texttt{dderive}, \texttt{digits\_sum}
(the digits recombine to the number), \texttt{DtabT\_zero} and \texttt{KtabT\_zero} (the defect tables
vanish for a valid witness), \texttt{dlayout}, \texttt{dbatch}, \texttt{dopen\_complete} and \texttt{dwf}.

%% =====================================================================================
\section{Layer 11: zero knowledge at $\mathrm{GF}(q^2)$ and for the executable provers}
\label{sec:fm-zkexec}
%% =====================================================================================

This layer has two halves. The first instantiates the generic zero-knowledge theorems for the two
protocols as implemented, at $\mathrm{GF}(q^2)$. The second shows that the executable provers send
exactly the messages of those models.

\subsection{The models: \texttt{ZkLayer.v}, \texttt{ZkOpenMix.v}, \texttt{ZkMixJoin.v}}

\texttt{ZkLayer.v} instantiates \texttt{layer\_zk} with the shuffle prover's tables: the seven masked
evaluations $(u,\uh,\sigma,h)$ at $r_A$ and $W_1,W_2,W_3$ at $r_z$, the unmasked summands $f_A,f_B,f_C$,
their Lagrange interpolants $pf_A,pf_B,pf_C$, and the verifier's claims $c_A,c_B,c_C$. It proves that
for a valid witness the true sums are what the verifier computes from the masked evaluations
(\texttt{HcA\_true}, \texttt{HcB\_true}, \texttt{HcC\_true}). \texttt{mix\_layer\_zk} is the result.

\texttt{ZkOpenMix.v} instantiates \texttt{open3\_zk} with $K=\Zq$, $F=\mathrm{GF}(q^2)$ and the embedding
$x\mapsto(x,0)$. The data rows are the rows holding $W$, $\sigma$, the halves of $u$ and $\uh$, and the
mask table; the opening coins are the pads $p_B,p_S$, the slot masks $m_B,m_S$ and the proximity
masks $q_B,q_S$. \texttt{kz} rebuilds the prover's coin record from these coins (a mask over $F$
becomes two integer rows, its real and imaginary parts). The check equations of the generic
theorem are proven from the layout (\texttt{Hmu\_mix}, \texttt{HshB\_mix}, \dots). \texttt{mix\_open\_zk} is
the result.

\texttt{ZkMixJoin.v} joins them with \texttt{zk\_join}. \texttt{kOf} puts the layer coins into the mask
table. The opening's simulator needs the batched target value $cv$; \texttt{cvV} computes it from the
layer transcript (the masked evaluations and the final values of the three sumchecks), and
\texttt{cv\_honest} proves that this is the true value for a valid witness.
\begin{theorem}[\texttt{mix\_zk}]
For fixed good challenges, the list of honest transcripts (layer and opening) over all prover coins
is a permutation of the list of simulated transcripts over all simulator coins.
\end{theorem}

\subsection{The models for decryption: \texttt{ZkDecOpen.v}, \texttt{ZkDecJoin.v}}

\texttt{ZkDecOpen.v} instantiates \texttt{open3\_zk} for the decryption layout (two tensor slots on the
big matrix, three functional slots on $s$, $e$ and the mask table). \texttt{ZkDecJoin.v} joins it with
\texttt{dec\_layer\_zk}; the mask table, flattened to integers, is the third data row of the small
matrix (\texttt{dSx}), and \texttt{cv\_honest\_dec} computes $cv$ from the layer transcript.
\texttt{dec\_zk} is the permutation statement for the decryption proof.

\subsection{The executable prover is the model: \texttt{ExecZkSC.v}, \texttt{ExecZkTables.v}, \texttt{ExecZkMsgs.v}}

The model's honest prover computes messages from functions on the cube; the executable prover folds
tables of integers. \texttt{ExecZkSC.v} proves they agree for one sumcheck.
\begin{itemize}[leftmargin=1.3em]
\item \texttt{fold\_mle}, \texttt{foldx\_entry}: after folding at a prefix $rs$, entry $e$ of a table is
  $\mle T(rs,\mathtt{pt}(e))$.
\item $\mathtt{Sfull}$ is the masked summand of the \emph{original} tables:
  $\mathtt{Sfull}(x)=\text{terms}\bigl((\mle{T_i}(x)+a_iZ(x))_i\bigr)$. \texttt{core\_rnd},
  \texttt{last\_rnd}: the executable round values are the round function of $\mathtt{Sfull}$; before the
  last round the masks vanish because $Z=0$ on Boolean points, and in the last round the term
  $a\,z\,X(1-X)$ gives $aZ(pt,X)$.
\item \texttt{sc\_msgs\_hmsg}: for any model summand $f$ equal to $\mathtt{Sfull}$ and any interpolant
  $pf$ of its round functions, the messages of a run of \texttt{sc\_loop}, as field elements, are the
  values of the model's honest messages \texttt{hmsg} at the nodes.
\end{itemize}
\texttt{ExecZkTables.v} proves $\mathtt{Sfull}=f_A,f_B,f_C$ for the prover's actual tables
(\texttt{SfullA}, \texttt{SfullB}, \texttt{SfullC}), and \texttt{ExecZkMsgs.v} combines: the messages of A, B
and C (\texttt{msgsA}, \texttt{msgsB}, \texttt{msgsC}), the mask totals $G$ (\texttt{G\_run}), and the masked
evaluations (\texttt{evals\_run}) are those of the model.

\subsection{The opening and the whole run: \texttt{ExecZkOpen.v}, \texttt{ExecZkMain.v}}

With the coins $\mathtt{kz}(\dots)$, the rows of the committed matrices that hold the masks are the
real and imaginary parts of the model's masks (\texttt{PF\_mB}, \texttt{PF\_mS}). Since
$x=\mathrm{re}(x)+t\,\mathrm{im}(x)$, the mask part of each slot vector is the model's mask
(\texttt{msP\_mask}); hence the slot vectors (\texttt{vecs\_run}), $\mu$ (\texttt{mu\_run}), the proximity
vectors (\texttt{pvB\_run}, \texttt{pvS\_run}) and the opened entries (\texttt{colsB\_data},
\texttt{colsB\_mask}, \texttt{colsS\_data}, \texttt{colsS\_mask}) are the model's.
\begin{theorem}[\texttt{xmix\_prove\_refines}]
For prover coins built from model coins, every run of the executable shuffle prover returns a proof
whose field elements are the model's honest transcript at the run's challenges: the masked
evaluations, the $G$ values and the sumcheck messages (as values at the nodes) are those of the
layer model, and $\mu$, the slot and proximity vectors and the opened entries are those of the
opening model.
\end{theorem}

\subsection{Coins: \texttt{ExecZkCoins.v}}

\texttt{coins\_onto}: every coin record with the prover's shapes and entries in $[0,q)$ is
$\mathtt{kz}(\mathtt{kOf}'(\dots))$ of the decoded model coins, which are well shaped
(\texttt{decode\_shaped}); \texttt{coins\_back}: decoding the coins built from model coins gives the
model coins back. For fixed salts the correspondence is one to one, so uniform prover coins are
uniform model coins.

\subsection{The same for decryption: \texttt{ExecZkDecL.v}, \texttt{ExecZkDecO.v}, \texttt{ExecZkDecM.v}, \texttt{ExecZkDecC.v}}

\texttt{ExecZkDecL.v} does the layer: \texttt{peval\_rcoefs\_F} (the expanded coefficients evaluate to
$\mathtt{rangep}$), \texttt{SfullG}, \texttt{SfullT}, \texttt{msgsG\_run}, \texttt{msgsT\_run}, \texttt{G\_runD},
\texttt{evals\_runD}. \texttt{ExecZkDecO.v} does the opening: \texttt{dkz} builds the prover's coins from
model coins (the mask table is the layer coins; the random rows hold the real and imaginary parts
of two slot masks and a proximity mask on the big matrix, three slot masks and a proximity mask on
the small matrix), and \texttt{PFD\_mB}, \texttt{PFD\_mS}, \texttt{PFD\_dataB}, \texttt{PFD\_dataS},
\texttt{vec\_run}, \texttt{mu\_run}, \texttt{pvB\_run}, \texttt{pvS\_run} and the four column lemmas follow.
\texttt{ExecZkDecM.v} gives \texttt{xdec\_prove\_refines}, and \texttt{ExecZkDecC.v} the coin
correspondence (\texttt{dcoins\_onto}, \texttt{ddecode\_shaped}, \texttt{dcoins\_back}).

%% =====================================================================================
\section{How the files combine into the protocol}
\label{sec:fm-together}
%% =====================================================================================

We now follow the three main results from the top down and name, at each step, the file and the
theorem that supply it.

\subsection{Soundness of the shuffle proof}

\begin{enumerate}[leftmargin=1.6em,itemsep=2pt]
\item \textbf{The statement.} \texttt{mix\_theorem\_5\_2} (\texttt{MixFinal.v}) bounds the number of
  accepting challenge tuples when the extracted $(\sigma,W)$ is not a witness;
  \texttt{mix\_theorem\_5\_2\_Fq2} (\texttt{ConcreteInstance.v}) fixes the field to $\mathrm{GF}(q^2)$.
\item \textbf{Layer and opening.} It is \texttt{mix\_full} (\texttt{MixSound.v}) with the opening hypothesis
  discharged by \texttt{Hopen\_F}. \texttt{mix\_full} itself is \texttt{card\_join} (\texttt{Card.v}) applied to
  \texttt{mix\_sound} and to \texttt{link}.
\item \textbf{The sumcheck layer.} \texttt{mix\_sound} splits into \texttt{soundC} (non-ternary entry),
  \texttt{soundA} (no shuffle, $H$ correct) and \texttt{soundB} (no shuffle, $H$ wrong).
  \begin{itemize}[leftmargin=1.2em]
  \item \texttt{soundA} uses \texttt{multiset\_step} (\texttt{MixMultiset.v} $\leftarrow$
    \texttt{logder\_count} of \texttt{LogDer.v} $\leftarrow$ \texttt{roots\_bound} and \texttt{shift\_count}),
    \texttt{A\_claim\_diff} with \texttt{eq\_batch\_card} (\texttt{EqBatch.v} $\leftarrow$
    \texttt{mle\_zero\_count} of \texttt{Mle.v}), and \texttt{sc\_sound} (\texttt{SumcheckSound.v} $\leftarrow$
    \texttt{round\_sound} of \texttt{Sumcheck.v}) with the degree certificate \texttt{certA}
    (\texttt{MixSummands.v} $\leftarrow$ \texttt{mstar\_is\_poly} of \texttt{MaskedExt.v}).
  \item \texttt{soundB} uses \texttt{B\_claim\_diff}, \texttt{affine\_card} and \texttt{sc\_sound} with
    \texttt{certB}.
  \item \texttt{soundC} uses \texttt{cube\_ternary} (\texttt{Bgv.v}), \texttt{C\_claim\_diff}, \texttt{eq\_batch\_card}
    and \texttt{sc\_sound} with \texttt{certC}.
  \item The tables of the sumchecks are tied to the ciphertexts by \texttt{bridge\_h} and
    \texttt{bridge\_hh} (\texttt{MixBridge.v} $\leftarrow$ \texttt{proj\_rerand}, \texttt{nmul\_adjoint}).
  \end{itemize}
\item \textbf{The opening.} \texttt{Hopen\_F} is \texttt{open\_sound\_claims} (\texttt{OpenClaims.v}) after
  rewriting the claims' values with \texttt{layout\_identity} (\texttt{MixLayout.v} $\leftarrow$
  \texttt{RowLayout.v}, \texttt{LayoutLemmas.v}). \texttt{open\_sound\_claims} combines \texttt{batch\_mixed}
  (\texttt{SlotKinds.v}) with \texttt{open\_sound\_dec} (\texttt{OpenSoundDec.v}), which is \texttt{open\_sound}
  (\texttt{OpenSound.v} $\leftarrow$ \texttt{batch\_count}, \texttt{c\_count}, \texttt{opening\_values} of
  \texttt{Opening.v} and \texttt{mat\_bad\_queries} of \texttt{MatQueries.v}) with the decoder of
  \texttt{RsDecode.v} and the proximity hypotheses.
\item \textbf{The hypotheses that remain.} One proximity statement per code (for $\Delta=3$:
  \texttt{rs\_prox\_affine} of \texttt{ReedSolomon.v}, which rests on the line lemma), the code
  parameters, and $\tau,N<q$.
\end{enumerate}

\subsection{Soundness of the decryption proof}

\texttt{dec\_theorem\_7\_1} (\texttt{DecFinal.v}) is \texttt{dec\_full} (\texttt{DecSound.v}) with the opening
discharged by \texttt{open\_sound\_claims} and \texttt{layout\_identityD} (\texttt{DecLayout.v}).
\texttt{dec\_sound} splits into the range test and the ternary test (\texttt{soundZ} of \texttt{ZeroTest.v},
the same chain as \texttt{soundC}) and the two compressions (\texttt{mle\_zero\_count}). The exact range
statement is \texttt{range\_exact}. The secret product $u_i\cdot s$ becomes a linear claim through
\texttt{nmul\_adjoint} (\texttt{Negacyclic.v}).

\subsection{Zero knowledge, from the running code to the simulator}

For the shuffle proof three statements fit together.
\begin{enumerate}[leftmargin=1.6em,itemsep=2pt]
\item \texttt{xmix\_prove\_refines} (\texttt{ExecZkMain.v}): for prover coins built from model coins
  $(x,y)$, every run outputs the model's honest transcript $(\mathit{hon}_1(x),\mathit{hon}_2(x,y))$
  at the run's challenges.
\item \texttt{coins\_onto}, \texttt{decode\_shaped}, \texttt{coins\_back} (\texttt{ExecZkCoins.v}): for fixed
  salts, prover coins and model coins correspond one to one, so uniform prover coins give uniform
  model coins.
\item \texttt{mix\_zk} (\texttt{ZkMixJoin.v}): at fixed good challenges, the honest model transcripts
  over all model coins are a permutation of the simulated ones. It is \texttt{zk\_join}
  (\texttt{ZkJoin.v}) of \texttt{mix\_layer\_zk} (\texttt{ZkLayer.v} $\leftarrow$ \texttt{layer\_zk} of
  \texttt{ZkTranscript.v} $\leftarrow$ \texttt{mix\_view\_inj}/\texttt{surj} of \texttt{ZkMix.v} $\leftarrow$
  \texttt{mview\_inj}/\texttt{surj} of \texttt{ZkSumcheck.v} and \texttt{mstar\_inj}/\texttt{surj} of
  \texttt{ZkBlocks.v}) and \texttt{mix\_open\_zk} (\texttt{ZkOpenMix.v} $\leftarrow$ \texttt{open3\_zk} of
  \texttt{ZkJoin.v} $\leftarrow$ \texttt{rs\_hiding} of \texttt{ReedSolomon.v}), glued by \texttt{cv\_honest}.
\end{enumerate}
Together: for fixed challenges and salts, the field elements output by the extracted prover with
uniform coins have the distribution of the simulator's output. The development states the three
pieces separately; it does not contain one theorem that composes them. For the decryption proof
the three pieces are \texttt{xdec\_prove\_refines} (\texttt{ExecZkDecM.v}), \texttt{dcoins\_onto},
\texttt{ddecode\_shaped}, \texttt{dcoins\_back} (\texttt{ExecZkDecC.v}) and \texttt{dec\_zk}
(\texttt{ZkDecJoin.v}). Two limits remain, as in the paper: the challenges are fixed, whereas the
non-interactive prover derives them from the hash, and hashes, salts and Merkle paths are not
modelled.

\subsection{The executable code}

\texttt{xmix\_back\_sound} (\texttt{ExecVerify.v}) says that when the extracted shuffle verifier's back end
accepts, the events counted by \texttt{mix\_theorem\_5\_2} hold for the challenges of the transcript.
It rests on \texttt{xderive\_sound} (\texttt{ExecDerive.v} $\leftarrow$ \texttt{xsc\_sound} of
\texttt{ExecSumcheck.v} $\leftarrow$ \texttt{esc\_sound} of \texttt{Interp.v}), \texttt{to\_F\_claims9}
(\texttt{ExecClaims.v}), \texttt{Hver\_tables} and \texttt{to\_F\_mle} (\texttt{ExecPublic.v} $\leftarrow$
\texttt{to\_F\_nadj} of \texttt{ExecNadj.v}), \texttt{xaccO\_sound} (\texttt{ExecOpen.v} $\leftarrow$
\texttt{xfft\_zenc\_spec} of \texttt{ExecFFT.v}), and at the bottom the field refinement lemmas of
\texttt{ExecField.v}. \texttt{xmix\_complete} (\texttt{ExecComplete.v}) closes the loop in the other
direction. The decryption proof has the same pair: \texttt{xdec\_back\_sound} and \texttt{xdec\_complete}.

\subsection{Checking it yourself}

\begin{itemize}[leftmargin=1.3em]
\item \texttt{dune build} compiles all files; Rocq rejects any incorrect proof.
\item \texttt{./audit.sh} compiles \texttt{audit/Audit.v}, which runs \texttt{Print Assumptions} on $335$ main
  theorems, and fails unless every one is ``Closed under the global context''.
\item \texttt{./run\_all.sh} runs the build, the audit and the tests of the extracted code.
\end{itemize}

% =====================================================================================
\section{Index of the files}
\label{sec:fm-index}
% =====================================================================================

All $112$ files, in an order in which each file comes after every file it imports. ``Level'' is the
length of the longest chain of imports below the file. Line counts include comments.

{\small
\begin{longtable}{@{}rlp{0.12\textwidth}rp{0.42\textwidth}@{}}
\toprule
level & file & directory & lines & content\\
\midrule
\endhead
\bottomrule
\endfoot
0 & \texttt{Hierarchy.v} & Algebra & 302 & algebraic classes: monoid, group, ring, field, vector space\\
0 & \texttt{Functions.v} & Utility & 502 & fast modular exponentiation\\
0 & \texttt{Sha256.v} & Utility & 884 & SHA-256 definition and test vectors\\
0 & \texttt{Util.v} & Utility & 2427 & vectors, finite indices, modular facts\\
0 & \texttt{Prob.v} & Probability & 218 & rational weights (unused here)\\
0 & \texttt{Counting.v} & Crypto & 385 & probability as counting; averaging\\
1 & \texttt{Monoid.v} & Algebra & 93 & monoid identities\\
1 & \texttt{Product\_space.v} & Algebra & 83 & products of vector spaces (unused here)\\
1 & \texttt{Vector\_product\_space.v} & Algebra & 231 & powers of vector spaces (unused here)\\
1 & \texttt{Fermat.v} & Utility & 1780 & Fermat's little theorem\\
1 & \texttt{Insert.v} & Utility & 159 & vector insertion (unused here)\\
1 & \texttt{Distr.v} & Probability & 1723 & finite distributions (unused here)\\
1 & \texttt{Card.v} & Crypto & 304 & counts of non-decidable events; composing layers\\
2 & \texttt{Group.v} & Algebra & 173 & group identities\\
2 & \texttt{Zpstar.v} & Utility & 1404 & the field $\mathbb{Z}_p$ (and $\mathbb{Z}_p^*$)\\
3 & \texttt{Ring.v} & Algebra & 257 & ring identities; link to the \texttt{ring} tactic\\
3 & \texttt{SchnorrOrder.v} & Utility & 173 & Schnorr group orders (unused here)\\
4 & \texttt{Integral\_domain.v} & Algebra & 46 & integral domains\\
5 & \texttt{Field.v} & Algebra & 184 & field facts; link to the \texttt{field} tactic\\
6 & \texttt{Vector\_space.v} & Algebra & 127 & vector spaces (unused here)\\
6 & \texttt{PolyRoots.v} & Crypto & 191 & polynomials; root bound\\
7 & \texttt{LogDer.v} & Crypto & 621 & logarithmic-derivative lemma\\
8 & \texttt{Mle.v} & Crypto & 176 & cube, multilinear extension, zero count\\
8 & \texttt{Negacyclic.v} & Crypto & 248 & the ring $F[X]/(X^N+1)$, adjoint\\
8 & \texttt{PolyFun.v} & Crypto & 134 & degree of functions\\
8 & \texttt{ConcreteField.v} & Examples & 262 & $q$ prime; $\mathrm{GF}(q^2)$\\
9 & \texttt{Bgv.v} & Crypto & 121 & ciphertexts, projection, $\Rmix$\\
9 & \texttt{CubeIndex.v} & Crypto & 203 & cube points and indices; linearity of mle\\
9 & \texttt{MaskedExt.v} & Crypto & 131 & masked extensions and their degrees\\
9 & \texttt{NttInverse.v} & Crypto & 238 & adjoint at the roots; inverse NTT\\
9 & \texttt{Proximity.v} & Crypto & 453 & proximity and consistency lemmas\\
9 & \texttt{SlotAttack.v} & Crypto & 143 & the slot attack as a theorem\\
9 & \texttt{Sumcheck.v} & Crypto & 144 & sumcheck verifier; one round\\
10 & \texttt{Interp.v} & Crypto & 286 & values at nodes vs.\ coefficients\\
10 & \texttt{MixMultiset.v} & Crypto & 177 & multiset step of the shuffle\\
10 & \texttt{ReedSolomon.v} & Crypto & 195 & the row code; hiding\\
10 & \texttt{SumcheckMask.v} & Crypto & 122 & the sumcheck mask\\
10 & \texttt{SumcheckSound.v} & Crypto & 147 & sumcheck soundness; round-by-round step\\
10 & \texttt{ExecField.v} & Examples & 175 & executable field; \texttt{to\_F}\\
11 & \texttt{MixSummands.v} & Crypto & 235 & summands $g_A,g_B,g_C$ and degrees\\
11 & \texttt{Opening.v} & Crypto & 434 & algebra of the opening\\
11 & \texttt{RsDecode.v} & Crypto & 203 & the extractor's decoder\\
11 & \texttt{ZkBlocks.v} & Crypto & 246 & elementary bijections\\
11 & \texttt{ExecFFT.v} & Examples & 243 & NTT encoding\\
11 & \texttt{ExecHash.v} & Examples & 239 & SHA-256, challenger, Merkle paths\\
12 & \texttt{EqBatch.v} & Crypto & 158 & batching with eq\\
12 & \texttt{MatQueries.v} & Crypto & 141 & column queries on one matrix\\
12 & \texttt{SlotKinds.v} & Crypto & 209 & tensor and functional slots\\
12 & \texttt{ZkOpen.v} & Crypto & 159 & opening bijection\\
12 & \texttt{ZkSumcheck.v} & Crypto & 260 & one masked sumcheck\\
13 & \texttt{MixIdentities.v} & Crypto & 234 & true sum minus claim\\
13 & \texttt{OpenSound.v} & Crypto & 379 & soundness of the opening\\
13 & \texttt{ZkMix.v} & Crypto & 160 & layer free-value bijection\\
14 & \texttt{MixBridge.v} & Crypto & 153 & tables vs.\ ciphertexts\\
14 & \texttt{MixSoundB.v} & Crypto & 139 & sumcheck B\\
14 & \texttt{MixSoundC.v} & Crypto & 179 & sumcheck C\\
14 & \texttt{OpenSoundDec.v} & Crypto & 121 & opening soundness with the decoder\\
14 & \texttt{ZeroTest.v} & Crypto & 300 & masked zero-test sumcheck\\
14 & \texttt{ZkTranscript.v} & Crypto & 571 & permutation-form ZK; layers\\
15 & \texttt{DecSound.v} & Crypto & 565 & decryption sumcheck layer\\
15 & \texttt{LayoutLemmas.v} & Crypto & 165 & evaluations as linear claims\\
15 & \texttt{MixSoundA.v} & Crypto & 320 & sumcheck A\\
15 & \texttt{OpenClaims.v} & Crypto & 154 & opening soundness for compiled claims\\
15 & \texttt{ZkOpenTranscript.v} & Crypto & 242 & opening transcript, one matrix\\
16 & \texttt{MixSound.v} & Crypto & 483 & shuffle sumcheck layer; \texttt{mix\_full}\\
16 & \texttt{RowLayout.v} & Crypto & 323 & rows of the implementation\\
16 & \texttt{ZkDecLayer.v} & Crypto & 262 & decryption layer ZK\\
16 & \texttt{ZkJoin.v} & Crypto & 546 & joining layers; subfield opening\\
17 & \texttt{DecLayout.v} & Crypto & 475 & decryption claims and layout\\
17 & \texttt{MixLayout.v} & Crypto & 300 & shuffle claims and layout\\
17 & \texttt{ExecSumcheck.v} & Examples & 250 & executable sumcheck verifier\\
18 & \texttt{DecFinal.v} & Crypto & 371 & Theorem 7.1\\
18 & \texttt{MixFinal.v} & Crypto & 468 & Theorem 5.2\\
18 & \texttt{ExecClaims.v} & Examples & 175 & executable shuffle claims\\
19 & \texttt{ConcreteInstance.v} & Examples & 43 & soundness at $\mathrm{GF}(q^2)$\\
19 & \texttt{ExecOpen.v} & Examples & 412 & executable opening checks\\
20 & \texttt{ExecNadj.v} & Examples & 287 & fast adjoint\\
21 & \texttt{ExecPublic.v} & Examples & 140 & verifier's public tables\\
22 & \texttt{ExecDerive.v} & Examples & 297 & executable sumcheck layer\\
23 & \texttt{ExecProveSC.v} & Examples & 120 & stand-alone sumcheck prover (unused)\\
23 & \texttt{ExecVerify.v} & Examples & 401 & shuffle verifier\\
24 & \texttt{ExecDec.v} & Examples & 544 & decryption verifier\\
24 & \texttt{ExecMerkle.v} & Examples & 131 & Merkle trees\\
25 & \texttt{ExecProveBase.v} & Examples & 295 & prover pieces; folding sumcheck\\
26 & \texttt{ExecDecProve.v} & Examples & 218 & decryption prover\\
26 & \texttt{ExecRows.v} & Examples & 210 & tables in rows\\
26 & \texttt{ExecSerial.v} & Examples & 205 & parser reads serialiser\\
27 & \texttt{ExecFold.v} & Examples & 673 & completeness of folding sumcheck\\
27 & \texttt{ExecNtt.v} & Examples & 207 & NTT facts, prover side\\
28 & \texttt{ExecFront.v} & Examples & 231 & front end accepts; challenger shapes\\
28 & \texttt{ExecZkSC.v} & Examples & 375 & executable sumcheck = model messages\\
29 & \texttt{ExecShape.v} & Examples & 327 & shape facts\\
30 & \texttt{ExecProve.v} & Examples & 632 & shuffle prover\\
31 & \texttt{ExecBack.v} & Examples & 527 & true sums for a witness\\
32 & \texttt{ExecBackD.v} & Examples & 495 & verifier's sumcheck layer accepts\\
33 & \texttt{ExecBackO.v} & Examples & 748 & opening checks accept\\
34 & \texttt{ExecComplete.v} & Examples & 82 & completeness, shuffle\\
34 & \texttt{ExecDecComplete.v} & Examples & 1620 & completeness, decryption\\
34 & \texttt{ZkLayer.v} & Examples & 443 & shuffle layer ZK at $\mathrm{GF}(q^2)$\\
34 & \texttt{ZkOpenMix.v} & Examples & 813 & shuffle opening ZK\\
35 & \texttt{ExecZkOpen.v} & Examples & 513 & executable opening = model\\
35 & \texttt{ExecZkTables.v} & Examples & 211 & executable summands = model\\
35 & \texttt{ZkDecOpen.v} & Examples & 470 & decryption opening ZK\\
35 & \texttt{ZkMixJoin.v} & Examples & 604 & \texttt{mix\_zk}\\
36 & \texttt{ExecZkCoins.v} & Examples & 419 & shuffle coin correspondence\\
36 & \texttt{ExecZkMsgs.v} & Examples & 239 & executable messages = model\\
36 & \texttt{ZkDecJoin.v} & Examples & 724 & \texttt{dec\_zk}\\
37 & \texttt{ExecZkMain.v} & Examples & 318 & shuffle prover refines model\\
38 & \texttt{ExecZkDecL.v} & Examples & 418 & decryption layer refinement\\
39 & \texttt{ExecZkDecO.v} & Examples & 521 & decryption opening refinement\\
40 & \texttt{ExecZkDecM.v} & Examples & 332 & decryption prover refines model\\
41 & \texttt{ExecZkDecC.v} & Examples & 296 & decryption coin correspondence\\
\end{longtable}}

%% =====================================================================================
\part{Reference}
%% =====================================================================================

Setting: a linear code $\mathcal C$ over $\F$ of length $n$, distance $d$, and $e\le\lfloor(d-1)/3\rfloor$ with
$e+1<|\F|$. ``Close'' means within distance $e$ of $\mathcal C$.

\begin{lemma}[Line lemma; Diamond--Posen Thm.~2.1]
If more than $e+1$ points $\gamma\in\F$ make $x+\gamma y$ close, then $x$ and $y$ agree with codewords outside a
common set of at most $e$ coordinates.
\end{lemma}

Goal: the same for the affine space $U=m_{\mathsf{prox}}+\mathrm{span}(\mathit{row}_1,\dots,\mathit{row}_R)$, assuming
more than an $(e+2)/|\F|$ fraction of $U$ is close.

\begin{enumerate}[leftmargin=1.5em]
\item \textbf{Every point is close.} Fix $u\in U$. The lines through $u$ partition $U\setminus\{u\}$; each has
  $|\F|-1$ points besides $u$. If every such line had at most $e+1$ close points besides $u$, the close
  fraction of $U$ would be at most $(e+1)/(|\F|-1)\cdot(1-1/|U|)+1/|U|\le(e+2)/|\F|$, contradicting the
  assumption. So some line through $u$ has at least $e+2>e+1$ close points; by the line lemma all its points
  agree with codewords outside a common small set, in particular $u$ is close.
\item \textbf{A common error set.} Any two points of $U$ lie on a line all of whose points are close (step~1),
  so by the line lemma they share an error set of size $\le e$. Since $2e<d$, decoding is unique, and the
  decoding map is affine on $U$; the error vectors (point minus its decoding) form an affine space of vectors
  of weight $\le e$.
\item \textbf{Joint support $\le e$.} If the union of the supports of these error vectors had more than $e$
  coordinates, pick $e+1$ of them and a combination of at most $e+1$ error vectors with coefficients avoiding at
  most $e+1$ bad values (possible because $e+1<|\F|$) that is non-zero on all of them: weight $>e$, contradiction.
  So all rows agree with codewords outside one set of at most $e$ coordinates, contradicting the hypothesis
  of the proximity lemma.
\end{enumerate}

%% =====================================================================================
\section{Glossary of symbols}
\label{sec:glossary}
%% =====================================================================================

\begin{longtable}{@{}p{0.2\textwidth}p{0.75\textwidth}@{}}
\toprule
symbol & meaning\\\midrule\endhead
$q$ & prime modulus; paper $2^{78}+100663297$, toy $7681$, mini $17$\\
$\Zq$ & integers modulo $q$ (a field)\\
$N$ & ring degree (power of two); paper 4096, toy 4\\
$\Rq$ & $\Zq[X]/(X^N+1)$, BGV's ring; has zero divisors\\
$p$ & plaintext modulus, $2$\\
$\coeff(x)$ & coefficient vector of $x$ (concatenated for tuples)\\
$(a,b)$ & public key, $b=as+pe$\\
$c=(u,v)$ & ciphertext; $c_i$ inputs, $\ch_j$ outputs\\
$w=(r,e_1,e_2)$ & ternary re-randomiser; $\EncZ(w)=(ar+pe_1,br+pe_2)$\\
$\sigma$ & permutation: output $j$ is input $\sigma(j)$\\
$\tau$ & number of ballots (power of two); $\ell=\log\tau$, $m=\ell+\log N$\\
$\Mpk$ & $2N\times3N$ matrix with $\coeff(\EncZ(w))=\Mpk\coeff(w)$\\
$\F$ & $\mathrm{GF}(q^2)=\Zq[t]/(t^2-\nu)$, $|\F|\approx2^{156}$\\
$\rho,\rho'$ & projection vector in $\F^{2N}$; $\rho'=\rho\Mpk\in\F^{3N}$\\
$z_i,\zpub_i,\zh_j$ & projections of re-randomised inputs, of inputs, of outputs\\
$\alpha,\beta$ & log-derivative challenges; $h_i=z_i+\beta i-\alpha$, $\hh_j=\zh_j+\beta\sigma_j-\alpha$\\
$u,\uh$ & committed inverses $1/h$, $1/\hh$\\
$W_r,W_{e_1},W_{e_2}$ & witness tables, coefficient $y$ of ballot $i$ at index $y+Ni$\\
$\mle T$, $\msk T$ & multilinear extension; masked extension $\mle T+a_TZ$\\
$\eq(t,x)$, $Z(x)$ & equality polynomial; vanishing polynomial of the cube\\
$\id(r)$ & $\sum_i2^ir_i$, extension of the table $(0,\dots,\tau-1)$\\
$g_A,g_B,g_C$; $G_A,G_B,G_C$ & sumcheck masks and their sums\\
$t_c,k_i,\rho_A$; $\rho_B$; $t_z,k^C_i,\rho_C$ & challenges of sumchecks A, B, C\\
$r,r_y,r_z$ & end points of sumchecks A, B, C\\
$U,\hat U,s,H$ & claimed $\msk u(r),\msk\uh(r),\msk\sigma(r),\msk h(r)$\\
$\gamma_A,\gamma_C$; $\mathit{cl}_B$ & derived mask values; final claim of B\\
$C,k,n,d$ & data symbols per row, message length $C+\nq$, code length, distance $n-k+1$\\
$\varrho$ & nominal code rate, $1/4$ or $1/2$\\
$\Delta$, $e$ & proximity regime (3 or 2), $e=\lfloor(d-1)/\Delta\rfloor$\\
$\nq$ & number of column queries $=$ number of pads per row\\
$\kappa,\mu,c,\gamma$ & opening: batching coefficients, mask value, combination challenge, proximity coefficients\\
$m_v$, $m_{\mathsf{prox}}$, $\bar u_v$, $p$ & opening mask rows, sent slot vectors, proximity vector\\
$\epsprox$ & $(e+2)/|\F|$ ($\Delta=3$) or $n/|\F|$ ($\Delta=2$)\\
$\epsrbr$ & round-by-round soundness error\\
$B$, $\mathit{base}$, $\delta_d$ & noise range $\mathit{base}^k$, digit base, digit tables\\
$\theta,\theta_1,\theta_2$ & compression challenges of the decryption proof\\
\bottomrule
\end{longtable}

%% =====================================================================================
\section*{One-page summary}
%% =====================================================================================

\begin{recap}[The whole paper]
\begin{enumerate}[leftmargin=1.5em,itemsep=1pt]
\item BGV's ring $\Rq$ has zero divisors (CRT slots), so ring-level identity tests can be fooled by swapping a
  slot; this broke the CCS 2023 shuffle proof.
\item The shuffle relation is linear over $\Zq$ ($\coeff(\EncZ(w))=\Mpk\coeff(w)$) plus two side conditions:
  permutation and ternary randomness.
\item A random projection $\ip{\rho}{\coeff(\cdot)}$ into $\F=\mathrm{GF}(q^2)$ turns each ciphertext into one field
  element; distinct ciphertexts collide with probability $1/|\F|$.
\item A tagged logarithmic derivative ($\sum1/h=\sum1/\hh$) proves the permutation without revealing anything
  that depends on it.
\item Three masked sumchecks reduce ``inverses correct + sums equal'', ``$h$ is the linear image of $W$'' and
  ``$w^3=w$'' to nine evaluations of committed tables.
\item A Ligero-style commitment (Reed--Solomon rows with pads, Merkle tree over columns) with a zero-knowledge
  batched opening settles the nine claims; size $O(\sqrt{\nq\tau N})$.
\item Knowledge error $\approx2^{-128}$, statistical HVZK, Fiat--Shamir via round-by-round soundness (classical
  and quantum ROM), no lattice assumption in the proof.
\item Extracted, verified code at 1024 ballots: 0.29\,s per ballot to prove, 17.6\,s to verify, 3.78\,MB.
\item Machine-checked in Rocq (no axioms): soundness and zero knowledge of both proofs, completeness of both
  provers and their refinement of the zero-knowledge model (Part~\ref{part:formal}).
\item The decryption proof uses the same tools plus an exact digit-based range proof for the flooding noise.
\end{enumerate}
\end{recap}

\end{document}
